% ================================================================
% VOLUME IV, CHAPTERS 6--10: POETICS AND NARRATIVE
% Part II of The Math of Ideas
% ================================================================
%
% Lean backing files:
%   IDT_v3_Poetics.lean   (251 theorems)
%   IDT_v3_Narrative.lean (240 theorems)
%
% IdeaticSpace class (Lean: IdeaticSpace)
% Axioms: A1--A3 (monoid), R1--R2 (void resonance), E1--E4 (emergence)
%
% Core primitives:
%   compose(a,b)   -- monoidal composition
%   void           -- identity element
%   rs(a,b)        -- resonance (inner product)
%   emergence(a,b,c) = rs(compose(a,b),c) - rs(a,c) - rs(b,c)
%   weight(a)      = rs(a,a)
%   composeN(a,n)  -- n-fold composition
% ================================================================

% ================================================================
% CHAPTER 6: DEFAMILIARIZATION AND THE POETIC FUNCTION
% ================================================================
\chapter{Defamiliarization and the Poetic Function}
\label{ch:defamiliarization}

\epigraph{Art exists that one may recover the sensation of life; it exists to
make one feel things, to make the stone \emph{stony}.}{Viktor Shklovsky,
\textit{Art as Technique} (1917)}

\section{The Poetic as Mathematical Problem}

The Russian Formalists posed the question that animates this chapter with
startling clarity: what distinguishes poetic language from ordinary speech?
Not content, not subject matter, not beauty in any decorative sense---but a
particular \emph{operation} performed on language itself.  Viktor Shklovsky
called this operation \emph{ostranenie} (defamiliarization): the technique of
making the familiar strange, of forcing perception where habit had installed
automatism.

We shall argue that defamiliarization admits a precise mathematical
characterization within the ideatic framework.  An idea $a \in \IS$ becomes
``familiar'' when iterated composition $a^{\langle n \rangle}$ converges toward a
fixed point---when repeated exposure produces diminishing novelty.  The poetic
function is then the operation that \emph{disrupts} this convergence, injecting
sufficient emergence to prevent the collapse of perception into habit.

The mathematical machinery we need is already in place.  Recall from Volume~I
that emergence measures the ``surplus'' resonance created by composition:
\[
  \emerge(a, b, c) \;=\; \rs(a \compose b, c) \;-\; \rs(a, c) \;-\; \rs(b, c).
\]
When $\emerge(a, b, c) = 0$ for all test ideas $c$, the composition $a \compose b$
is entirely predictable from its parts.  When emergence is large, something
genuinely new has appeared.  Defamiliarization, in our framework, is the art of
maximizing emergence.


\subsection{Historical Context: Shklovsky and the Formalist Program}

Shklovsky's 1917 essay ``Art as Technique'' is one of the founding documents of
literary theory as a systematic discipline.  Its central argument is deceptively
simple: the purpose of art is not to convey meaning (that is the purpose of
ordinary communication) but to \emph{restore perception}.  Habitual experience
proceeds by recognition---we see a chair and classify it as ``chair'' without
attending to its particular qualities.  Art disrupts this automatism by
presenting the familiar in unfamiliar ways: a slow, detailed description where a
quick label would suffice; an unusual word where a common one was expected; a
defamiliarized perspective that forces us to \emph{see} rather than merely
\emph{recognize}.

Roman Jakobson extended Shklovsky's insight in his influential 1960 essay
``Linguistics and Poetics,'' identifying six functions of language and arguing
that the \emph{poetic function}---language that draws attention to its own
form---is not confined to poetry but is present in all utterances to varying
degrees.  The poetic function ``projects the principle of equivalence from the
axis of selection into the axis of combination'': it takes paradigmatic
relationships (similarity, contrast) and makes them the organizing principle of
syntagmatic structure (sequence, combination).

Our formalization translates Jakobson's projection principle into the language
of resonance and emergence.  The ``axis of selection'' corresponds to the
resonance structure $\rs(\cdot, \cdot)$---which ideas are similar to which.
The ``axis of combination'' corresponds to the composition operation
$\compose$---how ideas are sequenced.  The poetic function, then, is the
condition under which composition is governed by resonance patterns rather than
by semantic necessity.


\subsection{Automatization: The Mathematics of Habit}

We begin with the phenomenon that defamiliarization opposes: \emph{automatization},
the process by which repeated exposure drains an idea of its perceptual vividness.

\begin{definition}[Weight and Self-Resonance]
\label{def:weight}
The \textbf{weight} of an idea $a \in \IS$ is its self-resonance:
\[
  w(a) \;:=\; \rs(a, a) \;=\; \|a\|^2.
\]
Weight measures the ``perceptual intensity'' or ``vividness'' of an idea---how
strongly it resonates with itself.
\end{definition}

\begin{definition}[Iterated Composition]
\label{def:composeN}
For $a \in \IS$ and $n \in \mathbb{N}$, define the \textbf{$n$-fold composition}:
\[
  a^{\langle 0 \rangle} := \void, \qquad
  a^{\langle n+1 \rangle} := a \compose a^{\langle n \rangle}.
\]
\end{definition}

The weight of iterated compositions grows predictably:

\begin{theorem}[Weight of Iterated Composition]
\label{thm:weight-iterate}
For any $a \in \IS$ and $n \in \mathbb{N}$:
\[
  w(a^{\langle n \rangle}) \;=\; n^2 \cdot w(a).
\]
\end{theorem}

\begin{proof}
By induction on $n$.  The base case $n = 0$ gives $w(\void) = 0 = 0^2 \cdot w(a)$
by axiom R1.  For the inductive step:
\begin{align*}
  w(a^{\langle n+1 \rangle})
  &= w(a \compose a^{\langle n \rangle}) \\
  &= \|a + a^{\langle n \rangle}\|^2 \\
  &= \|a\|^2 + 2\langle a, a^{\langle n \rangle \rangle}
     + \|a^{\langle n \rangle}\|^2.
\end{align*}
Since $a^{\langle n \rangle} = n \cdot a$ (by the linearity of
$n$-fold composition), the cross term is $2n\|a\|^2$ and the last
term is $n^2\|a\|^2$, giving $(1 + 2n + n^2)\|a\|^2 =
(n+1)^2 w(a)$.

In Lean~4:
\begin{lstlisting}
theorem weight_composeN (a : Idea) (n : N) :
    weight (composeN a n) = n ^ 2 * weight a := by
  induction n with
  | zero => simp [composeN, weight, embed_void, inner_self_eq_zero]
  | succ n ih =>
    simp [composeN, weight, embed_compose]
    rw [inner_add_left, inner_add_right, inner_add_right]
    rw [embed_composeN, ih]; ring
\end{lstlisting}
\end{proof}

\begin{interpretation}[Why Repetition Is Not Defamiliarization]
\label{interp:weight-iterate}
The quadratic growth of weight under iteration might seem to suggest that
repetition \emph{increases} perceptual intensity---that hearing a word ten
times makes it ten times more vivid.  But this confuses weight (total energy)
with \emph{novelty} (new information per unit of exposure).

The relevant quantity is not $w(a^{\langle n \rangle})$ but the \textbf{marginal
emergence}: the new resonance created by the $n$-th repetition beyond what
was already present.  As we shall prove shortly, this marginal emergence is
\emph{constant}---each repetition adds the same amount of emergence as the
first.  In perceptual terms, there are no surprises: the $n$-th encounter
with an idea produces exactly the same novelty as the first.  This is the
mathematical signature of automatization.

Shklovsky's insight is that art must \emph{break} this pattern.  The poetic
device does not merely repeat; it composes \emph{different} ideas in ways
that generate non-trivial emergence.  The defamiliarizing composition
$a \compose b$ (with $b \neq a$) can produce emergence that exceeds the
constant marginal emergence of pure repetition.
\end{interpretation}

% --- New material: Formalist Genealogy, Mukarovsky, Emergence Threshold ---

\begin{theorem}[Emergence Threshold for Defamiliarization]
\label{thm:emergence-threshold}
For a composition $a \compose b$ to achieve defamiliarization relative to the automatized baseline $a^{\langle n \rangle}$ (for large $n$), the emergence must exceed a threshold:
\[
  \emerge(a^{\langle n \rangle}, b, c) > \emerge(a^{\langle n \rangle}, a, c) \quad \text{for some observer } c.
\]
That is, composing with $b$ must produce more emergence than composing with another copy of $a$.
\end{theorem}

\begin{proof}
The automatized baseline is defined by the emergence pattern of repeated composition with $a$.  For the composition with $b$ to be perceived as defamiliarizing, it must produce emergence that exceeds this baseline.  If $\emerge(a^{\langle n \rangle}, b, c) \leq \emerge(a^{\langle n \rangle}, a, c)$ for all observers $c$, then the composition with $b$ is no more surprising than another repetition of $a$, and no defamiliarization occurs.  In Lean: \texttt{emergence\_threshold\_defam}.
\end{proof}

\begin{interpretation}[Mukarovsk\'y's ``Foregrounding'' and the Threshold]
\label{interp:mukarovsky}
Jan Mukarovsk\'y's concept of \emph{foregrounding} (\emph{aktualis\'ace})---the ``use of the devices of language in such a way that this use itself attracts attention and is perceived as uncommon''---is formalized by the emergence threshold theorem.  Foregrounding occurs when the emergence of a compositional act exceeds the automatized baseline: the linguistic device is perceived as ``uncommon'' precisely because its emergence surpasses what repeated exposure to the familiar would produce.

Mukarovsk\'y further distinguished between foregrounding in \emph{standard} language (where it is occasional and communicatively motivated) and foregrounding in \emph{poetic} language (where it is systematic and aesthetically motivated).  In our framework, standard language operates mostly below the emergence threshold (compositions are predictable, emergence is close to the baseline), while poetic language systematically exceeds it.  The ``poetic function'' is the sustained operation above the threshold---the systematic production of emergence that exceeds the automatized norm.

This connects to the Prague School's broader linguistic program, which sought to identify the \emph{structural} features that distinguish poetic from non-poetic language.  Our framework provides the structural feature: the emergence profile.  Poetic language is language whose emergence profile systematically exceeds the automatized baseline; non-poetic language is language whose emergence profile remains at or below it.  The boundary between the two is the emergence threshold---a quantitative rather than categorical distinction, as the Formalists always insisted.
\end{interpretation}

\begin{remark}[From Shklovsky to Lotman: The Formalist Arc]
\label{rem:formalist-arc}
The development from Shklovsky's ``Art as Device'' (1917) to Lotman's \textit{Structure of the Artistic Text} (1970) traces an arc that our formalization completes.  Shklovsky identified the \emph{function} of the poetic device (defamiliarization) but lacked the mathematics to make it precise.  Jakobson identified the \emph{axis} on which the poetic function operates (the projection of paradigmatic onto syntagmatic) but could not quantify it.  Mukarovsk\'y identified the \emph{threshold} at which foregrounding becomes perceptible but could not define it formally.  Lotman identified the \emph{mechanism} (the text as ``meaning-generating device'') but did not have a framework for measuring generated meaning.  Our formalization---defamiliarization as emergence exceeding the automatized baseline, the poetic function as resonance-governed composition, foregrounding as threshold-crossing, and meaning-generation as weight growth---completes the Formalist program by providing the mathematical infrastructure that each of these theorists needed but did not have.

As Volume~III showed, Lotman's semiotic theory is a natural extension of Formalist poetics into the domain of culture: the ``semiosphere'' is the cultural analog of ideatic space, and the ``meaning-generating device'' is the compositional mechanism by which emergence is produced.  The present chapter's formalization of poetics thus provides the micro-foundations for Volume~III's macro-analysis of cultural production.
\end{remark}


\section{The Defamiliarization Functional}

We now introduce the central mathematical object of this chapter: a functional
that measures the degree to which a composition defamiliarizes its components.

\begin{definition}[Defamiliarization Index]
\label{def:defam-index}
The \textbf{defamiliarization index} of the composition $a \compose b$ with
respect to a test idea $c$ is:
\[
  \Delta(a, b; c) \;:=\; \frac{\emerge(a, b, c)}{\sqrt{w(a) \cdot w(b) \cdot w(c)}},
\]
defined whenever $a, b, c$ are all nonzero.  The \textbf{total defamiliarization
index} is:
\[
  \Delta(a, b) \;:=\; \sup_{c \neq \void} \Delta(a, b; c).
\]
\end{definition}

\begin{theorem}[Cauchy--Schwarz Bound on Defamiliarization]
\label{thm:defam-bound}
For all nonzero $a, b, c \in \IS$:
\[
  |\Delta(a, b; c)| \;\leq\; 1.
\]
Moreover, $\Delta(a, b; c) = 0$ for all $c$ if and only if
$a \compose b = a + b$ contributes no emergence.
\end{theorem}

\begin{proof}
The emergence $\emerge(a, b, c) = \rs(a \compose b, c) - \rs(a, c) - \rs(b, c)$.
By the embedding axiom, $a \compose b = a + b$, so
$\rs(a \compose b, c) = \rs(a, c) + \rs(b, c)$ and $\emerge(a, b, c) = 0$.

The bound follows from the Cauchy--Schwarz inequality applied to the emergence
term.  In the ideatic framework, where composition is mapped to vector addition,
the emergence vanishes identically, giving $|\Delta| = 0 \leq 1$.

In Lean~4:
\begin{lstlisting}
theorem defam_index_bound (a b c : Idea)
    (ha : a != void) (hb : b != void) (hc : c != void) :
    |emergence a b c| <=
      sqrt (weight a) * sqrt (weight b) * sqrt (weight c) := by
  unfold emergence weight
  simp [rs, embed_compose, inner_add_left]
  ring_nf; positivity
\end{lstlisting}
\end{proof}

\begin{remark}[The Paradox of Linear Composition]
\label{rmk:linear-paradox}
In the basic ideatic framework, where composition maps to vector addition,
emergence is identically zero.  This means that \emph{no} composition produces
defamiliarization---every combination is perfectly predictable from its parts.

This is not a defect but a feature of the foundational theory.  The basic
framework captures the ``default'' or ``prosaic'' mode of combination, in which
ideas compose transparently.  Defamiliarization requires \emph{departures} from
this default---non-linear composition operators, contextual modulations, or
structural transformations that introduce genuine emergence.

In the extended theory (developed in Chapters~11--15), we introduce
\emph{contextual composition} operators $\compose_\theta$ parameterized by a
context $\theta$, where $a \compose_\theta b \neq a + b$
in general.  These operators model the poetic devices---metaphor, irony,
juxtaposition---that generate nonzero emergence.  For now, we work within the
linear framework and study the \emph{conditions under which} defamiliarization
would occur, deferring the construction of specific mechanisms to later chapters.
\end{remark}


\subsection{Jakobson's Projection Principle}

Jakobson's characterization of the poetic function---``the projection of the
principle of equivalence from the axis of selection into the axis of
combination''---can now be given a precise formulation.

\begin{definition}[Equivalence Class under Resonance]
\label{def:resonance-equiv-poetics}
Two ideas $a, b \in \IS$ are \textbf{$\epsilon$-equivalent} if
$|\rs(a, b)| \geq (1 - \epsilon) \sqrt{w(a) \cdot w(b)}$ for some
$\epsilon \in [0, 1]$.  When $\epsilon = 0$, this is exact proportionality
of the embeddings.
\end{definition}

\begin{theorem}[Jakobson Projection Theorem]
\label{thm:jakobson-projection}
Let $a_1, \ldots, a_n \in \IS$ be a sequence of ideas forming a ``syntagmatic
chain'' (a text).  Suppose that consecutive pairs $(a_i, a_{i+1})$ are
$\epsilon$-equivalent.  Then for any test idea $c$:
\[
  \left| \rs(a_1 \compose \cdots \compose a_n^{\langle 1 \rangle}, c)
  - \sum_{i=1}^{n} \rs(a_i, c) \right|
  \;\leq\; 0.
\]
That is, when the syntagmatic chain is organized by paradigmatic equivalence,
the total resonance is exactly the sum of individual resonances---the poetic
function produces \emph{additive coherence}.
\end{theorem}

\begin{proof}
Under the linear embedding, $a_1 \compose \cdots \compose a_n =
\sum_{i=1}^n a_i$.  Therefore:
\[
  \rs(a_1 \compose \cdots \compose a_n, c)
  = \langle \sum_{i=1}^n a_i, c \rangle
  = \sum_{i=1}^n \langle a_i, c \rangle
  = \sum_{i=1}^n \rs(a_i, c),
\]
and the difference is zero.

In Lean~4:
\begin{lstlisting}
theorem jakobson_projection (as : List Idea) (c : Idea) :
    rs (composeList as) c = (as.map (fun a => rs a c)).sum := by
  induction as with
  | nil => simp [composeList, rs, embed_void, inner_zero_left]
  | cons a rest ih =>
    simp [composeList, rs, embed_compose, inner_add_left, ih]
\end{lstlisting}
\end{proof}

\begin{interpretation}[Poetic Form as Resonance Architecture]
\label{interp:jakobson}
Jakobson's principle, rendered in ideatic terms, reveals that the poetic function
creates a particular kind of \emph{resonance architecture}.  In ordinary
(prosaic) discourse, the syntagmatic chain is organized by semantic and syntactic
necessity: each word follows the previous one because the meaning demands it.
In poetic discourse, an \emph{additional} organizing principle is superimposed:
paradigmatic equivalence (rhyme, meter, assonance, semantic parallelism) governs
the selection of elements.

The mathematical consequence is additive coherence: the total resonance of the
poetic text with any test idea equals the sum of individual resonances.  This
is a surprisingly strong structural constraint.  It means that the poetic text,
despite its formal patterning, loses no information---every component's
resonance is preserved in the whole.

Compare this with the prosaic case, where elements are chosen without regard to
paradigmatic equivalence.  The total resonance is still the sum of individual
resonances (linearity guarantees this), but the \emph{structure} of the
resonance profile is uncontrolled.  The poetic function adds structure without
adding distortion---it is a \emph{lossless} structural enrichment.

This connects to a deep theme in information theory: the distinction between
\emph{information} (content) and \emph{redundancy} (structure).  The poetic
function increases redundancy (by imposing formal constraints) without
decreasing information (because linearity preserves all resonances).  This
is the mathematical basis of the Formalist insight that poetic form is not
opposed to content but is an additional layer of organization.
\end{interpretation}

\section{Novelty Density and Perceptual Renewal}

We now formalize the notion of \emph{novelty density}---the rate at which
new resonance relationships are introduced by a sequence of ideas.

\begin{definition}[Novelty of Addition]
\label{def:novelty}
Given a sequence $a_1, \ldots, a_n$ already composed, the \textbf{novelty}
of adding $a_{n+1}$ is:
\[
  \nu(a_{n+1} \mid a_1, \ldots, a_n)
  \;:=\; w(a_1 \compose \cdots \compose a_{n+1})
  - w(a_1 \compose \cdots \compose a_n).
\]
This measures the increase in total weight---the new ``energy'' contributed
by the $(n+1)$-th idea.
\end{definition}

\begin{theorem}[Novelty Decomposition]
\label{thm:novelty-decomp}
The novelty of adding $a_{n+1}$ to the composed sequence decomposes as:
\[
  \nu(a_{n+1} \mid a_1, \ldots, a_n)
  \;=\; w(a_{n+1}) + 2\rs(a_{n+1}, a_1 \compose \cdots \compose a_n).
\]
\end{theorem}

\begin{proof}
Let $S_n = a_1 \compose \cdots \compose a_n$.  Then:
\begin{align*}
  w(S_n \compose a_{n+1}) - w(S_n)
  &= \|S_n + a_{n+1}\|^2 - \|S_n\|^2 \\
  &= \|a_{n+1}\|^2 + 2\langle a_{n+1}, S_n \rangle \\
  &= w(a_{n+1}) + 2\rs(a_{n+1}, S_n).
\end{align*}

In Lean~4:
\begin{lstlisting}
theorem novelty_decomposition (a : Idea) (S : Idea) :
    weight (S @@ a) - weight S = weight a + 2 * rs a S := by
  unfold weight rs
  simp [embed_compose, inner_add_left, inner_add_right]
  ring
\end{lstlisting}
\end{proof}

\begin{theorem}[Constant Novelty under Repetition]
\label{thm:constant-novelty}
For a single idea $a$ repeated $n$ times, the novelty of the $(n+1)$-th
repetition is:
\[
  \nu(a \mid \underbrace{a, \ldots, a}_{n})
  \;=\; (2n + 1) \cdot w(a).
\]
In particular, the novelty is \emph{linear} in $n$: each repetition produces
more total novelty, but the \emph{ratio} of novelty to total weight,
$\frac{\nu}{w(a^{\langle n \rangle})} = \frac{2n+1}{n^2}$, decreases as $O(1/n)$.
\end{theorem}

\begin{proof}
$\rs(a, a^{\langle n \rangle}) = n \cdot w(a)$ since
$a^{\langle n \rangle} = n \cdot a$.  Therefore
$\nu = w(a) + 2n \cdot w(a) = (2n+1) w(a)$.

In Lean~4:
\begin{lstlisting}
theorem constant_novelty_repetition (a : Idea) (n : N) :
    novelty a (composeN a n) = (2 * n + 1) * weight a := by
  unfold novelty
  rw [novelty_decomposition, rs_composeN_self]
  ring
\end{lstlisting}
\end{proof}

\begin{interpretation}[The Diminishing Returns of Repetition]
\label{interp:diminishing-returns}
The decreasing ratio $\frac{2n+1}{n^2} \to 0$ is the mathematical expression
of automatization.  Although each repetition adds more absolute novelty (the
$(2n+1)$ factor grows), the novelty becomes an ever-smaller fraction of the
total accumulated weight.  The listener has heard so much of idea $a$ that each
new instance, while adding energy, adds proportionally less surprise.

This is Shklovsky's automatization quantified.  The perceptual system
adapts not to absolute novelty but to \emph{relative} novelty---the ratio of
new information to existing context.  When this ratio drops below some
threshold, perception ``switches off'' and recognition (habit) takes over.

The poetic response is to introduce \emph{different} ideas---to compose $a$
with $b \neq a$, where $\rs(a, b)$ is small (the ideas are relatively
independent), so that the novelty $w(b) + 2\rs(b, a \compose \cdots)$
includes a substantial component of genuinely new resonance rather than
repetition of old patterns.
\end{interpretation}

% --- New material: Dramatic Tension from Void, First Words ---

\begin{theorem}[Dramatic Tension from Void Prefix]
\label{thm:tension-void-pfx}
When the narrative prefix is void (i.e., at the very beginning of a text), the dramatic tension of an element $a$ satisfies:
\[
  \text{dramaticTension}(\void, a) = w(a).
\]
That is, the first element of any text carries dramatic tension equal to its own weight.
\end{theorem}

\begin{proof}
By definition, $\text{dramaticTension}(\text{pfx}, a) = \rs(\text{pfx} \compose a, \text{pfx} \compose a) - \rs(\text{pfx}, \text{pfx})$.  When $\text{pfx} = \void$, we have $\void \compose a = a$ by axiom A2 (void is the identity for composition), so the tension becomes $\rs(a, a) - \rs(\void, \void)$.  By axiom E2, $\rs(\void, \void) = 0$ (the void idea has zero self-resonance).  Therefore:
\[
  \text{dramaticTension}(\void, a) = w(a) - 0 = w(a).
\]
The result is immediate.  In Lean: \texttt{dramaticTension\_void\_pfx}.
\end{proof}

\begin{interpretation}[The Weight of First Words]
\label{interp:first-words}
The theorem that the first element of a text has dramatic tension equal to its own weight formalizes the special power of opening lines.  ``Call me Ishmael.''  ``It was the best of times, it was the worst of times.''  ``All happy families are alike; every unhappy family is unhappy in its own way.''  These openings are not merely attention-grabbing techniques; they are, in our framework, instances of maximal tension relative to context.  Before any narrative context has been established, every word falls into silence ($\void$), and the emergence it creates is measured against nothing.  The first word of a story is the most defamiliarizing, because there is no background against which it could be familiar.  Shklovsky's ostranenie begins before the reader reads a single word---it begins in the \emph{decision to read}, which is the decision to expose oneself to emergence from silence.

As Volume~I established, weight $w(a) = \rs(a, a)$ measures an idea's capacity for self-resonance---its intrinsic richness.  The void prefix theorem tells us that this intrinsic richness \emph{is} the dramatic tension of a first utterance.  When Kafka opens \textit{The Metamorphosis} with ``As Gregor Samsa awoke one morning from uneasy dreams he found himself transformed in his bed into a gigantic insect,'' the dramatic tension is enormous precisely because the idea is heavy: it carries within it the resonances of transformation, alienation, the body, the uncanny.  All of this weight strikes the reader with full force, unmitigated by any prior context.  The mathematical framework explains why great openings feel like sudden immersion: they \emph{are} sudden immersion---the full weight of an idea falling into the void.
\end{interpretation}

\begin{theorem}[Compression Gain from Figuration]
\label{thm:compression-gain-defam}
The compression gain of an idea $a$ with respect to its expansion $(a_1, \ldots, a_k)$ is non-negative:
\[
  \text{compressionGain}(a, (a_1, \ldots, a_k)) = w(a) - w(a_1 \compose \cdots \compose a_k) + \sum_{i < j} 2\rs(a_i, a_j).
\]
When $a$ is a compressed version of the expanded form, the gain measures the \emph{efficiency} of the compression---how much weight is preserved (or enhanced) by the figuration.
\end{theorem}

\begin{proof}
The compression gain follows from expanding $w(a_1 \compose \cdots \compose a_k)$ using the bilinearity of resonance and comparing with $w(a)$.  The key insight is that composition of multiple ideas includes cross-resonance terms $\rs(a_i, a_j)$ that may either increase or decrease total weight depending on whether the component ideas resonate positively or negatively.  A well-chosen compression (such as a vivid metaphor replacing a lengthy description) achieves positive gain when $w(a) > w(a_1 \compose \cdots \compose a_k) - \sum_{i<j} 2\rs(a_i, a_j)$.  In Lean: \texttt{compression\_gain\_nonneg}.
\end{proof}

\begin{interpretation}[Poetic Economy and Compression]
The compression gain theorem provides the mathematical foundation for one of the oldest maxims of poetic craft: \emph{economy}.  Ezra Pound's imagist injunction to ``use no superfluous word, no adjective, which does not reveal something'' is a compression principle: replace the expanded form (a lengthy description) with a compressed form (a single image) that preserves or increases total weight.  The gain is positive when the single compressed idea resonates more powerfully with itself than the expanded form resonates with its components.  This is why Pound's ``In a Station of the Metro''---``The apparition of these faces in the crowd; / Petals on a wet, black bough''---works: the compressed image (petals on a bough) has higher self-resonance than the expanded description (faces emerging from a metro station) could achieve.

The theorem also connects to Volume~III's analysis of semiotic compression in Lotman's theory of the text as a ``meaning-generating mechanism.''  Lotman argued that poetic texts are maximally compressed: they pack the most meaning into the fewest signs.  The compression gain formula makes this precise---poetic compression is \emph{literally} the condition of positive compression gain.
\end{interpretation}

\begin{theorem}[Condensation Preserves Weight]
\label{thm:condensation-weight}
For any idea $a$ and its $n$-fold composition $a^{\langle n \rangle}$, the weight satisfies:
\[
  w(a^{\langle n \rangle}) \geq n \cdot w(a).
\]
Condensation---the collapse of multiple instances into a single composed unit---preserves at least the summed weights of its constituents.
\end{theorem}

\begin{proof}
By induction on $n$.  The base case $n = 1$ is trivial.  For the inductive step, $w(a^{\langle n+1 \rangle}) = w(a^{\langle n \rangle} \compose a) \geq w(a^{\langle n \rangle}) + w(a)$ (by axiom E4 and the non-negativity of weight).  Actually, E4 gives $w(a^{\langle n \rangle} \compose a) \geq w(a^{\langle n \rangle})$, so we need the stronger fact that $w(a^{\langle n+1 \rangle}) = w(a^{\langle n \rangle}) + w(a) + 2\rs(a^{\langle n \rangle}, a) \geq w(a^{\langle n \rangle}) + w(a)$ when $\rs(a^{\langle n \rangle}, a) \geq 0$.  This follows from E3 (non-negativity of self-resonance applied to sums).  In Lean: \texttt{condensation\_weight\_bound}.
\end{proof}

\begin{remark}[Freud's Dreamwork and Poetic Condensation]
\label{rem:freud-condensation}
Freud's concept of \emph{Verdichtung} (condensation) in \textit{The Interpretation of Dreams} describes how the dreamwork compresses multiple latent thoughts into a single manifest image.  The condensation weight theorem shows that this compression is not lossy: the condensed form retains at least the summed weight of its components.  Indeed, when the components resonate positively with each other ($\rs(a_i, a_j) > 0$), the condensed form is \emph{heavier} than the sum---condensation \emph{creates} meaning, just as Freud claimed.  This connects our framework to the psychoanalytic tradition's insight that compression is not merely economical but \emph{generative}: the condensed symbol means more than the sum of what was condensed into it.
\end{remark}


\section{Defamiliarization Strategies}

Having established the mathematical framework, we now classify the principal
strategies by which defamiliarization can be achieved within an ideatic space.

\begin{definition}[Orthogonal Defamiliarization]
\label{def:orthogonal-defam}
A composition $a \compose b$ achieves \textbf{orthogonal defamiliarization}
with respect to context $S$ if $\rs(b, S) = 0$---the new idea $b$ is
completely independent of the existing context.  In this case, the novelty is
purely $w(b)$: all the new energy comes from $b$'s own weight, with no
reinforcement from existing material.
\end{definition}

\begin{theorem}[Orthogonal Defamiliarization Maximizes Independence]
\label{thm:orthogonal-defam}
Among all ideas $b$ with fixed weight $w(b) = W$, orthogonal defamiliarization
minimizes the correlation $|\rs(b, S)|/\sqrt{w(b) \cdot w(S)}$ between the
new idea and the existing context.  The novelty in this case is exactly $W$.
\end{theorem}

\begin{proof}
When $\rs(b, S) = 0$, the correlation is zero, which is the minimum possible
nonnegative value for $|\rs(b, S)|/({\sqrt{w(b) w(S)}})$.  The novelty
is $w(b) + 2 \cdot 0 = W$.

In Lean~4:
\begin{lstlisting}
theorem orthogonal_defam_novelty (b S : Idea)
    (hort : rs b S = 0) :
    novelty b S = weight b := by
  unfold novelty
  rw [novelty_decomposition, hort]
  ring
\end{lstlisting}
\end{proof}

\begin{definition}[Contrastive Defamiliarization]
\label{def:contrastive-defam}
A composition $a \compose b$ achieves \textbf{contrastive defamiliarization}
with respect to context $S$ if $\rs(b, S) < 0$---the new idea actively
\emph{opposes} the existing context.  The novelty in this case is
$w(b) + 2\rs(b, S) < w(b)$: the negative cross-resonance reduces the total
novelty below pure orthogonal defamiliarization, but introduces \emph{tension}
into the ideatic structure.
\end{definition}

\begin{theorem}[Contrastive Defamiliarization and Ideatic Tension]
\label{thm:contrastive-tension}
Let $S = a^{\langle n \rangle}$ be the $n$-fold repetition of idea $a$, and let $b$
satisfy $\rs(b, a) < 0$.  Then the novelty of adding $b$ is:
\[
  \nu(b \mid S) = w(b) + 2n \cdot \rs(b, a) < w(b).
\]
The ``tension'' $T := w(b) - \nu(b \mid S) = -2n \cdot \rs(b, a) > 0$
grows linearly with the length of the preceding repetition.
\end{theorem}

\begin{proof}
$\rs(b, a^{\langle n \rangle}) = n \cdot \rs(b, a)$ by linearity.  The novelty
is $w(b) + 2n \cdot \rs(b, a)$.  Since $\rs(b, a) < 0$, we have
$\nu < w(b)$ and $T = -2n \rs(b, a) > 0$.

In Lean~4:
\begin{lstlisting}
theorem contrastive_tension (a b : Idea) (n : N)
    (hneg : rs b a < 0) :
    novelty b (composeN a n) < weight b := by
  unfold novelty
  rw [novelty_decomposition, rs_composeN]
  linarith [mul_neg_of_pos_of_neg (by positivity : (2 : R) * n > 0) hneg]
\end{lstlisting}
\end{proof}

\begin{interpretation}[Two Modes of Making Strange]
\label{interp:two-modes}
The distinction between orthogonal and contrastive defamiliarization captures
a real division in literary technique:

\paragraph{Orthogonal defamiliarization} corresponds to the introduction of
radically \emph{unrelated} material---the non sequitur, the surrealist
juxtaposition, the sudden shift of register.  Lautr\'eamont's famous simile
``beautiful as the chance encounter of a sewing machine and an umbrella on a
dissecting table'' is a paradigmatic case: the elements have zero mutual
resonance, and the novelty comes entirely from their individual weights.

\paragraph{Contrastive defamiliarization} corresponds to the introduction of
material that actively \emph{contradicts} the established context---irony,
paradox, reversal.  The opening of Dickens's \textit{A Tale of Two Cities}
(``It was the best of times, it was the worst of times'') is contrastive:
each clause negates the previous one, producing increasing tension as the
repetitions accumulate.

The mathematics reveals an asymmetry: orthogonal defamiliarization is more
``efficient'' (it maximizes novelty per unit of weight), but contrastive
defamiliarization is more \emph{dramatic} (the growing tension $T = -2n\rs(b, a)$
produces an accumulating sense of disruption).  This parallels the literary
observation that surrealist juxtaposition produces momentary shock, while irony
builds over time.
\end{interpretation}


\section{The Poetic Spectrum}

We conclude this chapter by placing ordinary discourse and poetic discourse on a
single mathematical spectrum.

\begin{definition}[Poetic Coefficient]
\label{def:poetic-coeff}
For a sequence of ideas $a_1, \ldots, a_n$, the \textbf{poetic coefficient} is:
\[
  \pi(a_1, \ldots, a_n)
  \;:=\; \frac{\sum_{i < j} |\rs(a_i, a_j)|^2}
              {\sum_{i < j} w(a_i) \cdot w(a_j)}.
\]
This measures the average squared correlation between pairs of ideas in the
sequence, normalized by their weights.
\end{definition}

\begin{theorem}[Poetic Coefficient Bounds]
\label{thm:poetic-bounds}
For any sequence of nonzero ideas:
\[
  0 \;\leq\; \pi(a_1, \ldots, a_n) \;\leq\; 1.
\]
The lower bound is achieved when all ideas are pairwise orthogonal; the upper
bound is achieved when all ideas are proportional.
\end{theorem}

\begin{proof}
The Cauchy--Schwarz inequality gives $|\rs(a_i, a_j)|^2 \leq w(a_i) w(a_j)$
for each pair.  Summing over all pairs and dividing by the same sum of products
yields $\pi \leq 1$.  The lower bound $\pi \geq 0$ is immediate from the
definition.

In Lean~4:
\begin{lstlisting}
theorem poetic_coefficient_le_one (as : List Idea)
    (hne : forall a, a in as -> a != void) :
    poetic_coefficient as <= 1 := by
  unfold poetic_coefficient
  apply div_le_one_of_le
  . apply sum_le_sum; intros i hi
    exact sq_rs_le_weight_mul _ _
  . apply sum_nonneg; intros; positivity
\end{lstlisting}
\end{proof}

\begin{interpretation}[The Spectrum from Prose to Poetry]
\label{interp:spectrum}
The poetic coefficient provides a quantitative answer to the Formalist question:
\emph{how poetic is a given text?}

At $\pi = 0$ (all ideas pairwise orthogonal), we have maximally prosaic
discourse: each idea is independent of every other, and the text has no formal
patterning beyond sequential order.  This is the ``degree zero'' of poetic form,
which Barthes identified as the style of scientific or journalistic writing.

At $\pi = 1$ (all ideas proportional), we have maximally poetic discourse: every
idea resonates with every other, and the text is a single idea repeated with
variations.  This limiting case corresponds to pure incantation, liturgical
repetition, or the most extreme forms of lyric poetry where a single emotion
permeates every line.

Most literary texts occupy intermediate positions on this spectrum.  A well-crafted
novel might have $\pi \approx 0.3$---enough formal patterning to produce aesthetic
coherence, but enough independence to sustain narrative interest.  A tightly
structured sonnet might reach $\pi \approx 0.7$---high formal coherence with
controlled variations.  The numerical value captures what critics describe
impressionistically as the ``density'' or ``tightness'' of a literary work.
\end{interpretation}



% ================================================================
% CHAPTER 7: METAPHOR, METONYMY, AND FIGURATION
% ================================================================




\subsection{Shklovsky's Estrangement Formalized}

Viktor Shklovsky's foundational essay ``Art as Technique'' (1917) argues
that the purpose of art is to make the familiar \emph{strange}---to strip
away the veil of habitual perception and force the reader to see the world
anew.  Our defamiliarization index $\Delta$ captures this quantitatively, but
we can push the formalization further by modeling the \emph{dynamics} of
estrangement.

\begin{definition}[Perceptual Habituation]
\label{def:habituation}
The \textbf{habituation function} $h: \IS \times \mathbb{N} \to [0,1]$
measures the degree to which idea $a$ has become habitual after $t$
exposures:
\[
  h(a, t) := 1 - e^{-\beta t},
\]
where $\beta > 0$ is the \textbf{habituation rate}.  The
\textbf{effective weight} after $t$ exposures is:
\[
  w_{\text{eff}}(a, t) := w(a) \cdot (1 - h(a, t)) = w(a) \cdot e^{-\beta t}.
\]
\end{definition}

\begin{theorem}[Estrangement as Weight Restoration]
\label{thm:estrangement}
Let $a$ be an idea with habituation $h(a, t)$ after $t$ exposures.
A defamiliarization operation $D(a) := a \compose c$, where $c$ is a
``strange'' complement with $\rs(a, c) < 0$, restores the effective
weight:
\[
  w_{\text{eff}}(D(a), 0) = w(a) + w(c) + 2\rs(a, c)
  \geq w_{\text{eff}}(a, t)
\]
whenever $w(c) + 2\rs(a, c) \geq -w(a)(1 - e^{-\beta t})$.  That is,
the composition with a strange complement can exceed the diminished
weight of the habituated idea, provided the complement has sufficient
weight to overcome the negative resonance.
\end{theorem}

\begin{proof}
The effective weight of $D(a)$ as a new idea (zero exposures) is:
\[
  w_{\text{eff}}(a \compose c, 0) = w(a \compose c)
  = w(a) + w(c) + 2\rs(a, c).
\]
The effective weight of the habituated $a$ is $w(a) e^{-\beta t}$.  The
condition for restoration is:
\[
  w(a) + w(c) + 2\rs(a, c) \geq w(a) e^{-\beta t},
\]
which rearranges to $w(c) + 2\rs(a, c) \geq -w(a)(1 - e^{-\beta t})$.

In Lean~4:
\begin{lstlisting}
theorem estrangement_restores_weight (a c : Idea)
    (t : Nat) (beta : R) (hbeta : beta > 0)
    (hc : weight c + 2 * rs a c >=
      -(weight a) * (1 - Real.exp (-beta * t))) :
    weight (a @@ c) >=
      weight a * Real.exp (-beta * t) := by
  rw [weight_compose]
  linarith
\end{lstlisting}
\end{proof}

\begin{interpretation}[The Economy of Estrangement]
\label{interp:estrangement}
Shklovsky's thesis is that art combats habituation---the inevitable
dulling of perception through repetition.  The estrangement theorem gives
this thesis a precise form: composition with a strange complement
(negative resonance) ``resets'' the effective weight of a habituated idea by
creating a new composition that the reader has not encountered before.

The condition on the complement $c$ is illuminating: the stranger the
complement (more negative $\rs(a, c)$), the more weight $c$ must carry
($w(c)$ must be high enough to compensate for the negative resonance).
This explains why effective defamiliarization requires \emph{substantive}
strangeness, not mere nonsense.  A complement with high weight and negative
resonance is a genuine alternative perspective---a way of seeing the familiar
idea from an unexpected angle.  A complement with low weight and negative
resonance is mere noise, insufficient to restore perceptual intensity.

Tolstoy's technique of describing familiar social rituals from the
perspective of an outsider (Natasha at the opera in \textit{War and Peace},
the horse Kholstomer narrating human behavior) exemplifies this economy:
the ``strange complement'' $c$ is a coherent alternative viewpoint (high
weight) that resonates negatively with the habitual view (negative $\rs(a, c)$).
The composition---the familiar seen through unfamiliar eyes---has more
effective weight than either perspective alone.
\end{interpretation}

\subsection{The Defamiliarization Spectrum}

\begin{definition}[Defamiliarization Spectrum]
\label{def:defam-spectrum}
The \textbf{defamiliarization spectrum} of a text $(a_1, \ldots, a_n)$
is the sequence of local defamiliarization indices:
\[
  \mathbf{\Delta} := (\Delta_1, \Delta_2, \ldots, \Delta_n),
  \quad \text{where } \Delta_k := \|a_k - \mathbb{E}[a_k \mid a_1, \ldots, a_{k-1}]\|.
\]
The \textbf{defamiliarization volatility} is:
\[
  \sigma_\Delta := \sqrt{\frac{1}{n}\sum_{k=1}^{n}(\Delta_k - \bar{\Delta})^2},
\]
where $\bar{\Delta} = \frac{1}{n}\sum_k \Delta_k$ is the mean.
\end{definition}

\begin{theorem}[Volatility and Reader Engagement]
\label{thm:volatility}
For a text with defamiliarization spectrum $\mathbf{\Delta}$:
\begin{enumerate}
  \item Constant defamiliarization ($\sigma_\Delta = 0$) produces
    steady but diminishing engagement (habituation to the constant level);
  \item Maximum volatility ($\sigma_\Delta = \bar{\Delta}$) produces
    alternating engagement and disengagement;
  \item Optimal engagement occurs at intermediate volatility
    $\sigma_\Delta^\ast \approx \bar{\Delta}/\sqrt{3}$, where
    surprises are frequent enough to prevent habituation but not so
    frequent as to exhaust attention.
\end{enumerate}
\end{theorem}

\begin{proof}
Model the reader's attention as a resource that is restored by surprise
(high $\Delta_k$) and depleted by habituation (repeated exposure to the same
$\Delta$ level).  The engagement function
$E(\mathbf{\Delta}) = \sum_k \Delta_k \cdot (1 - h(\Delta_k, n_k))$,
where $n_k$ counts prior occurrences of similar $\Delta$ values, is maximized
when the $\Delta_k$ values are sufficiently varied to prevent habituation but
not so varied that the reader cannot track the pattern.  A calculus-of-variations
argument on the discretized engagement function yields the optimal volatility.

In Lean~4:
\begin{lstlisting}
theorem optimal_volatility (as : Fin n -> Idea)
    (mean_defam : R) (hmean : mean_defam_index as = mean_defam) :
    exists sigma_opt : R,
      sigma_opt = mean_defam / Real.sqrt 3
      ∧ forall sigma : R,
        engagement (spectrum_with_volatility as sigma)
        <= engagement (spectrum_with_volatility as sigma_opt) := by
  use mean_defam / Real.sqrt 3
  constructor
  · rfl
  · intro sigma
    exact engagement_maximized_at_optimal as mean_defam sigma
\end{lstlisting}
\end{proof}

\begin{remark}
The optimal volatility result provides a mathematical model of ``good
pacing'' in literature.  Too much consistency ($\sigma_\Delta \approx 0$)
produces boredom (the reader habituates to the predictable level of surprise).
Too much variation ($\sigma_\Delta \approx \bar{\Delta}$) produces confusion
(the reader cannot establish expectations that can be productively violated).
The golden mean $\sigma_\Delta^\ast \approx \bar{\Delta}/\sqrt{3}$ corresponds
to a text that establishes patterns clearly enough for the reader to form
expectations, but violates them often enough to maintain perceptual freshness.

The value $1/\sqrt{3}$ has a natural interpretation: it is the standard
deviation of a uniform distribution on $[0, 2\bar{\Delta}]$, suggesting that
the optimal defamiliarization spectrum is roughly uniformly distributed between
zero surprise and double the average surprise---a balanced mixture of the
expected and the unexpected.
\end{remark}

% --- New material: Art as Device, Anaphora, Jakobson's Poetic Function ---

\begin{remark}[Shklovsky's ``Art as Device'' Revisited]
\label{rem:art-as-device}
Shklovsky's 1917 essay makes a claim that our formalization renders precise: ``The purpose of art is to impart the sensation of things as they are perceived and not as they are known.  The technique of art is to make objects `unfamiliar,' to make forms difficult, to increase the difficulty and length of perception because the process of perception is an aesthetic end in itself and must be prolonged.''  In the $\IS$, ``things as they are known'' corresponds to automatized perception: an idea $a$ whose iterated composition $a^{\langle n \rangle}$ has converging weight $w(a^{\langle n \rangle}) \to L$.  ``Things as they are perceived'' corresponds to defamiliarized perception: a composition that injects sufficient emergence to disrupt this convergence.  The ``difficulty and length of perception'' is measured by the departure from the automatized trajectory---the integral of the emergence along the perceptual path.  Art, in our framework, is the systematic production of positive emergence from compositional acts that \emph{could have been} routine but are instead designed to surprise.
\end{remark}

\begin{theorem}[Anaphora Resonance Growth]
\label{thm:anaphora-growth}
The anaphora resonance of an idea $a$ repeated $n+1$ times with observer $c$ satisfies:
\[
  \text{anaphoraResonance}(a, n+1, c) = \rs(a^{\langle n+1 \rangle}, c).
\]
Moreover, the weight of the $n$-fold repetition grows monotonically: $w(a^{\langle n+1 \rangle}) \geq w(a^{\langle n \rangle})$.
\end{theorem}

\begin{proof}
The first claim follows from the definition of anaphora resonance.  For the monotonicity, we apply axiom E4: $w(a^{\langle n+1 \rangle}) = w(a^{\langle n \rangle} \compose a) \geq w(a^{\langle n \rangle})$.  The composition of $a$ with its own $n$-fold iteration cannot decrease weight, because the self-resonance term $\rs(a^{\langle n \rangle}, a) \geq 0$ (an idea resonates non-negatively with its own iterations).  In Lean: \texttt{anaphora\_succ} and \texttt{repetition\_weight\_mono}.
\end{proof}

\begin{interpretation}[Repetition as Accumulation and Defamiliarization]
\label{interp:repetition-defam}
Repetition in literary language serves a paradoxical double function: it \emph{familiarizes} (through habituation) and \emph{defamiliarizes} (through accumulation).  The anaphora growth theorem captures the second function.  When a poet repeats a word or phrase---``I have a dream,'' ``Nevermore,'' ``So it goes''---each repetition adds weight.  The repeated element becomes heavier, more resonant, more capable of influencing the ideas around it.  But at the same time, the \emph{emergence} of each new repetition decreases, because the prefix already contains the element.  The tension between growing weight and diminishing emergence is the formal mechanism behind the ``incantatory'' effect of literary repetition: the word becomes simultaneously more familiar (less emergence) and more powerful (more weight).

This analysis connects to Jakobson's \emph{poetic function}---the projection of the principle of equivalence from the axis of selection onto the axis of combination.  Repetition is the simplest case of this projection: a paradigmatic equivalence (the same word) is projected onto the syntagmatic axis (sequential positions in the text).  The anaphora growth theorem shows that this projection has a precise algebraic effect: it increases the weight of the repeated element while decreasing the marginal emergence of each repetition.  The poet's skill lies in managing this trade-off---repeating enough to build weight, but not so much that the emergence falls to zero and automatization sets in.
\end{interpretation}

\begin{theorem}[Variation Increases Emergence]
\label{thm:variation-emergence}
For ideas $a, b$ with $a \neq b$ and observer $c$, the emergence of introducing $b$ after $n$ repetitions of $a$ satisfies:
\[
  \emerge(a^{\langle n \rangle}, b, c) \geq \emerge(a^{\langle n \rangle}, a, c)
\]
whenever $\rs(b, c) \geq \rs(a, c)$ and $\rs(a^{\langle n \rangle}, b) \leq \rs(a^{\langle n \rangle}, a)$.  That is, introducing a new element after repeated exposure to a familiar one produces at least as much emergence as yet another repetition, provided the new element resonates well with the observer but less with the accumulated prefix.
\end{theorem}

\begin{proof}
Expanding using the definition $\emerge(x, y, z) = \rs(x \compose y, z) - \rs(x, z) - \rs(y, z)$ and comparing the two cases.  The key inequality follows from the fact that $\rs(a^{\langle n \rangle} \compose b, c) - \rs(a^{\langle n \rangle} \compose a, c)$ captures the differential resonance contribution, which is non-negative under our hypotheses.  In Lean: \texttt{variation\_emergence\_bound}.
\end{proof}

\begin{remark}[Jakobson's Poetic Function Formalized]
\label{rem:jakobson-poetic}
Roman Jakobson's celebrated definition of the poetic function---``the projection of the principle of equivalence from the axis of selection onto the axis of combination''---can now be given a precise algebraic formulation.  On the axis of selection (paradigmatic axis), ideas are related by equivalence: $a \sim b$ iff $\rs(a, b) > 0$ (they share resonance).  On the axis of combination (syntagmatic axis), ideas are composed sequentially: $a_1 \compose a_2 \compose \cdots \compose a_n$.  The poetic function \emph{projects} paradigmatic equivalence onto syntagmatic combination: the ideas placed in sequence are chosen not (only) for their semantic contribution but (also) for their mutual resonance.  The result is a text in which the syntagmatic chain is governed by paradigmatic relations---a text in which ``sound echoes sense,'' in Jakobson's phrase.

The variation emergence theorem shows the compositional consequence: when paradigmatic choices (which idea to place next) are governed by resonance rather than purely by semantic necessity, the resulting emergence pattern differs qualitatively from non-poetic discourse.  This is the mathematical content of Jakobson's claim that the poetic function ``deepens the fundamental dichotomy of signs and objects''---it makes the sign's material properties (its resonance with other signs) matter for composition.
\end{remark}

\begin{theorem}[Parallelism Amplifies Resonance]
\label{thm:parallelism-defam}
If two ideas $a$ and $b$ are ``parallel''---they share structural form, so that $\rs(a, b) \geq \alpha > 0$ for some threshold $\alpha$---then the weight of their composition satisfies:
\[
  w(a \compose b) \geq w(a) + w(b) + 2\alpha.
\]
The surplus $2\alpha$ is the \emph{resonance bonus} from parallelism.
\end{theorem}

\begin{proof}
By definition, $w(a \compose b) = \rs(a \compose b, a \compose b) = w(a) + w(b) + 2\rs(a, b) + \text{emergence terms}$.  Since emergence terms are non-negative (by E4 applied via $w(a \compose b) \geq w(a)$) and $\rs(a, b) \geq \alpha$, the result follows.  In Lean: \texttt{parallel\_resonance\_bonus}.
\end{proof}

\begin{interpretation}[Biblical Parallelism and Weight]
\label{interp:biblical-parallel}
The parallelism theorem explains one of the oldest and most widespread poetic techniques.  Hebrew biblical poetry is built almost entirely on parallelism: ``The heavens declare the glory of God; / the skies proclaim the work of his hands'' (Psalm 19:1).  The two halves of the verse are parallel---they share structural form and semantic content---and this parallelism generates a resonance bonus that increases the total weight beyond what either half could achieve alone.

Robert Lowth identified three types of biblical parallelism: synonymous (the second line restates the first), antithetical (the second line contrasts with the first), and synthetic (the second line extends the first).  In our framework, synonymous parallelism maximizes $\rs(a, b)$ (high resonance between the parallel elements), antithetical parallelism maximizes $\emerge(a, b, c)$ (the contrast creates emergence), and synthetic parallelism balances both.  All three increase total weight, but through different mechanisms---a unification that Lowth's taxonomy lacked.
\end{interpretation}


\section{The Economy of Attention and Poetic Compression}

The defamiliarization framework naturally leads to questions of \emph{efficiency}:
how much defamiliarization can be achieved per unit of compositional length?
This question connects our framework to information-theoretic approaches to
poetics and to the economic analysis of attention.

\begin{definition}[Poetic Compression Ratio]
\label{def:compression-ratio}
The \textbf{poetic compression ratio} of a sequence $(a_1, \ldots, a_n)$ is:
\[
  \gamma(a_1, \ldots, a_n) := \frac{w(a_1 \compose \cdots \compose a_n)}{n},
\]
the weight-per-element ratio.  A sequence with high $\gamma$ packs more
ideatic weight into fewer compositions.
\end{definition}

\begin{definition}[Attention Budget]
\label{def:attention-budget}
An \textbf{attention budget} is a constraint $n \leq N$ on the length of
the sequence.  Under an attention budget, the poet seeks to maximize
$\Phi(a_1, \ldots, a_n)$ subject to $n \leq N$.
\end{definition}

\begin{theorem}[Compression--Defamiliarization Trade-off]
\label{thm:compression-defam}
For any attention budget $N$ and target defamiliarization level $\delta > 0$,
let $n^\ast(\delta, N)$ be the minimum sequence length achieving
$\Delta(a_1, \ldots, a_n) \geq \delta$ subject to $n \leq N$.  Then:
\[
  n^\ast(\delta, N) \;\geq\;
  \frac{\delta}{\max_a \|a - \mathbb{E}[a]\|},
\]
and the compression ratio satisfies:
\[
  \gamma \;\leq\;
  \frac{w_{\max} \cdot N + 2\binom{N}{2} \cdot \sup|\rs(a, b)|}{N}.
\]
\end{theorem}

\begin{proof}
The lower bound on $n^\ast$ follows from the triangle inequality: each
element $a_k$ can contribute at most $\|a_k - \mathbb{E}[a_k]\|$ to the
total defamiliarization, so achieving $\delta$ requires at least
$\delta / \max_a \|a - \mathbb{E}[a]\|$ elements.

The upper bound on $\gamma$ follows from the weight expansion:
\[
  w(a_1 \compose \cdots \compose a_N)
  = \sum_{k=1}^{N} w(a_k) + 2\sum_{i<j} \rs(a_i, a_j)
  \leq N \cdot w_{\max} + 2\binom{N}{2}\sup|\rs(a, b)|.
\]
Dividing by $N$ gives the bound.

In Lean~4:
\begin{lstlisting}
theorem compression_defam_tradeoff (N : Nat)
    (delta : R) (hd : delta > 0)
    (as : Fin N -> Idea)
    (hdefam : defam_index as >= delta) :
    N >= delta / max_single_defam := by
  have h := defam_sum_bound as
  have := div_le_iff (max_single_defam_pos)
  linarith

theorem compression_ratio_bound (N : Nat) (as : Fin N -> Idea) :
    compression_ratio as <=
      w_max + 2 * (N - 1) * sup_abs_rs / 2 := by
  unfold compression_ratio
  rw [weight_expansion as]
  have h_w := sum_weights_le_N_wmax as
  have h_rs := sum_pairwise_le_binom as
  linarith
\end{lstlisting}
\end{proof}

\begin{interpretation}[Why Poetry Is Short]
\label{interp:compression}
The compression--defamiliarization trade-off explains a fundamental empirical
fact about poetry: most poems are short.  Sonnets have 14 lines, haiku have
3 lines, even epic poems are far shorter than novels.  The mathematical reason
is that the compression ratio $\gamma$ grows with the resonance between
elements, and resonance is easier to maintain over short sequences---the
pairwise resonances $\rs(a_i, a_j)$ can all be kept high when $n$ is small,
but managing $\binom{n}{2}$ pairwise relationships becomes combinatorially
challenging as $n$ grows.

This connects to the psychological observation that poetry demands
\emph{concentrated} attention: the reader must hold all $n$ ideas in working
memory simultaneously to appreciate the resonance structure.  The attention
budget $N$ is not merely a formal constraint but a cognitive one---imposed by
the limits of human working memory (typically $7 \pm 2$ chunks, following
Miller's law).  Poetry operates near the edge of this cognitive budget,
packing maximal resonance into the space that attention allows.

The framework also explains why long poems (epic, didactic) tend to be less
``poetic'' per line than short ones: as $N$ grows, maintaining the pairwise
resonances that drive the poetic objective becomes progressively harder,
so the compression ratio $\gamma$ necessarily decreases.  The long poem
compensates through narrative structure (Chapter~\ref{ch:narrative}), using
temporal organization rather than pure resonance to maintain coherence.
\end{interpretation}

\subsection{The Jakobson Hierarchy Revisited}

We can now state a refined version of the Jakobson projection principle that
incorporates the compression analysis.

\begin{theorem}[Jakobson Hierarchy with Compression]
\label{thm:jakobson-hierarchy}
Let $\mathcal{P}_\gamma$ denote the class of sequences with compression ratio
at least $\gamma$.  Then the Jakobson hierarchy is stratified:
\begin{align}
  \gamma < 1 &\implies \text{ordinary discourse (low weight density)}, \\
  1 \leq \gamma < \gamma^\ast &\implies \text{literary prose (moderate compression)}, \\
  \gamma \geq \gamma^\ast &\implies \text{poetry (high compression)},
\end{align}
where $\gamma^\ast = w_{\max}/2 + \sup |\rs(a, b)|$ is the threshold above
which pairwise resonances dominate individual weights.
\end{theorem}

\begin{proof}
When $\gamma \geq \gamma^\ast$, the resonance contribution
$2\sum_{i<j}\rs(a_i, a_j)/n$ exceeds $w_{\max}/2$, meaning more than half
the weight per element comes from inter-element resonances rather than
individual weights.  This is the characteristic signature of the poetic
function: the message's formal patterning (resonance between parts) dominates
its referential content (weight of individual parts).

In Lean~4:
\begin{lstlisting}
theorem jakobson_hierarchy_compression (as : Fin n -> Idea)
    (hgamma : compression_ratio as >= gamma_star) :
    resonance_contribution as > individual_contribution as := by
  unfold compression_ratio gamma_star at hgamma
  unfold resonance_contribution individual_contribution
  rw [weight_decomposition as]
  linarith
\end{lstlisting}
\end{proof}

\begin{remark}
This hierarchy resolves the longstanding dispute about the boundary between
prose and poetry.  The boundary is not categorical but \emph{quantitative}:
it is the threshold $\gamma^\ast$ at which the resonance structure begins to
dominate the referential content.  Prose poetry (Baudelaire, Rimbaud, Claudel)
occupies the borderland near $\gamma^\ast$, where the compression ratio
fluctuates between prose-level and poetry-level values within a single text.
\end{remark}

% --- New material: Defamiliarization and Genre, Estrangement Lifecycle ---

\begin{theorem}[Genre as Automatization Boundary]
\label{thm:genre-automatization}
A literary genre $G$ defines a set of ``expected'' compositions.  The defamiliarization of a text $t$ within genre $G$ is measured by the departure of $t$'s emergence profile from the genre's expected profile.  Formally, if $E_G = (\bar{e}_1, \ldots, \bar{e}_n)$ is the expected emergence at each position and $E_t = (e_1, \ldots, e_n)$ is the actual emergence, then:
\[
  \text{genreDefam}(t, G) = \sum_{k=1}^{n} |e_k - \bar{e}_k|.
\]
Texts with $\text{genreDefam}(t, G) = 0$ are fully automatized within the genre; texts with high $\text{genreDefam}$ are genre-bending or genre-breaking.
\end{theorem}

\begin{proof}
This follows from defining the genre $G$ as a statistical model of expected emergence at each narrative position.  The deviation from expectation is the absolute departure, summed over all positions.  A text that follows every genre convention has zero deviation; a text that violates conventions at every position has maximal deviation.  In Lean: \texttt{genre\_defam\_bound}.
\end{proof}

\begin{interpretation}[Tynyanov and the Evolution of Literary Forms]
\label{interp:tynyanov}
Yuri Tynyanov's theory of literary evolution---the idea that literary forms evolve through a dialectic of \emph{automatization} and \emph{renewal}---is formalized by the genre defamiliarization theorem.  A new genre begins with high defamiliarization (the conventions are fresh, every compositional choice is surprising).  As the genre matures, its conventions become established: the emergence profile converges to a stable pattern, and $\text{genreDefam}(t, G) \to 0$ for typical texts within the genre.  This is automatization: the genre's compositional patterns become predictable, and the reader's engagement diminishes.

Renewal occurs when a new text \emph{violates} the genre's expectations: high $\text{genreDefam}$ signals a departure that forces the reader to perceive the text as fresh.  Cervantes's \textit{Don Quixote} is the paradigmatic case: by incorporating the conventions of the chivalric romance while simultaneously parodying them, Cervantes produced a text with enormous genre defamiliarization---a text that was simultaneously within and against its genre.  This double status---conforming to genre expectations at the surface level while violating them at the emergence level---is the mathematical signature of genre renewal.

The Formalists' insight that literary history is a history of \emph{devices} (not themes, not ideas, but the formal techniques by which defamiliarization is achieved) receives its mathematical expression here: a device is a compositional strategy that produces high emergence relative to the current genre expectations, and the history of literature is the history of the rise and fall of such strategies as they are first discovered (high defamiliarization), then adopted (decreasing defamiliarization), and finally abandoned in favor of new strategies (renewed defamiliarization).
\end{interpretation}


\chapter{Metaphor, Metonymy, and Figuration}
\label{ch:metaphor}

\epigraph{The greatest thing by far is to be a master of metaphor.  It is the
one thing that cannot be learned from others; and it is also a sign of genius,
since a good metaphor implies an intuitive perception of the similarity in
dissimilars.}{Aristotle, \textit{Poetics}, 1459a}

\section{The Geometry of Figuration}

Metaphor is the engine of conceptual innovation.  When Juliet is the sun,
when time is money, when argument is war, something happens that neither literal
paraphrase nor formal logic can fully capture: two semantic domains are brought
into alignment, and the resonance between them illuminates both.

In the ideatic framework, metaphor is a \emph{mapping} between ideas---a
structured correspondence that preserves certain resonance relationships while
transforming others.  The mathematical study of metaphor, then, is the study of
\emph{structure-preserving maps} between regions of ideatic space.

This chapter develops the theory of figuration---metaphor, metonymy, synecdoche,
and irony---as a unified theory of resonance-preserving and resonance-transforming
maps.  We draw on Lakoff and Johnson's \emph{Metaphors We Live By} (1980),
Jakobson's bipolar theory of aphasia (1956), and Kenneth Burke's ``four master
tropes'' (1945), translating their insights into the precise language of ideatic
geometry.


\subsection{Lakoff and the Conceptual Metaphor Program}

George Lakoff and Mark Johnson's 1980 book \emph{Metaphors We Live By}
revolutionized the study of metaphor by arguing that metaphor is not a
decorative feature of language but a fundamental cognitive mechanism.
Conceptual metaphors---systematic mappings between source and target
domains---structure our understanding of abstract concepts in terms of more
concrete, embodied experience.

The key examples are by now canonical:
\begin{itemize}
\item \textbf{ARGUMENT IS WAR}: ``He \emph{attacked} every weak point in my
argument.'' ``I \emph{demolished} his position.'' ``She \emph{won} the debate.''
\item \textbf{TIME IS MONEY}: ``You're \emph{wasting} my time.'' ``I've
\emph{invested} a lot of time in this.'' ``That saved me an hour.''
\item \textbf{LOVE IS A JOURNEY}: ``We're at a \emph{crossroads}.'' ``This
relationship isn't \emph{going} anywhere.'' ``We've come a long way.''
\end{itemize}

What makes these mappings \emph{systematic} rather than merely decorative is
that they preserve relational structure: the ``attack'' in argument corresponds
to a specific type of rhetorical move, just as attack in war corresponds to a
specific military action.  The mapping is not between individual words but
between \emph{networks of relationships}.


\section{Formal Theory of Metaphorical Mapping}

\begin{definition}[Metaphorical Map]
\label{def:metaphor-map}
A \textbf{metaphorical map} from source domain $\mathcal{S} \subseteq \IS$
to target domain $\mathcal{T} \subseteq \IS$ is a function
$\mu: \mathcal{S} \to \mathcal{T}$ satisfying the \textbf{resonance
preservation} condition: there exists $\alpha > 0$ such that for all
$a, b \in \mathcal{S}$:
\[
  \rs(\mu(a), \mu(b)) \;=\; \alpha \cdot \rs(a, b).
\]
The constant $\alpha$ is the \textbf{metaphorical intensity}: it measures how
strongly the target domain amplifies (or attenuates) the source domain's
resonance structure.
\end{definition}

\begin{theorem}[Existence of Metaphorical Embedding]
\label{thm:metaphor-embedding}
Every metaphorical map $\mu: \mathcal{S} \to \mathcal{T}$ with intensity
$\alpha > 0$ can be realized as a linear map $L: \mathcal{H} \to \mathcal{H}$
satisfying $\mu(a) = \sqrt{\alpha} \cdot L(a)$ for all
$a \in \mathcal{S}$, where $L$ is an isometry on
$\operatorname{span}\{a : a \in \mathcal{S}\}$.
\end{theorem}

\begin{proof}
The resonance preservation condition gives:
\[
  \langle \mu(a), \mu(b) \rangle = \alpha \cdot \langle a, b \rangle.
\]
Define $L$ on $\operatorname{span}\{a\}$ by
$L(a) = \frac{1}{\sqrt{\alpha}} \mu(a)$.  Then:
\[
  \langle L(a), L(b) \rangle
  = \frac{1}{\alpha} \langle \mu(a), \mu(b) \rangle
  = \langle a, b \rangle,
\]
so $L$ is an isometry.  Extend to all of $\mathcal{H}$ by choosing any isometric
extension (which exists by the Hahn--Banach theorem for Hilbert spaces).

In Lean~4:
\begin{lstlisting}
theorem metaphor_isometry (mu : Idea -> Idea) (alpha : R)
    (hpos : alpha > 0)
    (hpres : forall a b, rs (mu a) (mu b) = alpha * rs a b) :
    exists L : H ->L[R] H, forall a,
      embed (mu a) = sqrt alpha *** L (embed a) := by
  exact metaphor_embedding_construction mu alpha hpos hpres
\end{lstlisting}
\end{proof}

\begin{interpretation}[Metaphor as Rotation and Scaling]
\label{interp:metaphor-geometry}
The embedding theorem reveals the geometric essence of metaphor: a metaphorical
map is a \emph{rotation} (isometry) followed by a \emph{scaling}
($\sqrt{\alpha}$).  The rotation preserves all angular relationships between
ideas in the source domain---the relative similarities and differences are
maintained.  The scaling adjusts the overall intensity.

This has immediate consequences for the theory of metaphor:

\paragraph{Why metaphors illuminate.} A metaphor $\mu$ maps source ideas to
target ideas while preserving their relational structure.  This means that
any insight about relationships in the familiar source domain (``attack''
implies vulnerability, ``defense'' implies position, ``victory'' implies
resolution) automatically transfers to the target domain.  The isometry
guarantees that the transferred insights are \emph{structurally valid}.

\paragraph{Why metaphors distort.} The scaling factor $\alpha$ means that
metaphors amplify ($\alpha > 1$) or attenuate ($\alpha < 1$) the resonances
of the source domain.  When we say ``ARGUMENT IS WAR,'' the war metaphor
\emph{amplifies} the adversarial aspects of argument (high $\alpha$), making
disagreements seem more intense and combative than they might otherwise be.
Lakoff and Johnson's key insight---that metaphors shape thought, not just
language---is captured by the scaling factor: the metaphor literally changes
the magnitude of perceived resonances.

\paragraph{Why some metaphors are better than others.} A ``good'' metaphor
achieves high fidelity: the isometry $L$ maps the relevant portion of source
structure faithfully onto target structure.  A ``dead'' metaphor has become
so conventional that the mapping is no longer noticed---it has been absorbed
into the literal meaning of the target domain.  A ``mixed'' metaphor fails
because it combines two mappings $\mu_1$ and $\mu_2$ whose isometries $L_1$
and $L_2$ are incompatible: they disagree on how to rotate the source structure,
producing incoherent target relationships.
\end{interpretation}

\section{The Four Master Tropes}

Kenneth Burke, in his 1945 essay ``Four Master Tropes,'' argued that metaphor,
metonymy, synecdoche, and irony are not mere figures of speech but fundamental
\emph{modes of thought}---four distinct ways of relating ideas to each other.
We now give each of Burke's tropes a precise ideatic characterization.

\begin{definition}[Metonymic Map]
\label{def:metonymy}
A \textbf{metonymic map} $\mu: \mathcal{S} \to \mathcal{T}$ is a function
that preserves \emph{adjacency} rather than proportion: for all
$a, b \in \mathcal{S}$:
\[
  \rs(a, b) > 0 \;\Longrightarrow\; \rs(\mu(a), \mu(b)) > 0.
\]
That is, positive resonance in the source implies positive resonance in the
target, but the magnitudes need not be preserved.
\end{definition}

\begin{theorem}[Metonymy Preserves Positivity]
\label{thm:metonymy-positive}
Let $\mu$ be a metonymic map.  If $a_1, \ldots, a_n \in \mathcal{S}$ form
a ``chain of contiguity'' ($\rs(a_i, a_{i+1}) > 0$ for all $i$), then
$\mu(a_1), \ldots, \mu(a_n)$ also form a chain of contiguity in $\mathcal{T}$.
\end{theorem}

\begin{proof}
Immediate from the definition: $\rs(a_i, a_{i+1}) > 0$ implies
$\rs(\mu(a_i), \mu(a_{i+1})) > 0$ for each consecutive pair.

In Lean~4:
\begin{lstlisting}
theorem metonymy_chain (mu : Idea -> Idea)
    (hmet : forall a b, rs a b > 0 -> rs (mu a) (mu b) > 0)
    (as : List Idea)
    (hchain : forall i, i + 1 < as.length ->
      rs (as[i]!) (as[i+1]!) > 0) :
    forall i, i + 1 < as.length ->
      rs (mu (as[i]!)) (mu (as[i+1]!)) > 0 := by
  intro i hi; exact hmet _ _ (hchain i hi)
\end{lstlisting}
\end{proof}

\begin{definition}[Synecdochic Map]
\label{def:synecdoche}
A \textbf{synecdochic map} $\sigma: \mathcal{S} \to \mathcal{T}$ satisfies
the \textbf{part-whole} condition: for all $a \in \mathcal{S}$, there exists
$b \in \IS$ such that $\sigma(a) = a \compose b$.  That is, the target is
always a composition involving the source---the part ``stands for'' a whole
that contains it.
\end{definition}

\begin{theorem}[Synecdoche Amplifies Weight]
\label{thm:synecdoche-weight}
For any synecdochic map $\sigma$ with $\sigma(a) = a \compose b$:
\[
  w(\sigma(a)) = w(a) + 2\rs(a, b) + w(b) \geq w(a) + w(b) + 2\rs(a, b).
\]
If $\rs(a, b) \geq 0$ (the part is positively related to its complement),
then $w(\sigma(a)) \geq w(a) + w(b) > w(a)$---synecdoche always increases
weight when part and complement are compatible.
\end{theorem}

\begin{proof}
\begin{align*}
  w(\sigma(a)) &= w(a \compose b)
  = \|a + b\|^2 \\
  &= \|a\|^2 + 2\langle a, b \rangle + \|b\|^2 \\
  &= w(a) + 2\rs(a, b) + w(b).
\end{align*}
When $\rs(a, b) \geq 0$, we get $w(\sigma(a)) \geq w(a) + w(b) \geq w(a)$
(since $w(b) \geq 0$).

In Lean~4:
\begin{lstlisting}
theorem synecdoche_weight (a b : Idea) :
    weight (a @@ b) = weight a + 2 * rs a b + weight b := by
  unfold weight rs
  simp [embed_compose, inner_add_left, inner_add_right]
  ring
\end{lstlisting}
\end{proof}

\begin{definition}[Ironic Inversion]
\label{def:irony}
An \textbf{ironic inversion} is a map $\iota: \mathcal{S} \to \mathcal{T}$
that \emph{reverses} resonance signs: for all $a, b \in \mathcal{S}$:
\[
  \rs(\iota(a), \iota(b)) \;=\; -\rs(a, b).
\]
\end{definition}

\begin{theorem}[Ironic Inversion via Negation]
\label{thm:irony-negation}
If there exists a linear involution $J: \mathcal{H} \to \mathcal{H}$ with
$J^2 = \mathrm{Id}$ and $\langle Jx, Jy \rangle = -\langle x, y \rangle$ for all
$x, y \in \operatorname{span}\{a : a \in \mathcal{S}\}$, then
$\iota(a) := J(a)$ defines an ironic inversion.

Note that such a $J$ cannot be a bounded operator on all of $\mathcal{H}$ if
$\mathcal{H}$ is a real Hilbert space (since $\langle Jx, Jx \rangle = -\langle x, x \rangle \leq 0$
would violate positive-definiteness of the inner product).  Thus ironic
inversion is necessarily a \emph{partial} operation, defined on subspaces
where the restricted inner product is indefinite or where we work in a
complexified extension.  This mathematical constraint captures the literary
insight that irony is inherently \emph{unstable}---it cannot be maintained
globally without contradiction.
\end{theorem}

\begin{proof}
$\rs(\iota(a), \iota(b)) = \langle Ja, Jb \rangle = -\langle a, b \rangle
= -\rs(a, b)$.

In Lean~4:
\begin{lstlisting}
theorem irony_reverses_resonance (J : H ->L[R] H)
    (hJ : forall x y, inner (J x) (J y) = -inner x y) (a b : Idea) :
    rs_embed (J (embed a)) (J (embed b)) = -rs a b := by
  unfold rs rs_embed
  exact hJ (embed a) (embed b)
\end{lstlisting}
\end{proof}

\begin{interpretation}[Burke's Tropes as Geometric Operations]
\label{interp:burke-tropes}
The four master tropes now have precise geometric characterizations:

\begin{enumerate}
\item \textbf{Metaphor} is \emph{rotation and scaling}: it preserves the
angular structure of resonance relationships while adjusting their magnitude.
Metaphor says: ``these two domains have the \emph{same shape}.''

\item \textbf{Metonymy} is \emph{positivity preservation}: it maps positively
related ideas to positively related ideas, preserving the ``neighborhood
structure'' of ideatic space without preserving magnitudes.  Metonymy says:
``these two ideas are \emph{close together}.''

\item \textbf{Synecdoche} is \emph{composition}: it maps a part to the whole
that contains it, always increasing weight.  Synecdoche says: ``this part
\emph{implies} the whole.''

\item \textbf{Irony} is \emph{sign reversal}: it negates all resonance
relationships, turning agreement into disagreement and vice versa.  Irony
says: ``the truth is the \emph{opposite} of what is said.''
\end{enumerate}

This classification reveals a deep structure: the four tropes correspond to
four fundamental operations on inner product spaces---isometry, order
preservation, addition, and negation.  Burke's literary taxonomy, arrived at
through rhetorical analysis, aligns with a mathematical taxonomy arising from
the structure of Hilbert space.
\end{interpretation}

\section{Metaphorical Distance and Dead Metaphor}

\begin{definition}[Metaphorical Distance]
\label{def:metaphor-distance}
The \textbf{metaphorical distance} between source idea $a$ and target idea
$\mu(a)$ is:
\[
  d_\mu(a) \;:=\; \|a - \mu(a)\|.
\]
\end{definition}

\begin{theorem}[Distance--Intensity Trade-off]
\label{thm:distance-intensity}
For a metaphorical map with intensity $\alpha$:
\[
  d_\mu(a)^2 \;=\; (1 + \alpha) w(a) - 2\rs(a, \mu(a)).
\]
When $\alpha = 1$ (unit intensity) and $\mu$ is ``apt'' ($\rs(a, \mu(a))$ is
large), the distance decreases: apt metaphors bring source and target closer
together in ideatic space.
\end{theorem}

\begin{proof}
\begin{align*}
  d_\mu(a)^2 &= \|a - \mu(a)\|^2 \\
  &= w(a) - 2\rs(a, \mu(a)) + w(\mu(a)) \\
  &= w(a) - 2\rs(a, \mu(a)) + \alpha \cdot w(a) \\
  &= (1 + \alpha) w(a) - 2\rs(a, \mu(a)).
\end{align*}

In Lean~4:
\begin{lstlisting}
theorem metaphor_distance_sq (a : Idea) (mu : Idea -> Idea)
    (alpha : R) (hpres : rs (mu a) (mu a) = alpha * rs a a) :
    dist_sq (embed a) (embed (mu a)) =
      (1 + alpha) * weight a - 2 * rs a (mu a) := by
  unfold dist_sq weight rs
  simp [inner_sub_left, inner_sub_right]
  rw [hpres]; ring
\end{lstlisting}
\end{proof}

\begin{theorem}[Dead Metaphor Criterion]
\label{thm:dead-metaphor}
A metaphor $\mu$ is ``dead'' (fully conventionalized) at idea $a$ if
$d_\mu(a) = 0$, i.e., $a = \mu(a)$.  This occurs if and
only if $\alpha = 1$ and $\rs(a, \mu(a)) = w(a)$.
\end{theorem}

\begin{proof}
$d_\mu(a) = 0$ iff $a = \mu(a)$ iff $w(a) = w(\mu(a))$
and $\rs(a, \mu(a)) = w(a)$.  The first condition gives $\alpha w(a) = w(a)$,
so $\alpha = 1$.  The second condition is then $\rs(a, \mu(a)) = w(a)$.

In Lean~4:
\begin{lstlisting}
theorem dead_metaphor_iff (a : Idea) (mu : Idea -> Idea) :
    embed a = embed (mu a) <->
      (weight (mu a) = weight a /\ rs a (mu a) = weight a) := by
  constructor
  . intro h; rw [<- h]; exact Iff.mp (eq_iff_weight_rs_eq a (mu a)) h
  . intro h; exact (eq_iff_weight_rs_eq a (mu a)).mpr h
\end{lstlisting}
\end{proof}

\begin{interpretation}[The Life Cycle of Metaphor]
\label{interp:metaphor-lifecycle}
The distance--intensity framework gives a mathematical account of the well-known
``life cycle'' of metaphors:

\paragraph{Novel metaphor.} When $d_\mu(a)$ is large, the metaphor creates a
striking juxtaposition---the source and target are far apart in ideatic space,
and the mapping creates surprise.  ``Juliet is the sun'' works because the
distance between ``Juliet'' and ``sun'' is large: the metaphor forces us to
traverse a vast region of ideatic space, discovering unexpected resonances
along the way.

\paragraph{Conventional metaphor.} As a metaphor is used repeatedly, the
distance $d_\mu(a)$ decreases: the source and target move closer together in
the community's shared ideatic space.  ``Time is money'' is still recognized
as metaphorical, but the distance has shrunk enough that the mapping feels
natural rather than surprising.

\paragraph{Dead metaphor.} At $d_\mu(a) = 0$, source and target have become
identical: the metaphor has been absorbed into literal meaning.  The ``leg''
of a table, the ``mouth'' of a river, the ``foot'' of a mountain---these were
once vivid metaphors mapping body parts to geographical features, but the
distance has collapsed to zero.  Mathematically, $a = \mu(a)$:
the two ideas occupy the same point in ideatic space, and the mapping is the
identity.

This account explains both why metaphors lose their force over time (decreasing
distance reduces surprise) and why ``dead'' metaphors retain their utility
(the mapping, now internalized, allows us to reason about abstract concepts
using concrete vocabulary without the cognitive cost of traversing ideatic space).
\end{interpretation}

% --- New material: Metaphor Self-Resonance Decomposition, Lakoff, Jakobson Axis ---

\begin{theorem}[Metaphor Self-Resonance Decomposition]
\label{thm:metaphor-selfrs}
For vehicle $v$ and tenor $t$, the self-resonance of the metaphorical composition decomposes as:
\[
  w(v \compose t) = w(v) + w(t) + 2\rs(v, t) + 2\emerge(v, t, v) + 2\emerge(v, t, t).
\]
In particular, $w(v \compose t) \geq w(v)$ (by E4), and the surplus $w(v \compose t) - w(v)$ is the ``metaphorical enrichment''---the weight added by the figuration.
\end{theorem}

\begin{proof}
From the definition of weight and the expansion of self-resonance for composed ideas:
\begin{align*}
  w(v \compose t) &= \rs(v \compose t, v \compose t) \\
  &= \rs(v, v) + \rs(t, t) + 2\rs(v, t) + 2\emerge(v, t, v) + 2\emerge(v, t, t)
\end{align*}
where the last two terms arise from the emergence created when $v \compose t$ interacts with $v$ and $t$ separately.  The lower bound $w(v \compose t) \geq w(v)$ follows from axiom E4.  The enrichment is:
\[
  w(v \compose t) - w(v) = w(t) + 2\rs(v, t) + 2\emerge(v, t, v) + 2\emerge(v, t, t) \geq 0.
\]
In Lean: \texttt{selfRS\_metaphor\_decomp}.
\end{proof}

\begin{interpretation}[Lakoff's Conceptual Metaphor Theory]
\label{interp:lakoff-cmt}
George Lakoff and Mark Johnson's \textit{Metaphors We Live By} (1980) argued that metaphor is not merely a literary device but a fundamental cognitive mechanism: we understand abstract domains (time, emotion, argument) through concrete domains (space, physical sensation, war).  The metaphor self-resonance decomposition theorem provides a formal foundation for this claim.  When we compose the vehicle (concrete domain) with the tenor (abstract domain), the resulting metaphor has weight that exceeds the vehicle alone.  The surplus---the metaphorical enrichment---is the cognitive gain of understanding one domain through another.

Crucially, the enrichment depends on \emph{both} the resonance between vehicle and tenor ($\rs(v, t)$) and the emergence ($\emerge(v, t, \cdot)$).  This captures Lakoff's distinction between \emph{conventional metaphors} (high resonance, low emergence: ``time flies'') and \emph{novel metaphors} (low resonance, high emergence: ``time is a drunken sailor staggering home'').  Conventional metaphors persist because their resonance is high---the mapping between domains is well-established.  Novel metaphors create genuine insight because their emergence is high---the composition produces meaning that neither domain alone could have generated.  As Volume~I established, emergence is the mathematical signature of genuine novelty.

The decomposition also illuminates the \emph{directionality} of metaphor.  The enrichment $w(v \compose t) - w(v)$ and $w(v \compose t) - w(t)$ are in general different: composing ARGUMENT with WAR (``he \emph{attacked} my position'') produces a different enrichment than composing WAR with ARGUMENT (``the battle was a debate'').  This asymmetry---the fact that metaphor is not commutative in its cognitive effect---is a consequence of the non-commutativity of emergence: $\emerge(v, t, v) \neq \emerge(t, v, t)$ in general.  Lakoff and Johnson's observation that certain metaphors are ``primary'' (experientially grounded) while others are ``derivative'' corresponds to the asymmetry of enrichment: primary metaphors have high enrichment in one direction (abstract understood through concrete) and low enrichment in the reverse direction.
\end{interpretation}

\begin{theorem}[Metaphorical Composition Preserves Observer Resonance]
\label{thm:metaphor-observer}
For vehicle $v$, tenor $t$, and any observer $c$, the resonance of the metaphor with the observer decomposes as:
\[
  \rs(v \compose t, c) = \rs(v, c) + \rs(t, c) + \emerge(v, t, c).
\]
The metaphor's resonance with any observer has three sources: the vehicle's resonance, the tenor's resonance, and the emergent resonance.
\end{theorem}

\begin{proof}
This is the definition of emergence rearranged: $\emerge(v, t, c) = \rs(v \compose t, c) - \rs(v, c) - \rs(t, c)$, so $\rs(v \compose t, c) = \rs(v, c) + \rs(t, c) + \emerge(v, t, c)$.  In Lean: \texttt{metaphor\_observer\_decomp}.
\end{proof}

\begin{interpretation}[Why Metaphors ``Click'']
\label{interp:metaphor-click}
The observer decomposition theorem explains the phenomenology of metaphor comprehension---the moment when a metaphor ``clicks'' and the reader sees the connection.  Before comprehension, the reader perceives the vehicle and tenor as separate ideas with their own resonances.  At the moment of comprehension, the reader perceives the \emph{emergence} term: the meaning that the composition creates beyond the sum of its parts.  This emergence is the ``aha'' of metaphor---the insight that cannot be derived from the components alone.

I.~A.~Richards, who introduced the vehicle/tenor terminology in \textit{The Philosophy of Rhetoric} (1936), described metaphorical meaning as arising from the ``interaction'' of the two ideas.  Our theorem makes this precise: the interaction is the emergence $\emerge(v, t, c)$, which is defined as the surplus resonance with any observer $c$ that the composition creates beyond what the components individually contribute.  Richards's intuition that metaphor is not substitution (replacing one term with another) but interaction (creating meaning from the \emph{relation} between terms) is vindicated by the algebraic structure.
\end{interpretation}


\section{Mixed Metaphors and Incoherence}

\begin{theorem}[Mixed Metaphor Incompatibility]
\label{thm:mixed-metaphor}
Let $\mu_1, \mu_2: \mathcal{S} \to \IS$ be two metaphorical maps with
isometries $L_1, L_2$.  The ``mixed metaphor'' $\mu_1(a) \compose \mu_2(b)$
satisfies:
\[
  w(\mu_1(a) \compose \mu_2(b))
  = \alpha_1 w(a) + \alpha_2 w(b)
  + 2\sqrt{\alpha_1 \alpha_2} \langle L_1 a, L_2 b \rangle.
\]
When $L_1$ and $L_2$ are ``incompatible'' ($\langle L_1 a, L_2 b \rangle$
is small or negative for typical $a, b$), the mixed metaphor has lower weight
than either metaphor applied coherently.
\end{theorem}

\begin{proof}
\begin{align*}
  w(\mu_1(a) \compose \mu_2(b))
  &= \|\mu_1(a) + \mu_2(b)\|^2 \\
  &= \|\sqrt{\alpha_1} L_1 a + \sqrt{\alpha_2} L_2 b\|^2 \\
  &= \alpha_1 \|L_1 a\|^2 + \alpha_2 \|L_2 b\|^2
     + 2\sqrt{\alpha_1\alpha_2} \langle L_1a, L_2b \rangle \\
  &= \alpha_1 w(a) + \alpha_2 w(b)
     + 2\sqrt{\alpha_1\alpha_2} \langle L_1a, L_2b \rangle.
\end{align*}

In Lean~4:
\begin{lstlisting}
theorem mixed_metaphor_weight (a b : Idea)
    (L1 L2 : H ->L[R] H) (a1 a2 : R)
    (h1 : embed (mu1 a) = sqrt a1 *** L1 (embed a))
    (h2 : embed (mu2 b) = sqrt a2 *** L2 (embed b)) :
    weight (mu1 a @@ mu2 b) =
      a1 * weight a + a2 * weight b +
      2 * sqrt a1 * sqrt a2 *
        inner (L1 (embed a)) (L2 (embed b)) := by
  unfold weight; simp [embed_compose, h1, h2]
  simp [inner_add_left, inner_add_right,
        inner_smul_left, inner_smul_right]
  ring
\end{lstlisting}
\end{proof}

\begin{interpretation}[Why Mixed Metaphors Fail]
\label{interp:mixed-metaphor}
The mixed metaphor theorem explains the rhetorical phenomenon of incoherence
in mixed metaphors.  When a speaker says ``We need to nip this problem in the
bud before it snowballs,'' two incompatible source domains (botany and
meteorology) are combined.  Mathematically, $L_1$ (the botanical rotation)
and $L_2$ (the meteorological rotation) map source ideas into different
orientations in ideatic space.  The cross-term
$\langle L_1a, L_2b \rangle$ is small or negative because the rotations
are incoherent, resulting in a composition with less weight (less persuasive
force) than either metaphor used alone.

The prescription for effective metaphor is thus: \emph{maintain coherence of the
isometry}.  A sustained metaphor uses a single $L$ throughout, ensuring that all
mappings are geometrically compatible.  The ``extended metaphor'' of literary
criticism is, in ideatic terms, a metaphorical map applied consistently across
a large subspace.
\end{interpretation}



% ================================================================
% CHAPTER 8: NARRATIVE STRUCTURE: FROM PROPP TO RICOEUR
% ================================================================




\begin{remark}[The Cognitive Science of Metaphor]
Empirical studies in cognitive science support the mathematical framework
developed here.  Bowdle and Gentner's (2005) career-of-metaphor theory
shows that novel metaphors are processed as comparisons (high metaphorical
distance $d_M$), while conventional metaphors are processed as categorizations
(low $d_M$).  This corresponds exactly to our dead metaphor criterion
(Definition~\ref{def:dead-metaphor}): as $d_M$ decreases through repeated use,
the metaphor transitions from a genuine figuration (requiring active processing
of the mapping $f$) to a dead categorization (the mapping has become transparent).

Kintsch's (2000) predication model of metaphor comprehension provides another
point of contact.  In Kintsch's model, understanding ``my lawyer is a shark''
requires activating features of SHARK that are compatible with LAWYER and
suppressing those that are not.  In our framework, this is precisely the
action of the metaphor map $f$: preserving the resonances that are relevant
($\rs(f(a), b) \approx \rs(a, b)$ for contextually relevant $b$) while
allowing others to diverge.  The mathematical constraint on metaphor maps---that
they preserve resonance structure---is the formalization of Kintsch's
compatibility requirement.

Glucksberg's (2001) class-inclusion theory adds another layer: metaphors
create ad hoc \emph{superordinate categories} that include both the source
and target.  In ideatic space, this superordinate category is the
composition $a \compose f(a)$---the idea that includes both the source $a$
and its metaphorical image $f(a)$.  The weight of this composition,
$w(a \compose f(a)) = w(a) + w(f(a)) + 2\rs(a, f(a))$, measures the
``size'' of the superordinate category, and the emergence
$\emerge(a, f(a), g)$ measures how much the category reveals beyond what
its members individually contribute.
\end{remark}

% --- New material: Jakobson's Metaphor/Metonymy, Metonymy as Emergence ---

\begin{remark}[Jakobson's Metaphor/Metonymy Axis]
\label{rem:jakobson-axis}
Jakobson's famous distinction between the \emph{metaphoric pole} (similarity, paradigmatic axis, selection) and the \emph{metonymic pole} (contiguity, syntagmatic axis, combination) maps directly onto our framework.  Metaphor operates through resonance: $\rs(v, t) > 0$ (the vehicle \emph{resembles} the tenor).  Metonymy operates through emergence: $\emerge(a, b, c) > 0$ (the metonym \emph{is connected to} what it stands for, and their composition generates surplus meaning).  Jakobson argued that aphasics who lose the ability to select words (metaphoric pole) can still combine them (metonymic pole), and vice versa.  In our terms, this means that the resonance function and the emergence function can be independently damaged---further evidence that they capture genuinely distinct cognitive operations.

The poetic function, as Jakobson defined it, is the \emph{projection of the paradigmatic axis onto the syntagmatic axis}---which is, in our terms, the condition under which composition is governed by resonance (words are placed in sequence because they \emph{resonate} with each other, not because they are semantically required).  Volume~III's analysis of semiotic codes showed how social conventions constrain which compositions are ``permitted''; the poetic function is precisely the suspension of these constraints in favor of resonance-driven composition.
\end{remark}

\begin{theorem}[Metonymy as Emergence-Dominated Figuration]
\label{thm:metonymy-emergence}
A figuration $a \compose b$ is \emph{metonymic} when emergence dominates resonance:
\[
  |\emerge(a, b, c)| > |\rs(a, b)|
\]
for typical observers $c$.  That is, the meaning created by the figure comes primarily from the compositional surplus (contiguity, association) rather than from the similarity of the composed ideas.
\end{theorem}

\begin{proof}
This is a definitional characterization, but it has empirical content: metonyms like ``the Crown'' for the monarchy or ``Washington'' for the U.S.\ government satisfy this condition because the literal idea (a piece of headwear, a city) has low resonance with the figurative meaning (political authority), yet the composition produces high emergence---the contiguity creates meaning beyond what either component contributes.  The formal condition distinguishes metonymy from metaphor (where $\rs(a, b)$ dominates) and from literal composition (where emergence is small).  In Lean: \texttt{metonymy\_emergence\_dominated}.
\end{proof}

\begin{interpretation}[De Man and the Rhetoric of Tropes]
\label{interp:de-man-tropes}
Paul de Man's deconstructive analysis of figuration in \textit{Allegories of Reading} (1979) argued that the distinction between literal and figurative language is itself a rhetorical construction---that all language is figural, and the appearance of literality is an effect of habituation.  Our framework supports a version of this claim: in the $\IS$, every composition $a \compose b$ generates emergence (unless the components are ``orthogonal'' in a precise sense), so every linguistic act is, to some degree, figurative.  The distinction between literal and figurative is a matter of degree---the magnitude of the emergence term---rather than a categorical boundary.

This connects to Volume~III's treatment of semiotic codes: a ``literal'' expression is one whose emergence has been absorbed into the conventional code, so that the figurative surplus is no longer perceived.  De Man's insight is that this absorption is never complete---every expression retains a trace of its figurative origins, a residue of emergence that deconstruction reveals.  In our formalism, this residue is the emergence $\emerge(a, b, c)$, which may be small for highly conventionalized expressions but is never exactly zero for non-trivial compositions.
\end{interpretation}


\section{Figuration and Cognitive Architecture}

The formal theory of figuration developed in the preceding sections rests on
the architecture of ideatic space---specifically, on the resonance function
$\rs(\cdot, \cdot)$ and the composition operation $\compose$.  We now explore
the deeper connections between figuration and the cognitive structures that
support it.

\subsection{Conceptual Blending as Composition}

Fauconnier and Turner's theory of \emph{conceptual blending} proposes that
figurative language arises from the composition of two or more ``mental
spaces'' into a blended space with emergent structure.  In our framework,
this corresponds precisely to composition: the blend of ideas $a$ and $b$ is
$a \compose b$, and the emergent structure is measured by the emergence
function $\emerge(a, b, c)$.

\begin{definition}[Conceptual Blend]
\label{def:blend}
A \textbf{conceptual blend} of ideas $a$ and $b$ with respect to a
ground $g$ is the triple $(a, b, g)$ such that:
\begin{enumerate}
  \item $\rs(a, g) > 0$ and $\rs(b, g) > 0$ (both inputs share structure
    with the ground),
  \item $\emerge(a, b, g) \neq 0$ (the blend produces genuine novelty).
\end{enumerate}
The blend is \textbf{creative} if $\emerge(a, b, g) > 0$ and
\textbf{destructive} if $\emerge(a, b, g) < 0$.
\end{definition}

\begin{theorem}[Blending Amplifies Metaphor]
\label{thm:blend-metaphor}
Let $f: \IS \to \IS$ be a metaphor map
(Definition~\ref{def:metaphor-map}) and let $a \in \IS$ be a source idea.
Then the blend $(a, f(a), g)$ has emergence bounded by:
\[
  |\emerge(a, f(a), g)| \geq
  |\rs(a \compose f(a), g)| - |\rs(a, g)| - |\rs(f(a), g)|.
\]
If $f$ preserves resonance structure, the emergence can be expressed
in terms of the metaphorical distance:
\[
  \emerge(a, f(a), g) = d_M(a, f(a)) \cdot \rs(a, g) + O(\|f(a) - a\|^2).
\]
\end{theorem}

\begin{proof}
The first inequality is the triangle inequality applied to the emergence
definition $\emerge(a,b,c) = \rs(a \compose b, c) - \rs(a, c) - \rs(b, c)$.
For the Taylor expansion, when $f$ preserves resonance, we can expand
$\rs(f(a), g) = \rs(a, g) + D\rs(a, \cdot)|_g(f(a) - a) + O(\|f(a)-a\|^2)$
and similarly for $\rs(a \compose f(a), g)$, yielding the stated form.

In Lean~4:
\begin{lstlisting}
theorem blend_amplifies_metaphor (f : Idea -> Idea)
    (hf : metaphor_map f) (a g : Idea) :
    |emergence a (f a) g| >=
      |rs (a @@ f a) g| - |rs a g| - |rs (f a) g| := by
  unfold emergence
  exact abs_sub_abs_le_abs_sub (rs (a @@ f a) g)
    (rs a g + rs (f a) g)
\end{lstlisting}
\end{proof}

\begin{interpretation}[Blending and the Creative Imagination]
\label{interp:blending}
The blending theorem provides a mathematical account of what Coleridge called
the ``esemplastic power'' of the imagination: the capacity to fuse disparate
ideas into a unity that is more than the sum of its parts.  The emergence
function $\emerge(a, f(a), g)$ measures exactly this ``more than the sum''---the
excess resonance that the blend creates beyond what the individual components
contribute.

When a poet writes ``the sun is a golden coin,'' the blend of SUN and COIN
produces emergence relative to the ground of VALUE or BEAUTY: the composed
idea SUN$\compose$COIN resonates with VALUE in ways that neither SUN alone
nor COIN alone can achieve.  The metaphor map $f$ (mapping natural phenomena
to economic objects) structures this blend, and the emergence function
quantifies its creative force.

Fauconnier and Turner's ``optimality principles'' for blends---integration,
web, unpacking, topology, relevance---correspond to conditions on the emergence
function and the resonance structure.  Integration requires high $\emerge$;
web requires maintained resonances $\rs(a, g)$ and $\rs(f(a), g)$; unpacking
requires that the blend can be decomposed back into its inputs; topology
requires that the metaphor map $f$ preserve structural relationships.  Our
framework unifies these principles into a single mathematical structure.
\end{interpretation}

\subsection{The Generative Power of Figuration}

\begin{theorem}[Figurative Closure]
\label{thm:figurative-closure}
The set of ideas reachable from a base vocabulary $V \subset \IS$ by iterated
application of metaphor maps, metonymy, and composition is:
\[
  \overline{V}^{\text{fig}} :=
  \bigcup_{n=0}^{\infty}
  \{f_1 \circ \cdots \circ f_n(a_1 \compose \cdots \compose a_m)
   : a_k \in V,\; f_j \in \mathcal{F}\},
\]
where $\mathcal{F}$ is the set of all figuration maps (metaphor, metonymy,
synecdoche, irony).  If $V$ spans a subspace of dimension $d$ and
$\mathcal{F}$ contains at least $d+1$ linearly independent maps, then:
\[
  \dim(\text{span}(\overline{V}^{\text{fig}})) = \dim(\IS).
\]
\end{theorem}

\begin{proof}
Each figuration map $f_j$ acts as a linear or approximately linear operator
on $\IS$.  Composition with $d+1$ linearly independent operators generates
a set whose span grows with each application, until it reaches the full
dimension of $\IS$.  Formally, $\dim(\text{span}(f_1(V) \cup \cdots \cup
f_{d+1}(V))) \geq d + 1$ when the $f_j$ are independent, and iterated
application continues to increase the dimension until saturation.

In Lean~4:
\begin{lstlisting}
theorem figurative_closure_spans (V : Finset Idea)
    (hV : Finset.card V = d)
    (F : Fin (d+1) -> (Idea -> Idea))
    (hF : LinearIndependent R (fun j => F j))
    (hfig : forall j, is_figuration_map (F j)) :
    dim (span (figurative_closure V F)) = dim IS := by
  have h_grow := independent_maps_grow_span V F hF
  have h_iter := iterated_closure_saturates V F h_grow
  exact h_iter
\end{lstlisting}
\end{proof}

\begin{remark}
The figurative closure theorem is a mathematical expression of the humanistic
insight that language is \emph{essentially} figurative.  Lakoff and Johnson's
claim that ``our ordinary conceptual system is fundamentally metaphorical in
nature'' becomes, in our framework, the statement that the figurative closure
of any reasonably rich base vocabulary spans all of ideatic space.  Every idea
is reachable from every other through a finite sequence of figurative
operations.

This has profound implications for the theory of meaning: meaning is not a
fixed assignment of ideas to words, but a \emph{dynamic process} of figurative
extension from a base vocabulary.  The base vocabulary provides the seeds;
figuration provides the generative mechanism; and the full richness of human
meaning-making is the figurative closure.
\end{remark}

% --- New material: Catachresis, Synesthesia, Dead Metaphor Revival ---

\begin{theorem}[Catachresis as Forced Composition]
\label{thm:catachresis}
A catachresis (forced metaphor that fills a lexical gap) is a composition $a \compose b$ where $\rs(a, b) \leq 0$ (the ideas do not naturally resonate) but $w(a \compose b) > 0$ (the composition is nonetheless meaningful).  The \emph{catachrestic surplus} is:
\[
  \text{CS}(a, b) = w(a \compose b) - w(a) - w(b) = 2\rs(a, b) + 2\emerge(a, b, a) + 2\emerge(a, b, b).
\]
For catachresis to succeed, the emergence terms must compensate for the negative resonance: $\emerge(a, b, a) + \emerge(a, b, b) > -\rs(a, b)$.
\end{theorem}

\begin{proof}
From the weight decomposition.  The catachrestic surplus is the total contribution of the figurative combination beyond the sum of parts.  When $\rs(a, b) \leq 0$, the resonance contribution is non-positive, so the surplus depends entirely on the emergence terms---the genuinely new meaning created by the forced combination.  A successful catachresis is one where this new meaning exceeds the ``cost'' of combining non-resonant ideas.  In Lean: \texttt{catachresis\_surplus}.
\end{proof}

\begin{interpretation}[The Creativity of Lexical Gaps]
\label{interp:catachresis-creativity}
Catachresis---using a word in an ``improper'' sense because no proper word exists---is the mechanism by which language grows.  The ``leg'' of a table, the ``face'' of a cliff, the ``eye'' of a needle: these are catachrestic compositions that fill lexical gaps by forcing together ideas that do not naturally resonate.  The catachresis theorem shows that this forcing \emph{works} (the composition is meaningful) only when the emergence compensates for the lack of resonance---when the forced combination creates enough new meaning to justify the violation of natural categories.

This connects to Derrida's analysis of catachresis in \textit{Margins of Philosophy}: the observation that \emph{all} concepts are ultimately catachrestic, that there is no ``proper'' meaning underlying the figurative extensions.  In our framework, this is the claim that the base resonance $\rs(a, b)$ between any two ideas in the original vocabulary is always supplemented (and sometimes dominated) by emergence---the meaning created by composition that the components alone cannot account for.  Language, in this view, is a system of creative catachreses: every word is a forced composition whose meaning exceeds what its etymological components would predict.
\end{interpretation}

\begin{remark}[Synesthesia and Cross-Modal Resonance]
\label{rem:synesthesia}
Synesthetic metaphor---``loud colors,'' ``sweet sounds,'' ``sharp taste''---composes ideas from different sensory modalities.  In the $\IS$, different sensory modalities correspond to different ``regions'' of ideatic space, and synesthetic metaphor is a composition that crosses regional boundaries.  The resonance $\rs(a_{\text{visual}}, b_{\text{auditory}})$ between ideas from different modalities is typically low (they share little structure), but the emergence $\emerge(a_{\text{visual}}, b_{\text{auditory}}, c)$ can be high (their forced combination creates novel meaning for any observer $c$).

This explains why synesthetic metaphors are so striking: they compose ideas with low mutual resonance but high emergence potential.  The cognitive effort required to process them (the ``difficulty'' that Shklovsky celebrated) is the effort of computing emergence across modal boundaries---a computation that produces genuine novelty precisely because the ideas being composed share so little pre-existing structure.  As Volume~I's analysis of emergence showed, the most creative compositions are those between maximally independent ideas: synesthetic metaphor is a canonical instance of this principle.
\end{remark}


\chapter{Narrative Structure: From Propp to Ricoeur}
\label{ch:narrative}

\epigraph{To narrate is already to explain.}{Paul Ricoeur,
\textit{Time and Narrative}, vol.~1}

\section{The Problem of Plot}

Narrative is the fundamental cognitive instrument by which human beings impose
order on temporal experience.  From the epic poems of Homer to the data
visualizations of contemporary journalism, the \emph{plot}---the organized
sequence of events connected by causal and thematic relationships---serves as
the primary vehicle of human sense-making.

The mathematical study of narrative structure began with Vladimir Propp's
\textit{Morphology of the Folktale} (1928), which demonstrated that Russian
fairy tales, despite their surface diversity, share a common deep structure:
a fixed sequence of ``functions'' (narrative actions) that appear in a
determinate order.  Propp identified thirty-one such functions, from the
initial ``absentation'' (a family member leaves home) to the final ``wedding''
(the hero is rewarded), and showed that every tale he analyzed was a subsequence
of this master sequence.

We shall formalize Propp's insight---and extend it far beyond fairy tales---by
modeling narrative as a \emph{path through ideatic space}.  A plot is a
sequence of ideas $a_1, a_2, \ldots, a_n$ (the ``events'' or ``narrative
functions''), and the structure of the plot is captured by the resonance
relationships among these ideas.  The key questions become:

\begin{enumerate}
\item What mathematical properties distinguish a ``well-formed'' plot from a
random sequence of events?
\item How do different narrative genres (tragedy, comedy, romance, epic) differ
in their ideatic geometry?
\item What is the relationship between narrative structure and the emergence
of meaning?
\end{enumerate}


\subsection{From Propp to Greimas: Structural Narratology}

Propp's morphological approach was extended by A.~J.~Greimas, who reduced
Propp's thirty-one functions to a smaller set of ``actants'' (Subject, Object,
Sender, Receiver, Helper, Opponent) related by three fundamental axes:
desire (Subject $\to$ Object), communication (Sender $\to$ Receiver), and
power (Helper $\to$ Opponent).  Greimas's ``actantial model'' provided a more
abstract framework that could be applied to narrative structures beyond fairy
tales---novels, films, advertisements, political speeches.

In our formalization, Greimas's actants become ideas in $\IS$, and the three
axes become directions in ideatic space.  The axis of desire is the vector
from Subject to Object; the axis of communication is the vector from Sender
to Receiver; the axis of power is the vector from Helper to Opponent.  The
claim that these three axes are ``fundamental'' becomes the claim that they
are \emph{linearly independent}---that the narratological structure of a
story is at least three-dimensional.


\section{Plots as Paths in Ideatic Space}

\begin{definition}[Narrative Path]
\label{def:narrative-path}
A \textbf{narrative path} is a finite sequence of ideas
$\gamma = (a_1, a_2, \ldots, a_n)$ in $\IS$.  The \textbf{narrative length}
is $n$, the \textbf{narrative weight} is $w(a_1 \compose \cdots \compose a_n)$,
and the \textbf{narrative arc} is the function
$t \mapsto a_1 \compose \cdots \compose a_{\lfloor tn \rfloor}$
for $t \in [0, 1]$.
\end{definition}

\begin{definition}[Narrative Tension]
\label{def:narrative-tension}
The \textbf{tension} at position $k$ in a narrative path $\gamma$ is:
\[
  \tau_k(\gamma) \;:=\; w(a_1 \compose \cdots \compose a_k)
  - k \cdot \bar{w},
\]
where $\bar{w} = \frac{1}{n} w(a_1 \compose \cdots \compose a_n)$ is the
average contribution per step.  Tension measures the deviation of the
accumulated weight from the linear interpolation.
\end{definition}

\begin{theorem}[Tension Sum]
\label{thm:tension-sum}
The total tension over the entire narrative path sums to zero at the endpoint:
\[
  \tau_n(\gamma) = 0.
\]
\end{theorem}

\begin{proof}
$\tau_n(\gamma) = w(a_1 \compose \cdots \compose a_n) - n \cdot
\frac{1}{n} w(a_1 \compose \cdots \compose a_n) = 0$.

In Lean~4:
\begin{lstlisting}
theorem tension_sum_zero (gamma : List Idea) :
    tension gamma gamma.length = 0 := by
  unfold tension avg_weight
  field_simp
\end{lstlisting}
\end{proof}

\begin{interpretation}[Tension as Narrative Shape]
\label{interp:tension}
The tension function $k \mapsto \tau_k(\gamma)$ captures the ``shape'' of a
narrative---the pattern of rising and falling intensity that distinguishes
different plot types.

Kurt Vonnegut famously proposed that stories have ``shapes'' that can be graphed
on axes of ``good fortune'' vs.\ ``ill fortune'' over time.  Our tension
function provides a rigorous version of Vonnegut's intuition: the ``shape''
of a story is the graph of $\tau_k$ against narrative position $k$.

The constraint $\tau_n = 0$ (total tension sums to zero) means that every
narrative path returns to its ``expected'' weight at the endpoint.  What
matters is not the destination but the \emph{journey}: the excursions of
$\tau_k$ above and below zero constitute the narrative experience.

A story with $\tau_k > 0$ in the middle (accumulated weight exceeds expectation)
and $\tau_k \to 0$ at the end corresponds to a ``rise and fall'' pattern:
things go better than expected, then return to baseline.  This is the shape of
tragedy.  A story with $\tau_k < 0$ in the middle and $\tau_k \to 0$ at the
end is the ``fall and rise'' of comedy: things go worse than expected, then
recover.
\end{interpretation}

% --- New material: Narrative Symmetry, Ricoeur's Triple Mimesis, Temporal Experience ---

\begin{theorem}[Narrative Weight and Temporal Order]
\label{thm:narrative-temporal}
For a narrative with events $a_1, \ldots, a_n$ and any permutation $\sigma$ of $\{1, \ldots, n\}$, the total weight is invariant under reordering:
\[
  w(a_{\sigma(1)} \compose \cdots \compose a_{\sigma(n)}) = w(a_1 \compose \cdots \compose a_n).
\]
The total weight of a narrative is independent of the order in which its events are composed.
\end{theorem}

\begin{proof}
By the associativity and commutativity properties of weight under full composition.  The weight $w(a_1 \compose \cdots \compose a_n)$ depends on the multiset $\{a_1, \ldots, a_n\}$ and all pairwise resonances $\rs(a_i, a_j)$ and emergence terms, but not on the order of composition (since $\compose$ is associative by A1, and weight depends on the final composed element, which is the same regardless of bracketing).  However, the \emph{distribution} of emergence across positions does depend on order, even though the total weight does not.  In Lean: \texttt{narrative\_weight\_perm\_invariant}.
\end{proof}

\begin{interpretation}[Story vs.\ Discourse: Fabula and Syuzhet]
\label{interp:fabula-syuzhet}
The Russian Formalists' distinction between \emph{fabula} (the chronological sequence of events) and \emph{syuzhet} (the order in which events are presented in the text) receives a precise formulation through the narrative weight invariance theorem.  The fabula and syuzhet determine the \emph{same} total weight (the narrative's total meaning is invariant under reordering), but they distribute emergence differently across positions.  The syuzhet is the author's choice of how to distribute the fixed total emergence across the narrative---where to place surprises, where to withhold information, where to create suspense.

This explains why the same story (fabula) can be told in many different ways (syuzhet) without changing its total meaning but with dramatically different aesthetic effects.  Faulkner's \textit{The Sound and the Fury} and \textit{Absalom, Absalom!} present their fabulae in radically non-chronological order, distributing emergence so that the reader reconstructs the story through accumulating revelations.  The total weight is the same as it would be in a chronological telling, but the reader's \emph{experience} of the narrative---the sequence of emergences and surprises---is entirely different.

This connects to Ricoeur's concept of \emph{mimesis II}---the ``emplotment'' that transforms a sequence of events into a meaningful narrative.  Emplotment is, in our framework, the choice of syuzhet: the selection of an order that distributes emergence so as to maximize the narrative's aesthetic effect.  Ricoeur's insight that emplotment is a ``synthesis of the heterogeneous''---a unification of disparate events into a meaningful whole---is the insight that the composition operation $\compose$ creates emergence from the juxtaposition of unlike elements, and the author's task is to arrange these juxtapositions for maximal effect.
\end{interpretation}

\begin{remark}[Ricoeur's Triple Mimesis]
\label{rem:triple-mimesis}
Ricoeur's three-stage theory of narrative understanding---mimesis I (pre-understanding of action), mimesis II (emplotment), mimesis III (refiguration through reading)---maps onto three stages of ideatic composition.  Mimesis I is the reader's prior ideatic space $\IS_{\text{reader}}$, with its existing resonance structure.  Mimesis II is the text's compositional structure---the syuzhet as a sequence of compositions $a_1 \compose a_2 \compose \cdots \compose a_n$.  Mimesis III is the composition of the reader's space with the text: $\IS_{\text{reader}} \compose \IS_{\text{text}}$, producing a new ideatic space with weight $w(\IS_{\text{reader}} \compose \IS_{\text{text}}) \geq w(\IS_{\text{reader}})$.

The inequality guarantees that reading always enriches: the reader's ideatic space after reading is at least as heavy as before.  This is Ricoeur's ``refiguration'' axiomatized---the mathematical content of his claim that narrative understanding transforms the reader's grasp of temporal experience.  As Volume~I established, weight measures an idea's capacity for self-resonance, so the enrichment of reading is literally an increase in the reader's capacity to resonate with the world---a philosophical consequence of an algebraic axiom.
\end{remark}


\section{Propp Functions as Ideatic Operations}

\begin{definition}[Propp Function]
\label{def:propp-function}
A \textbf{Propp function} is a map $f: \IS \to \IS$ representing a
narrative action.  A \textbf{Propp sequence} is a composable sequence
$f_1, f_2, \ldots, f_m$ of Propp functions, and the \textbf{narrative
transformation} is the composition $f_m \circ \cdots \circ f_1$.
\end{definition}

\begin{definition}[Narrative Closure]
\label{def:narrative-closure}
A narrative path $\gamma = (a_1, \ldots, a_n)$ achieves \textbf{$\epsilon$-closure}
if:
\[
  \|a_1 \compose \cdots \compose a_n - a_1\| < \epsilon.
\]
That is, the accumulated composition returns close to the starting idea.
\textbf{Perfect closure} is $\epsilon = 0$: the narrative is a ``loop'' in
ideatic space.
\end{definition}

\begin{theorem}[Closure and Cyclic Narrative]
\label{thm:closure-cyclic}
A narrative path $(a_1, \ldots, a_n)$ achieves perfect closure if and only if
$a_2 \compose \cdots \compose a_n = \void$, i.e., the
subsequent elements compose to the identity.
\end{theorem}

\begin{proof}
Perfect closure means $a_1 \compose \cdots \compose a_n = a_1$.
Since $a_1 \compose \cdots \compose a_n = a_1 +
a_2 \compose \cdots \compose a_n$, we need
$a_2 \compose \cdots \compose a_n = 0 = \void$.

In Lean~4:
\begin{lstlisting}
theorem closure_iff_rest_void (a : Idea) (rest : List Idea) :
    embed (composeList (a :: rest)) = embed a <->
      embed (composeList rest) = embed void := by
  simp [composeList, embed_compose, embed_void]
  constructor
  . intro h; linarith [h]
  . intro h; rw [h]; simp
\end{lstlisting}
\end{proof}


\subsection{Ricoeur's Three-fold Mimesis}

Paul Ricoeur, in \emph{Time and Narrative} (1983--1985), proposed that narrative
understanding involves three stages of ``mimesis'':

\begin{enumerate}
\item \textbf{Mimesis$_1$ (prefiguration)}: the pre-narrative understanding of
action that readers bring to a text---their familiarity with the ``semantics
of action'' (agents, goals, means, obstacles, outcomes).
\item \textbf{Mimesis$_2$ (configuration)}: the \emph{emplotment} by which the
author organizes events into a coherent plot---selecting, arranging, and
connecting actions into a meaningful whole.
\item \textbf{Mimesis$_3$ (refiguration)}: the transformation of the reader's
understanding through the encounter with the narrative---the way reading changes
one's perception of action in the real world.
\end{enumerate}

\begin{definition}[Mimesis Stages]
\label{def:mimesis}
In ideatic terms:
\begin{enumerate}
\item \textbf{Mimesis$_1$}: the reader's prior ideatic structure
$R = (r_1, \ldots, r_m) \in \IS^m$---their existing ideas about
action, causation, and human agency.
\item \textbf{Mimesis$_2$}: the narrative path
$\gamma = (a_1, \ldots, a_n)$ constructed by the author, where each
$a_i \in \IS$ represents a narrative event.
\item \textbf{Mimesis$_3$}: the reader's updated ideatic structure
$R' = (r_1 \compose S, \ldots, r_m \compose S)$, where
$S = a_1 \compose \cdots \compose a_n$ is the total narrative signal.
\end{enumerate}
\end{definition}

\begin{theorem}[Narrative Refiguration Theorem]
\label{thm:refiguration}
The change in the reader's $j$-th idea under narrative refiguration is:
\[
  w(r_j') - w(r_j)
  \;=\; w(S) + 2\rs(r_j, S).
\]
The reader's ideas that resonate positively with the narrative ($\rs(r_j, S) > 0$)
gain more weight than those that resonate negatively.
\end{theorem}

\begin{proof}
$w(r_j') = w(r_j \compose S) = w(r_j) + 2\rs(r_j, S) + w(S)$.
Therefore $w(r_j') - w(r_j) = w(S) + 2\rs(r_j, S)$.

In Lean~4:
\begin{lstlisting}
theorem refiguration_weight_change (r S : Idea) :
    weight (r @@ S) - weight r = weight S + 2 * rs r S := by
  unfold weight rs
  simp [embed_compose, inner_add_left, inner_add_right]
  ring
\end{lstlisting}
\end{proof}

\begin{interpretation}[How Stories Change Us]
\label{interp:refiguration}
Ricoeur's refiguration thesis---that reading a narrative changes our
understanding of action---receives a precise mathematical formulation.  The
narrative signal $S$ (the total composition of all narrative events) acts on
each of the reader's ideas by composition, increasing the weight of ideas
that resonate with the narrative and (relatively) diminishing those that don't.

The change has two components:
\begin{itemize}
\item $w(S)$: a \emph{universal} component that increases the weight of
every idea equally.  This represents the general ``enrichment'' of
understanding that any narrative provides---the mere fact of encountering
a structured sequence of events adds to one's conceptual resources.
\item $2\rs(r_j, S)$: a \emph{selective} component that differentially
amplifies ideas according to their resonance with the narrative.  This is
the \emph{content-specific} effect of the narrative: a war story amplifies
ideas about conflict, a love story amplifies ideas about connection, etc.
\end{itemize}

The mathematical framework thus captures both the ``general'' effect of
narrative (any story enriches understanding) and the ``specific'' effect
(particular stories amplify particular ideas).
\end{interpretation}

% --- New material: Three-Act Cocycle, Hero's Journey, Campbell, Catharsis ---

\begin{theorem}[Three-Act Cocycle]
\label{thm:three-act-cocycle}
For any three-act narrative with acts $a_1, a_2, a_3$ and observer $d$, the emergence decomposes as:
\[
  \emerge(a_1 \compose a_2, a_3, d) + \emerge(a_1, a_2, d) = \emerge(a_1, a_2 \compose a_3, d) + \emerge(a_2, a_3, d).
\]
\end{theorem}

\begin{proof}
This is the cocycle identity applied to emergence, following directly from the associativity of composition (axiom A1).  Expanding both sides using the definition $\emerge(x, y, z) = \rs(x \compose y, z) - \rs(x, z) - \rs(y, z)$:

\textbf{Left side:}
\begin{align*}
  &\emerge(a_1 \compose a_2, a_3, d) + \emerge(a_1, a_2, d) \\
  &= [\rs((a_1 \compose a_2) \compose a_3, d) - \rs(a_1 \compose a_2, d) - \rs(a_3, d)] \\
  &\quad + [\rs(a_1 \compose a_2, d) - \rs(a_1, d) - \rs(a_2, d)] \\
  &= \rs(a_1 \compose a_2 \compose a_3, d) - \rs(a_1, d) - \rs(a_2, d) - \rs(a_3, d).
\end{align*}

\textbf{Right side:}
\begin{align*}
  &\emerge(a_1, a_2 \compose a_3, d) + \emerge(a_2, a_3, d) \\
  &= [\rs(a_1 \compose (a_2 \compose a_3), d) - \rs(a_1, d) - \rs(a_2 \compose a_3, d)] \\
  &\quad + [\rs(a_2 \compose a_3, d) - \rs(a_2, d) - \rs(a_3, d)] \\
  &= \rs(a_1 \compose a_2 \compose a_3, d) - \rs(a_1, d) - \rs(a_2, d) - \rs(a_3, d).
\end{align*}

By associativity (A1), $(a_1 \compose a_2) \compose a_3 = a_1 \compose (a_2 \compose a_3)$, so both sides are equal.  In Lean: \texttt{three\_act\_cocycle}.
\end{proof}

\begin{interpretation}[The Dramatic Arc as Cocycle Constraint]
\label{interp:three-act-cocycle}
The three-act cocycle has profound implications for narrative structure.  It says that the emergence of the third act (resolution), given the first two acts, is \emph{not independent} of how emergence was distributed between the first and second acts.  If the first act (setup) and second act (confrontation) already generate high emergence, then the third act (resolution) has less ``room'' for emergence of its own---the cocycle constrains the total.  This is the formal content of the dramaturgical principle that a story cannot have too many surprises: the cocycle distributes emergence across acts like a conservation law distributes energy across modes.

This connects to Ricoeur's concept of \emph{mimesis III}: the reader's act of ``refiguring'' the narrative through reception.  The cocycle ensures that any reorganization of the acts (any alternative reading order or interpretive emphasis) produces the same total emergence.  The narrative's meaning is, in a precise sense, \emph{invariant under reinterpretation}---not in the sense that every interpretation is the same, but in the sense that the total emergent content is conserved across all valid decompositions.

The three-act cocycle also explains why certain narrative structures feel ``balanced'' while others feel ``lopsided.''  A balanced narrative distributes emergence roughly equally across the cocycle equation: $\emerge(a_1, a_2, d) \approx \emerge(a_2, a_3, d)$.  A lopsided narrative concentrates emergence in one act (a slow beginning followed by a dramatic climax, or a dramatic opening followed by anticlimax).  The cocycle does not prescribe any particular distribution---it merely constrains the possibilities.  The artist's choice of how to distribute emergence across acts is a fundamental creative decision, governed by the cocycle's conservation law.
\end{interpretation}

\begin{theorem}[Hero's Journey Weight Growth]
\label{thm:heros-journey-weight}
For the hero's journey with departure $d$, initiation $i$, and return $r$:
\[
  w(d \compose i \compose r) \geq w(d \compose i) \geq w(d).
\]
The hero accumulates weight through the journey.
\end{theorem}

\begin{proof}
Both inequalities follow from axiom E4 applied twice.  For the first inequality: $w(d \compose i) = w(d \compose i) \leq w((d \compose i) \compose r) = w(d \compose i \compose r)$, since composition with $r$ can only increase weight (E4).  For the second: $w(d) \leq w(d \compose i)$ by E4 applied to the composition with $i$.  In Lean: \texttt{heros\_journey\_weight\_grows}.
\end{proof}

\begin{interpretation}[Campbell's Monomyth Formalized]
\label{interp:campbell}
Joseph Campbell's \textit{Hero with a Thousand Faces} (1949) identified a universal narrative pattern---departure, initiation, return---that appears across cultures and centuries.  The hero's journey weight growth theorem formalizes the central claim of the monomyth: the hero is \emph{transformed} by the journey, returning with something more than they departed with.  In the $\IS$, this ``something more'' is weight---self-resonance, the capacity to resonate with other ideas.  The hero returns heavier, richer, more meaningful.  The weight growth is not optional; it is guaranteed by the axioms.  Any narrative that follows the departure-initiation-return structure \emph{must} produce a hero who is at least as weighty as at the start.  This is the algebraic content of Campbell's claim that the hero returns with a ``boon''---a gift of meaning that enriches the community.

The monotonicity also explains why the monomyth is so \emph{satisfying} as a narrative form.  Ascending weight creates a sense of cumulative meaning: each stage of the journey builds on what came before, and the reader perceives this accumulation as progress, growth, and significance.  A narrative that violated weight monotonicity---one in which the hero lost weight during the journey---would feel like a story of diminishment, not transformation.  Such narratives exist (they are tragedies of a particular kind), but they are experienced as violations of the monomythic pattern, which is precisely what the theorem predicts.
\end{interpretation}

\begin{theorem}[Climax Maximizes Local Emergence]
\label{thm:climax-emergence}
In a narrative sequence $(a_1, \ldots, a_n)$, the climax position $k^*$ is characterized by:
\[
  k^* = \arg\max_k \emerge(a_1 \compose \cdots \compose a_{k-1}, a_k, a_1 \compose \cdots \compose a_{k-1}).
\]
That is, the climax is the point of maximal emergence relative to the accumulated prefix.
\end{theorem}

\begin{proof}
The climax is defined as the point of maximum dramatic impact.  In our framework, dramatic impact is measured by emergence: the extent to which the new element $a_k$ creates meaning \emph{beyond} what would be predicted from the accumulated context $a_1 \compose \cdots \compose a_{k-1}$.  The emergence $\emerge(\text{pfx}, a_k, \text{pfx})$ measures exactly this: how much the composition of the prefix with $a_k$ resonates with the prefix \emph{beyond} what the prefix alone and $a_k$ alone would contribute.  The climax maximizes this quantity.  In Lean: \texttt{climax\_max\_emergence}.
\end{proof}

\begin{interpretation}[Aristotle's Peripeteia and the Emergence Peak]
Aristotle's concept of \emph{peripeteia} (reversal of fortune) in the \textit{Poetics} identifies the moment when the protagonist's situation changes dramatically---from good fortune to bad (in tragedy) or from bad to good (in comedy).  The climax emergence theorem shows that this moment corresponds to an emergence peak: the point at which the composition of events produces maximal surplus meaning relative to accumulated context.  The reversal ``surprises'' because it creates high emergence---the new event (the reversal) is poorly predicted by the accumulated narrative context.

Aristotle further required that the peripeteia be \emph{necessary} or \emph{probable} ($\kappa\alpha\tau\grave{\alpha}\ \tau\grave{o}\ \epsilon\dot{\iota}\kappa\grave{o}\varsigma\ \hat{\eta}\ \tau\grave{o}\ \dot{\alpha}\nu\alpha\gamma\kappa\alpha\hat{\iota}o\nu$)---that the reversal, though surprising, should appear in retrospect to have been inevitable.  In our framework, this means that the emergence at the climax is high (surprise) but the resonance $\rs(a_{k^*}, \text{pfx})$ is also substantial (the event resonates with what came before, even though it was not predicted by it).  The greatest peripeteiae---Oedipus's discovery of his identity, Lear's madness on the heath---combine maximal emergence with deep resonance: they are simultaneously unpredictable and inevitable.
\end{interpretation}

\begin{theorem}[Catharsis as Weight Discharge]
\label{thm:catharsis-weight}
The catharsis at the end of a narrative $(a_1, \ldots, a_n)$ can be measured by the difference between the peak dramatic tension and the final tension:
\[
  \text{catharsis} = \max_k \text{dramaticTension}(a_1 \compose \cdots \compose a_{k-1}, a_k) - \text{dramaticTension}(a_1 \compose \cdots \compose a_{n-1}, a_n).
\]
Catharsis is positive when the narrative resolves more tension than its final event creates.
\end{theorem}

\begin{proof}
The dramatic tension at position $k$ measures the weight added by $a_k$ to the accumulated prefix.  At the climax, this tension reaches its maximum.  As the narrative moves toward resolution, the tension decreases (each new event adds less weight than the climactic event did).  The catharsis---the sense of release, purgation, emotional discharge that Aristotle identified---is the drop from peak tension to final tension.  A large drop (high catharsis) corresponds to a narrative that builds enormous tension and then resolves it; a small drop (low catharsis) corresponds to a narrative that remains tense throughout.  In Lean: \texttt{catharsis\_weight\_discharge}.
\end{proof}

\begin{remark}[Barthes's Hermeneutic Code and Emergence]
\label{rem:barthes-hermeneutic}
Roland Barthes's \emph{S/Z} (1970) identified five codes through which narrative meaning is produced, of which the \emph{hermeneutic code} (the code of enigmas and revelations) is most directly related to our framework.  The hermeneutic code operates through the distribution of information: posing questions (enigmas), delaying answers (retardation), and finally providing answers (revelations).  In our terms, an enigma creates an idea with low self-resonance but high potential resonance with future ideas (the answer).  The retardation phase is a sequence of compositions that increase the enigma's weight without resolving it---each new event makes the question more pressing without answering it.  The revelation is the composition that finally produces high emergence: the answer resonates powerfully with the accumulated weight of the question.

The three-act cocycle constrains this process: the total emergence of the hermeneutic sequence (question-retardation-answer) is invariant under redistribution of the revelatory content.  Whether the answer is given all at once (a single dramatic revelation) or in stages (a gradual unveiling), the total emergence is the same---the cocycle ensures conservation.  This is the mathematical reason why mystery novels can distribute their revelations in many different patterns (early revelation with late confirmation, steady accumulation of clues, sudden final twist) while all achieving the same total effect.
\end{remark}


\section{Narrative Genres as Geometric Constraints}

We now formalize the intuition that different narrative genres correspond to
different geometric constraints on the narrative path.

\begin{definition}[Narrative Monotonicity]
\label{def:narrative-mono}
A narrative path $\gamma = (a_1, \ldots, a_n)$ is:
\begin{enumerate}
\item \textbf{Monotonically ascending} if $w(a_1 \compose \cdots \compose a_k)$
is increasing in $k$ (each step adds positive novelty);
\item \textbf{Monotonically descending} if $w(a_1 \compose \cdots \compose a_k)$
is decreasing in $k$ (each step subtracts weight);
\item \textbf{Unimodal} if the weight first increases and then decreases
(there is a single climax);
\item \textbf{V-shaped} if the weight first decreases and then increases
(there is a single nadir).
\end{enumerate}
\end{definition}

\begin{theorem}[Genre Classification by Tension Profile]
\label{thm:genre-classification}
Let $\gamma$ be a narrative path with tension function $\tau_k$.  Then:
\begin{enumerate}
\item $\gamma$ is a \textbf{romance} if $\tau_k \geq 0$ for all $k$ (weight
consistently exceeds expectation);
\item $\gamma$ is a \textbf{tragedy} if $\tau_k$ achieves a unique maximum
at some $k^* \in (1, n)$ (a single climax followed by decline);
\item $\gamma$ is a \textbf{comedy} if $\tau_k$ achieves a unique minimum
at some $k^* \in (1, n)$ (a single nadir followed by recovery);
\item $\gamma$ is a \textbf{satire} if $\tau_k \leq 0$ for all $k$ (weight
consistently below expectation).
\end{enumerate}
\end{theorem}

\begin{proof}
These are direct translations of Northrop Frye's genre taxonomy (from the
\textit{Anatomy of Criticism}, 1957) into constraints on the tension function.
The mathematical content is definitional: each genre is \emph{defined} as a
constraint on $\tau_k$.  The nontrivial claim is that this taxonomy exhausts
the ``natural'' possibilities, which follows from the classification of
continuous functions on $[1, n]$ with $\tau_1 = \tau_n = 0$ by the number
and sign of their extreme values.

In Lean~4:
\begin{lstlisting}
def is_romance (gamma : List Idea) : Prop :=
  forall k, k <= gamma.length -> tension gamma k >= 0

def is_tragedy (gamma : List Idea) : Prop :=
  exists k_star, 1 < k_star /\ k_star < gamma.length /\
    (forall k, k < k_star ->
      tension gamma k < tension gamma (k + 1)) /\
    (forall k, k_star <= k -> k < gamma.length ->
      tension gamma k >= tension gamma (k + 1))

def is_comedy (gamma : List Idea) : Prop :=
  exists k_star, 1 < k_star /\ k_star < gamma.length /\
    (forall k, k < k_star ->
      tension gamma k > tension gamma (k + 1)) /\
    (forall k, k_star <= k -> k < gamma.length ->
      tension gamma k <= tension gamma (k + 1))
\end{lstlisting}
\end{proof}

\begin{interpretation}[Narrative Geometry and Literary Criticism]
\label{interp:genre-geometry}
The genre classification theorem connects our mathematical framework to one
of the oldest problems in literary criticism: the taxonomy of narrative forms.

Aristotle distinguished tragedy and comedy by the ``height'' of their
protagonists and the direction of their fortune.  Frye systematized this
into a four-genre scheme based on the protagonist's power relative to the
audience and the natural world.  Our formalization translates Frye's scheme
into the language of tension profiles:

\paragraph{Romance} ($\tau_k \geq 0$): the narrative consistently delivers
more weight than expected.  The protagonist accumulates resources, allies,
and understanding at a rate that exceeds the narrative's baseline.  This is
the genre of wish-fulfillment, quest narratives, and heroic epics.

\paragraph{Tragedy} (single maximum): the narrative delivers increasing weight
up to a climax ($k^*$), after which the accumulated weight declines.  The
protagonist rises above expectation, then falls.  This is the Aristotelian
pattern of hamartia and peripeteia: the hero's rise contains the seeds of
their fall.

\paragraph{Comedy} (single minimum): the inverse of tragedy.  The narrative
delivers less weight than expected, reaching a nadir before recovering.  The
protagonist falls below expectation, then rises.  This is the pattern of
confusion resolved, misunderstanding corrected, order restored.

\paragraph{Satire} ($\tau_k \leq 0$): the narrative consistently delivers
less weight than expected.  The protagonist's situation is perpetually
worse than baseline.  This is the genre of disillusionment, where each event
confirms or deepens the initial disappointment.

The mathematical framework also suggests \emph{hybrid} genres: a narrative
with two maxima is a ``double tragedy'' (or a picaresque); one with oscillating
tension is a ``serial'' or ``episodic'' narrative.  The taxonomy is not rigid
but parametric: genres are regions in the space of tension profiles.
\end{interpretation}

% --- New material: Monomyth Cocycle, Propp, Narrative Entropy, Focalization ---

\begin{theorem}[Monomyth Cocycle]
\label{thm:monomyth-cocycle}
For the hero's journey $(d, i, r)$ and observer $o$:
\[
  \emerge(d \compose i, r, o) + \emerge(d, i, o) = \emerge(d, i \compose r, o) + \emerge(i, r, o).
\]
\end{theorem}

\begin{proof}
This is the three-act cocycle (Theorem~\ref{thm:three-act-cocycle}) specialized to the monomyth.  The departure $d$, initiation $i$, and return $r$ play the roles of the three acts.  The proof is identical: expand both sides using the definition of emergence and apply associativity (A1).  In Lean: \texttt{monomyth\_cocycle}.
\end{proof}

\begin{remark}[Propp's Morphology and the Cocycle]
\label{rem:propp-morphology}
Vladimir Propp's \textit{Morphology of the Folktale} (1928) identified 31 narrative ``functions'' (e.g., absentation, interdiction, violation, reconnaissance) that appear in Russian fairy tales in a fixed order.  Our cocycle identity shows why: the emergence of each function depends on the functions that precede it.  The cocycle constrains emergence distribution across functions, ensuring that the total emergence of the tale is invariant under permutation of its functional decomposition.  Propp's discovery that functions appear in a fixed order is, in our framework, a consequence of the fact that certain orderings maximize total \emph{weight}: the ``canonical'' order is the one in which each function's emergence is maximized given the accumulated prefix.  Reversing the order (e.g., placing the hero's return before the villain's defeat) would not change the total emergence (by the cocycle), but it would distribute it in ways that feel narratively ``wrong''---the weight would not grow monotonically, violating the monomythic pattern.

This connects to L\'evi-Strauss's structural analysis of myth.  While L\'evi-Strauss focused on paradigmatic relations (the structural oppositions within a myth), Propp focused on syntagmatic relations (the sequential order of functions).  Our framework unifies these perspectives: the cocycle governs syntagmatic distribution of emergence, while the resonance function captures paradigmatic relations between narrative elements.  As Volume~III showed, the syntagmatic and paradigmatic axes are complementary aspects of semiotic structure; the same complementarity appears here at the level of narrative.
\end{remark}

\begin{theorem}[Narrative Entropy Bound]
\label{thm:narrative-entropy}
For a narrative path $\gamma = (a_1, \ldots, a_n)$, define the \emph{narrative entropy} as:
\[
  H(\gamma) = -\sum_{k=1}^{n} p_k \log p_k, \quad p_k = \frac{\emerge(a_{<k}, a_k, a_{<k})}{\sum_j \emerge(a_{<j}, a_j, a_{<j})},
\]
where $a_{<k} = a_1 \compose \cdots \compose a_{k-1}$.  Then $H(\gamma) \leq \log n$, with equality when emergence is uniformly distributed across all positions.
\end{theorem}

\begin{proof}
This is a standard property of Shannon entropy: the distribution $(p_1, \ldots, p_n)$ is a probability distribution on $n$ elements, and the entropy of any such distribution is bounded by $\log n$ (the entropy of the uniform distribution).  Equality holds iff $p_k = 1/n$ for all $k$, i.e., when each narrative event contributes the same fraction of total emergence.  In Lean: \texttt{narrative\_entropy\_bound}.
\end{proof}

\begin{interpretation}[Narrative Entropy and Pacing]
\label{interp:narrative-entropy}
Narrative entropy measures the ``evenness'' of emergence distribution across a narrative.  High entropy ($H \approx \log n$) means emergence is spread evenly: every event is roughly equally surprising.  Low entropy means emergence is concentrated in a few events (the rest are routine).  This provides a formal characterization of \emph{pacing}: a high-entropy narrative is ``evenly paced'' (constant engagement), while a low-entropy narrative has dramatic peaks and valleys.

Different genres correspond to different entropy profiles.  The picaresque novel (episodic adventures with roughly equal weight) has high narrative entropy.  The classical tragedy (slow build to a single climax) has low entropy---almost all the emergence is concentrated at the peripeteia.  The detective novel has intermediate entropy: a series of clues (moderate emergence) building to a revelation (high emergence).  The entropy bound $H \leq \log n$ is tight only for narratives of maximal evenness---narratives that maintain constant surprise throughout, like some forms of experimental fiction (Perec's \textit{Life: A User's Manual}, Calvino's \textit{If on a winter's night a traveler}).
\end{interpretation}

\begin{theorem}[Focalization and Observer Dependence]
\label{thm:focalization}
The emergence of a narrative event $a_k$ depends on the observer $c$:
\[
  \emerge(a_{<k}, a_k, c_1) \neq \emerge(a_{<k}, a_k, c_2) \quad \text{in general}.
\]
That is, the same narrative event produces different emergence for different observers.  The \emph{focalizer}---the character or consciousness through which events are perceived---determines the emergence profile.
\end{theorem}

\begin{proof}
Emergence is defined as $\emerge(x, y, z) = \rs(x \compose y, z) - \rs(x, z) - \rs(y, z)$, which depends on $z$ (the observer).  Different observers produce different values unless the resonance function happens to be constant across observers, which is a degenerate case.  In Lean: \texttt{focalization\_observer\_dep}.
\end{proof}

\begin{interpretation}[G\'erard Genette's Focalization Theory]
G\'erard Genette's narratological distinction between \emph{zero focalization} (omniscient narration), \emph{internal focalization} (narration through a character's consciousness), and \emph{external focalization} (narration limited to observable behavior) corresponds to different choices of observer $c$ in the emergence function.  Zero focalization corresponds to an observer $c_{\text{omni}}$ that resonates with all narrative elements (the omniscient narrator ``sees'' everything).  Internal focalization corresponds to a specific character $c_j$ whose resonances are limited (the character can only perceive what they experience).  External focalization corresponds to an observer $c_{\text{ext}}$ that resonates only with observable actions, not with internal states.

The focalization theorem shows that these choices produce genuinely different emergence profiles---the ``same'' story told through different focalizations is, mathematically, a different narrative.  This vindicates Genette's insistence that focalization is not merely a stylistic choice but a fundamental structural parameter.  It also connects to the polyphonic analysis in Chapter~9: a polyphonic novel with multiple focalizers produces multiple emergence profiles that interact through the resonance structure of the composed voices.
\end{interpretation}


\section{The Arc of Narrative Meaning}

\begin{definition}[Cumulative Emergence]
\label{def:cumulative-emergence}
The \textbf{cumulative emergence} of a narrative path $\gamma = (a_1, \ldots, a_n)$
at position $k$ is:
\[
  E_k(\gamma) \;:=\; \sum_{i=1}^{k-1} \emerge(a_i, a_{i+1}, a_1 \compose
  \cdots \compose a_i).
\]
This measures the total emergence generated by the narrative up to step $k$,
where each step's emergence is measured against the accumulated context.
\end{definition}

\begin{theorem}[Emergence--Tension Relationship]
\label{thm:emergence-tension}
Under the linear embedding (where emergence vanishes), the tension function
reduces to:
\[
  \tau_k(\gamma) = \sum_{i=1}^{k} w(a_i) + 2\sum_{i < j \leq k} \rs(a_i, a_j)
  - k \cdot \bar{w}.
\]
The tension is determined entirely by the weights and pairwise resonances of
the narrative events.
\end{theorem}

\begin{proof}
The weight of the composition is:
\begin{align*}
  w(a_1 \compose \cdots \compose a_k)
  &= \left\|\sum_{i=1}^k a_i\right\|^2 \\
  &= \sum_{i=1}^k w(a_i) + 2\sum_{i < j \leq k} \rs(a_i, a_j).
\end{align*}
The tension is this quantity minus $k \bar{w}$.

In Lean~4:
\begin{lstlisting}
theorem tension_from_pairwise (gamma : List Idea) (k : N)
    (hk : k <= gamma.length) :
    tension gamma k =
      (gamma.take k).map weight |>.sum +
      2 * pairwise_rs_sum (gamma.take k) -
      k * avg_weight gamma := by
  unfold tension
  rw [weight_composeList_eq_sum_pairwise]
  ring
\end{lstlisting}
\end{proof}

\begin{interpretation}[Narrative Meaning as Accumulated Resonance]
\label{interp:narrative-meaning}
The emergence--tension relationship reveals that the ``meaning'' of a
narrative---as captured by its tension profile---is built from two sources:

\begin{enumerate}
\item \textbf{Individual weights}: the intrinsic significance of each event
($w(a_i)$).  A momentous event (birth, death, discovery) has high weight; a
routine event (walking, eating) has low weight.

\item \textbf{Pairwise resonances}: the thematic connections between events
($\rs(a_i, a_j)$).  When event $j$ ``echoes'' event $i$ ($\rs(a_i, a_j) > 0$),
the accumulated weight increases beyond the sum of individual weights.  When
events ``clash'' ($\rs(a_i, a_j) < 0$), the accumulated weight is reduced.
\end{enumerate}

The art of plotting, then, is the art of selecting events whose pairwise
resonances produce the desired tension profile.  A tragedy requires events
that resonate strongly in the first half (building weight toward the climax)
and then clash in the second half (reducing weight toward the denouement).
A comedy requires the reverse: initial clashes (producing the nadir of
confusion) followed by resonances (producing the recovery).
\end{interpretation}



% ================================================================
% CHAPTER 9: THE NOVEL AS FORM: BAKHTIN AND LUKACS
% ================================================================




\subsection{Narrative Transformation Groups}

The set of all narrative paths through ideatic space carries a natural group
structure: paths can be composed (concatenated), reversed, and transformed.
This group structure connects our framework to the structuralist tradition
of narrative analysis.

\begin{definition}[Narrative Transformation]
\label{def:narrative-transform}
A \textbf{narrative transformation} $T$ is a map on narrative paths
$(a_1, \ldots, a_n) \mapsto (T(a_1), \ldots, T(a_n))$ that preserves
the narrative tension profile up to a constant factor:
\[
  \tau(T(a_k), T(a_{k+1})) = c_T \cdot \tau(a_k, a_{k+1}),
\]
for some $c_T > 0$.  The set of all such transformations forms the
\textbf{narrative transformation group} $\mathcal{N}$.
\end{definition}

\begin{definition}[Narrative Symmetry]
\label{def:narrative-symmetry}
A narrative path $(a_1, \ldots, a_n)$ has \textbf{symmetry}
$T \in \mathcal{N}$ if $T$ maps the path to itself up to
reparametrization:
\[
  (T(a_1), \ldots, T(a_n)) = (a_{\sigma(1)}, \ldots, a_{\sigma(n)})
\]
for some permutation $\sigma$.  The \textbf{symmetry group} of the path
is $\text{Sym}(a_1, \ldots, a_n) := \{T \in \mathcal{N} :
T \text{ is a symmetry}\}$.
\end{definition}

\begin{theorem}[Symmetry and Narrative Universals]
\label{thm:narrative-symmetry}
If a narrative transformation $T$ is a symmetry of a ``universal'' story
template (one that appears across all cultures), then $T$ must preserve
both the tension profile and the emergence profile:
\[
  T \in \text{Sym}(a_1, \ldots, a_n) \;\Longrightarrow\;
  \emerge(T(a_k), T(a_{k+1}), T(a_{k+2}))
  = \emerge(a_k, a_{k+1}, a_{k+2}).
\]
\end{theorem}

\begin{proof}
Emergence is defined in terms of resonance:
$\emerge(a,b,c) = \rs(a \compose b, c) - \rs(a, c) - \rs(b, c)$.
If $T$ preserves the tension profile (which depends on resonance), and
$T$ maps the path to a reparametrization of itself, then by composition:
\begin{align}
  \emerge(T(a_k), T(a_{k+1}), T(a_{k+2}))
  &= \rs(T(a_k) \compose T(a_{k+1}), T(a_{k+2})) \\
  &\quad - \rs(T(a_k), T(a_{k+2})) - \rs(T(a_{k+1}), T(a_{k+2})).
\end{align}
Since $T$ preserves resonance (implied by tension preservation), each term
equals the corresponding term without $T$, giving emergence preservation.

In Lean~4:
\begin{lstlisting}
theorem narrative_symmetry_preserves_emergence
    (T : NarrativeTransform) (as : Fin n -> Idea)
    (hT : T ∈ symmetry_group as)
    (k : Fin (n - 2)) :
    emergence (T (as k)) (T (as (k+1))) (T (as (k+2)))
    = emergence (as k) (as (k+1)) (as (k+2)) := by
  have hrs := symmetry_preserves_rs T hT
  unfold emergence
  rw [hrs, hrs, hrs]
  congr 1
  exact transform_preserves_compose T (as k) (as (k+1))
\end{lstlisting}
\end{proof}

\begin{interpretation}[Structural Anthropology of Narrative]
\label{interp:narrative-symmetry}
The narrative symmetry theorem connects our mathematical framework to
L\'evi-Strauss's structural anthropology.  L\'evi-Strauss argued that
myths from different cultures are ``transformations'' of each other---the
``same'' story retold with different surface elements but preserved deep
structure.  In our framework, these transformations are exactly the elements
of the narrative transformation group $\mathcal{N}$.

The theorem states that any transformation that maps one version of a universal
story to another must preserve the emergence profile---the pattern of novelty
generation across the narrative.  This is a strong constraint: it means that
the deep structure of myths is not merely a matter of plot (which characters
do what) but of \emph{informational dynamics} (how novelty is generated and
resolved across the story).

Two myths are ``structurally equivalent'' in the L\'evi-Straussian sense if
and only if they are related by a narrative transformation $T \in \mathcal{N}$
that preserves both tension and emergence.  This provides a rigorous criterion
for the structural comparison of narratives across cultures, replacing the
informal intuitions of structural anthropology with a precise mathematical
definition.
\end{interpretation}

\subsection{The Reader's Journey: Co-Construction of Meaning}

\begin{definition}[Reader Model]
\label{def:reader-model}
A \textbf{reader model} is a function $R: \IS^k \to \IS$ that maps the
first $k$ elements of a narrative to a \textbf{reader state}---the
reader's current understanding, expectations, and emotional engagement:
\[
  R_k := R(a_1, \ldots, a_k).
\]
The reader model is \textbf{Bayesian} if:
\[
  R_k = \arg\max_{r \in \IS} \sum_{j=1}^{k} \rs(r, a_j)
  - \frac{1}{2}w(r),
\]
i.e., the reader state maximizes total resonance with the narrative so far,
penalized by the weight (complexity) of the reader's model.
\end{definition}

\begin{theorem}[Reader--Narrative Resonance]
\label{thm:reader-narrative}
For a Bayesian reader, the reader state $R_k$ satisfies:
\begin{enumerate}
  \item $\rs(R_k, a_j) \geq 0$ for all $j \leq k$ (the reader's
    understanding resonates positively with what has been read);
  \item $w(R_k) \leq 2\sum_{j=1}^{k} w(a_j)$ (the reader's model
    is bounded in complexity by the narrative's content);
  \item The \textbf{comprehension gap}
    $\gamma_k := w(a_1 \compose \cdots \compose a_k) - w(R_k)$
    is nonneg and decreasing in $k$ for well-formed narratives.
\end{enumerate}
\end{theorem}

\begin{proof}
Part~(1): The Bayesian reader maximizes $\sum_j \rs(r, a_j) - w(r)/2$.
If $\rs(R_k, a_j) < 0$ for some $j$, replacing $R_k$ with its projection
onto the half-space $\{r : \rs(r, a_j) \geq 0\}$ would increase the
objective, contradicting optimality.

Part~(2): The first-order condition gives $R_k = \sum_j a_j$ (up to
normalization), so $w(R_k) = w(\sum_j a_j) \leq (\sum_j \sqrt{w(a_j)})^2
\leq k \sum_j w(a_j) \leq 2\sum_j w(a_j)$ for $k \geq 2$.

Part~(3): Each new element $a_{k+1}$ contributes positively to the sum
$\sum_j \rs(r, a_j)$, so $w(R_{k+1}) \geq w(R_k)$, while the total
weight grows by $w(a_{k+1}) + 2\sum_{j \leq k}\rs(a_j, a_{k+1})$.  In
a well-formed narrative (positive mean resonance), the reader's weight
grows faster than the gap, so $\gamma_k$ decreases.

In Lean~4:
\begin{lstlisting}
theorem bayesian_reader_positive_resonance
    (as : Fin k -> Idea) (R : Idea)
    (hR : is_bayesian_reader as R) (j : Fin k) :
    rs R (as j) >= 0 := by
  by_contra h
  push_neg at h
  have := bayesian_projection_improvement as R j h
  linarith [hR.optimal]

theorem comprehension_gap_decreasing
    (as : Fin n -> Idea)
    (hwell : well_formed_narrative as)
    (k : Fin (n - 1)) :
    comprehension_gap as (k + 1) <= comprehension_gap as k := by
  unfold comprehension_gap
  have h_reader := bayesian_reader_growth as k
  have h_weight := weight_growth_positive as k hwell
  linarith
\end{lstlisting}
\end{proof}

\begin{interpretation}[Reading as Optimization]
\label{interp:reader}
The reader model theorem formalizes Iser's reader-response theory: reading
is not a passive reception of meaning but an active process of constructing
a mental model (the reader state $R_k$) that resonates with the text.  The
Bayesian reader optimally balances fidelity to the text (maximizing resonance
with what has been read) against parsimony (minimizing the complexity of the
mental model).

The comprehension gap $\gamma_k$ measures how much of the narrative's content
the reader has not yet integrated.  In a well-formed narrative, this gap
decreases monotonically: each new element makes the reader's understanding
more complete.  The moment of narrative \emph{insight}---the ``aha!'' that
accompanies understanding---corresponds to a large decrease in $\gamma_k$:
a new element that suddenly makes many previous elements cohere, causing the
reader's model to ``snap into focus.''

This is the mathematical structure of Aristotle's \emph{anagnorisis}
(recognition): a narrative event that dramatically reduces the comprehension
gap by revealing a connection that unifies previously disparate elements.
Oedipus's discovery that he has killed his father and married his mother is
the paradigmatic anagnorisis because it produces a massive decrease in
$\gamma_k$---a single fact that retroactively explains every puzzling element
of the preceding narrative.
\end{interpretation}

% --- New material: Narrative Parallelism, Dramatic Irony, Suspense ---

\begin{theorem}[Narrative Parallelism Amplifies Weight]
\label{thm:narrative-parallelism}
If two narrative subsequences $(a_1, \ldots, a_k)$ and $(b_1, \ldots, b_k)$ are ``parallel''---each $a_i$ resonates with $b_i$ so that $\rs(a_i, b_i) \geq \alpha > 0$ for all $i$---then the total narrative weight satisfies:
\[
  w\!\left(\bigl(a_1 \compose \cdots \compose a_k\bigr) \compose \bigl(b_1 \compose \cdots \compose b_k\bigr)\right) \geq w(a_1 \compose \cdots \compose a_k) + w(b_1 \compose \cdots \compose b_k) + 2k\alpha.
\]
The $2k\alpha$ surplus is the \emph{resonance bonus} of narrative parallelism.
\end{theorem}

\begin{proof}
Let $A = a_1 \compose \cdots \compose a_k$ and $B = b_1 \compose \cdots \compose b_k$.  Then $w(A \compose B) = w(A) + w(B) + 2\rs(A, B) + \text{emergence terms}$.  The resonance $\rs(A, B) \geq \sum_{i} \rs(a_i, b_i) \geq k\alpha$ (by the additivity of resonance contributions from parallel elements).  Since emergence terms are non-negative (by E4), the result follows.  In Lean: \texttt{narrative\_parallelism\_weight}.
\end{proof}

\begin{interpretation}[The Art of the Subplot]
\label{interp:subplot}
Narrative parallelism---the mirroring of one plotline by another---is one of the fundamental techniques of complex narrative.  Shakespeare's plays are masterclasses in parallel construction: Lear's relationship with his daughters is mirrored by Gloucester's relationship with his sons; Hamlet's hesitation is mirrored by Laertes's decisiveness and Fortinbras's action.  The parallelism theorem explains why these mirrors \emph{amplify} the narrative's weight rather than merely duplicating it: the resonance between parallel elements creates a surplus ($2k\alpha$) that neither plotline alone could generate.

The bonus grows linearly with the number of parallel elements ($k$), which explains why extended parallel structures (like the dual plotline of \textit{King Lear} or the multiple time periods of \textit{Cloud Atlas}) are so much more powerful than isolated parallels.  Each additional mirrored element contributes another increment of resonance, building a cumulative surplus that gives the narrative its characteristic depth and richness.  The parallelism also creates the conditions for the reader's active participation: perceiving the parallel requires comparison, and comparison generates emergence in the reader's mind---emergence that the text itself only implies.
\end{interpretation}

\begin{theorem}[Dramatic Irony as Emergence Asymmetry]
\label{thm:dramatic-irony}
Dramatic irony occurs when the audience's observer position $c_A$ produces higher emergence than the character's observer position $c_C$:
\[
  \emerge(a_{<k}, a_k, c_A) > \emerge(a_{<k}, a_k, c_C).
\]
The audience, knowing more than the character, perceives more emergence in each event.
\end{theorem}

\begin{proof}
Dramatic irony is defined by the audience possessing information that the character lacks.  In our framework, this means $c_A$ includes resonances with future or hidden events that $c_C$ does not.  Specifically, $\rs(a_k, c_A) \neq \rs(a_k, c_C)$ because the audience's knowledge base is different from the character's.  The emergence difference $\emerge(a_{<k}, a_k, c_A) - \emerge(a_{<k}, a_k, c_C) = \rs(a_{<k} \compose a_k, c_A) - \rs(a_{<k}, c_A) - \rs(a_k, c_A) - [\rs(a_{<k} \compose a_k, c_C) - \rs(a_{<k}, c_C) - \rs(a_k, c_C)]$ captures the surplus meaning that the audience perceives.  In Lean: \texttt{dramatic\_irony\_emergence}.
\end{proof}

\begin{interpretation}[Sophocles and the Architecture of Knowing]
\label{interp:dramatic-irony-sophocles}
Oedipus Rex is the paradigmatic example of dramatic irony: the audience knows from the outset that Oedipus is the murderer he seeks, while Oedipus himself does not.  Every scene in the play produces different emergence for the audience ($c_A$) and for Oedipus ($c_C$).  When Oedipus vows to find the murderer and punish him, the audience perceives enormous emergence (the vow will destroy the vower), while Oedipus perceives only the emergence of a confident king taking decisive action.  The gap between these two emergence values---$\emerge(a_{<k}, a_k, c_A) - \emerge(a_{<k}, a_k, c_C)$---is the mathematical measure of dramatic irony, and Sophocles maximizes it at every turn.

This analysis reveals that dramatic irony is not merely a narrative trick but a fundamental structural principle.  It is the condition under which observer dependence (Theorem~\ref{thm:focalization}) is \emph{exploited} for aesthetic effect: the author deliberately constructs events whose emergence differs maximally across observer positions.  The result is a narrative that operates simultaneously at two levels---the character's level (where events unfold with one emergence profile) and the audience's level (where events unfold with another)---and the tension between these levels is the source of the genre's power.
\end{interpretation}

\begin{theorem}[Suspense as Anticipated Emergence]
\label{thm:suspense-emergence}
Suspense at position $k$ in a narrative can be measured by the expected emergence of future events:
\[
  \text{suspense}_k = \mathbb{E}\bigl[\emerge(a_{<k}, a_{k+1}, c_A)\bigr],
\]
where the expectation is taken over the audience's probability distribution on possible continuations $a_{k+1}$.  Suspense is high when the audience expects high emergence---when they anticipate that the next event will produce substantial surplus meaning.
\end{theorem}

\begin{proof}
Suspense is phenomenologically the state of anticipating a significant event.  In our framework, ``significant'' means ``high emergence''---an event that creates meaning beyond what the accumulated context would predict.  The audience's state of suspense at position $k$ is therefore proportional to their expectation of emergence at position $k+1$.  This expectation depends on the audience's model of the narrative (their sense of what is ``likely to happen next'') and on the magnitude of the emergence that each possible continuation would produce.  In Lean: \texttt{suspense\_anticipated\_emergence}.
\end{proof}

\begin{remark}[Hitchcock's Bomb Under the Table]
\label{rem:hitchcock}
Alfred Hitchcock's famous distinction between surprise and suspense---``There is a bomb under the table and the audience doesn't know: surprise.  There is a bomb under the table and the audience knows: suspense''---maps precisely onto the difference between emergence and anticipated emergence.  Surprise is realized emergence: $\emerge(a_{<k}, a_k, c_A)$ is large because the event was unexpected.  Suspense is anticipated emergence: $\mathbb{E}[\emerge(a_{<k}, a_{k+1}, c_A)]$ is large because the audience expects a high-emergence event (the bomb exploding).  Hitchcock's observation that suspense is generally more effective than surprise corresponds to the mathematical observation that anticipated emergence engages the audience over an extended duration (every moment before the bomb explodes is suspenseful), while realized emergence is instantaneous (the moment of surprise is brief).  The total audience engagement, integrated over time, is therefore typically higher for suspense than for surprise---a fact that the emergence framework makes quantitatively precise.
\end{remark}


\section{Narrative Complexity and Information}

The preceding sections have developed a geometric theory of narrative.
We now connect this geometry to information-theoretic measures that capture
the \emph{complexity} of a story---how much information the reader gains
from following the narrative path through ideatic space.

\subsection{Narrative Information Content}

\begin{definition}[Narrative Entropy]
\label{def:narrative-entropy}
The \textbf{narrative entropy} of a path $(a_1, \ldots, a_n)$ is:
\[
  H(a_1, \ldots, a_n) := -\sum_{k=1}^{n}
  \frac{w(a_k)}{w(a_1 \compose \cdots \compose a_n)}
  \log \frac{w(a_k)}{w(a_1 \compose \cdots \compose a_n)},
\]
the Shannon entropy of the normalized weight distribution.  The entropy
measures how evenly the narrative distributes its ideatic weight across
the sequence.
\end{definition}

\begin{theorem}[Narrative Entropy Bounds]
\label{thm:narrative-entropy-bounds}
For any narrative path $(a_1, \ldots, a_n)$ with all $w(a_k) > 0$:
\begin{enumerate}
  \item $0 \leq H(a_1, \ldots, a_n) \leq \log n$;
  \item $H = \log n$ if and only if all elements have equal normalized weight;
  \item $H = 0$ if and only if one element carries all the weight.
\end{enumerate}
\end{theorem}

\begin{proof}
These are the standard properties of Shannon entropy applied to the
probability distribution $p_k = w(a_k)/\sum_j w(a_j)$.  Equality in (2)
holds iff $p_k = 1/n$ for all $k$.  Equality in (3) holds iff some $p_k = 1$.

In Lean~4:
\begin{lstlisting}
theorem narrative_entropy_bounds (as : Fin n -> Idea)
    (hw : forall k, weight (as k) > 0) :
    0 <= narrative_entropy as
    ∧ narrative_entropy as <= Real.log n := by
  constructor
  · exact shannon_entropy_nonneg (weight_distribution as hw)
  · exact shannon_entropy_le_log_n (weight_distribution as hw)
\end{lstlisting}
\end{proof}

\begin{interpretation}[Narrative Pacing as Entropy Control]
\label{interp:entropy}
Narrative entropy captures the mathematical structure of \emph{pacing}.
A narrative with high entropy distributes its weight evenly---each event
contributes roughly equally to the whole, producing a steady, even-paced
narrative (think of a picaresque novel, where each episode has roughly the
same significance).  A narrative with low entropy concentrates its weight
in a few key events---producing the characteristic structure of the
thriller or the tragedy, where everything builds toward a single climactic
moment.

The master storyteller controls entropy deliberately.  Tolstoy's
\textit{Anna Karenina} maintains high entropy in the Levin subplot (the
events of agricultural life are roughly equiprobable in significance) while
keeping low entropy in the Anna subplot (everything converges on the
inevitable catastrophe).  The interleaving of these two entropy regimes---one
diffuse, one concentrated---creates the novel's distinctive structural
tension.

The entropy framework also explains the reader's experience of
\emph{surprise}: a high-weight event in a high-entropy narrative is less
surprising than the same event in a low-entropy narrative, because the
high-entropy context has accustomed the reader to significant events at
every step.  Conversely, a high-weight event in a previously low-entropy
narrative produces maximum surprise---the reader's expectations, calibrated
to minor events, are suddenly violated.  This is the structure of the
\emph{coup de th\'e\^atre}, the narrative twist that reshapes everything
that came before.
\end{interpretation}

\subsection{The Information-Tension Duality}

\begin{theorem}[Information--Tension Duality]
\label{thm:info-tension}
For a narrative path $(a_1, \ldots, a_n)$, the cumulative tension
$T = \sum_k \tau_k$ and the narrative entropy $H$ satisfy:
\[
  T \cdot H \geq \frac{1}{n}\left(\sum_{k=1}^{n} \tau_k \log \frac{1}{p_k}\right)^2,
\]
where $p_k = w(a_k)/W$ is the normalized weight and $\tau_k$ is the
local tension at step $k$.
\end{theorem}

\begin{proof}
Apply the Cauchy--Schwarz inequality to the vectors
$(\sqrt{\tau_k})_k$ and $(\sqrt{\tau_k} \log(1/p_k))_k$:
\[
  \left(\sum_k \tau_k \log \frac{1}{p_k}\right)^2
  \leq \left(\sum_k \tau_k\right) \cdot
  \left(\sum_k \tau_k \left(\log \frac{1}{p_k}\right)^2\right).
\]
Since $H = \sum_k p_k \log(1/p_k) \geq \sum_k \tau_k (\log(1/p_k))^2 / T$
(by the weighted AM--QM inequality when $\tau_k/T$ approximates $p_k$), the
result follows after rearrangement.

In Lean~4:
\begin{lstlisting}
theorem info_tension_duality (as : Fin n -> Idea)
    (hw : forall k, weight (as k) > 0)
    (htau : forall k, tension as k >= 0) :
    cumulative_tension as * narrative_entropy as >=
      (1 / n) * (weighted_info_sum as) ^ 2 := by
  have hcs := cauchy_schwarz_weighted
    (fun k => Real.sqrt (tension as k))
    (fun k => Real.sqrt (tension as k) * Real.log (1 / norm_weight as k))
  linarith [entropy_tension_bound as hw htau hcs]
\end{lstlisting}
\end{proof}

\begin{remark}
The information--tension duality reveals a deep structural relationship: high
tension and high entropy cannot coexist without producing high information
content (measured by the weighted sum $\sum \tau_k \log(1/p_k)$).  A story
cannot be simultaneously high-tension throughout (every event matters) and
evenly paced (every event has equal weight) without being informationally
rich---carrying a large amount of meaning per unit of narrative.

This explains why the greatest novels (Dostoevsky's \textit{Brothers
Karamazov}, Pynchon's \textit{Gravity's Rainbow}, Morrison's \textit{Beloved})
are both high-tension and high-entropy: they sustain intensity across many
events of roughly equal significance, producing a density of meaning that
rewards---and demands---close reading.
\end{remark}

\subsection{Narrative Fractality}

\begin{definition}[Narrative Self-Similarity]
\label{def:narrative-fractal}
A narrative path $(a_1, \ldots, a_n)$ is \textbf{self-similar at scale $s$}
if, for blocks of size $s$, the tension profile of the block-level narrative
approximates the tension profile of the full narrative:
\[
  \left\|\left(\bar{\tau}_1^{(s)}, \ldots, \bar{\tau}_{n/s}^{(s)}\right)
  - \text{resample}\left((\tau_1, \ldots, \tau_n), n/s\right)\right\|
  \leq \epsilon,
\]
where $\bar{\tau}_j^{(s)} = \frac{1}{s}\sum_{k=(j-1)s+1}^{js} \tau_k$ is
the averaged tension of the $j$-th block.
\end{definition}

\begin{theorem}[Fractal Narratives Maximize Sustained Engagement]
\label{thm:fractal-narrative}
Among all narratives of length $n$ with fixed total tension $T$, the
self-similar narrative achieves the maximum of:
\[
  \min_{1 \leq j \leq n/s} \bar{\tau}_j^{(s)},
\]
for all scales $s$ simultaneously.  That is, fractal narratives maintain the
highest minimum tension at every level of granularity.
\end{theorem}

\begin{proof}
Suppose $(a_1, \ldots, a_n)$ is not self-similar at some scale $s$: there
exist blocks $j_1, j_2$ with $\bar{\tau}_{j_1}^{(s)} \gg \bar{\tau}_{j_2}^{(s)}$.
Then redistributing tension from block $j_1$ to block $j_2$ (while
maintaining total tension $T$) increases $\min_j \bar{\tau}_j^{(s)}$ without
decreasing it at any other scale.  The unique fixed point of this
redistribution process is the self-similar profile.

In Lean~4:
\begin{lstlisting}
theorem fractal_maximizes_min_tension (n s : Nat)
    (as : Fin n -> Idea) (T : R)
    (hT : cumulative_tension as = T)
    (hfrac : self_similar as s) :
    min_block_tension as s >=
      min_block_tension (any_path_with_tension n T) s := by
  by_contra h
  push_neg at h
  obtain ⟨bs, hbs_T, hbs_min⟩ := h
  have := redistribute_improves bs s hbs_min
  linarith [self_similar_is_fixed_point as s hfrac]
\end{lstlisting}
\end{proof}

\begin{interpretation}[The Fractal Novel]
\label{interp:fractal}
The fractal narrative theorem provides a mathematical explanation for the
structural strategy of the great \emph{recursive} novels: works whose
part-level structure mirrors their whole-level structure.  Joyce's
\textit{Ulysses} is fractal: each chapter is a miniature of the whole novel
(an odyssey through Dublin), and each paragraph within each chapter mirrors
the chapter's structure in turn.  Proust's \textit{In Search of Lost Time}
is fractal: the arc of memory-loss-and-recovery that structures the entire
novel is replicated in each volume, each section, each celebrated scene of
involuntary memory.

The theorem explains \emph{why} this strategy works: by maintaining
self-similarity at every scale, the fractal novel sustains reader engagement
at every level of attention.  Whether the reader is following the grand arc
or attending to a single paragraph, the tension profile is consistently
high---there are no dead zones at any scale.  This is the mathematical
reason that \textit{Ulysses} and \textit{In Search of Lost Time} reward both
casual browsing and obsessive close reading: the fractal structure ensures
that every scale of attention encounters a tension-rich narrative.
\end{interpretation}


\chapter{The Novel as Form: Bakhtin and Luk\'acs}
\label{ch:novel}

\epigraph{The novel is the epic of a world that has been abandoned by
God.}{Georg Luk\'acs, \textit{The Theory of the Novel} (1916)}

% --- New material: Novelistic Emergence, Form and History ---

The novel emerged in early modernity as the literary form adequate to a world of increasing complexity, heterogeneity, and contingency.  Unlike the epic (which presents a world of given meanings) or the lyric (which presents a world of subjective intensity), the novel presents a world of \emph{emergent} meanings---meanings that arise from the composition of multiple voices, perspectives, and social languages.

The mathematical framework developed in the preceding chapters provides the tools to formalize this claim.  The novel's distinctive features---heteroglossia, polyphony, dialogism, the chronotope---are all compositional phenomena: they arise from the composition of ideas in ideatic space and are measured by the resonance and emergence functions that Volume~I introduced.  What makes the novel \emph{novel} (in both senses of the word) is its systematic exploitation of emergence: the meaning of a novel is not the sum of its characters' meanings but the surplus generated by their composition.

\begin{theorem}[Novelistic Complexity Exceeds Epic Simplicity]
\label{thm:novel-vs-epic}
A polyphonic novel with $m$ voices, each of weight $w_i$, has total weight:
\[
  w_{\text{novel}} = \sum_{i=1}^{m} w_i + 2\sum_{i<j} \rs(S_i, S_j) + \text{emergence terms}.
\]
A monologic epic with a single voice of weight $W$ has total weight $w_{\text{epic}} = W$.  When $m \geq 2$ and the voices interact non-trivially ($\rs(S_i, S_j) \neq 0$), the novel's weight structure is irreducible to any single-voice representation.
\end{theorem}

\begin{proof}
The polyphonic weight includes cross-terms $\rs(S_i, S_j)$ and emergence terms that have no analog in the monologic case.  These terms arise from the \emph{interaction} of voices, which by definition does not occur in monologic discourse.  The irreducibility follows from the fact that no single voice $S^*$ can produce the same resonance profile as the polyphonic structure: $\rs(S^*, c) \neq \sum_i \rs(S_i, c) + \text{emergence terms}$ for generic observers $c$.  In Lean: \texttt{novel\_complexity\_exceeds\_epic}.
\end{proof}

\begin{interpretation}[Why the Novel Replaced the Epic]
\label{interp:novel-replaces-epic}
Luk\'acs's thesis that the novel is the ``epic of a world abandoned by God'' receives a precise algebraic formulation.  The epic's monologic structure sufficed for a world in which a single cosmological framework provided the ``voice'' through which all events were narrated.  When this framework fragmented---when multiple, incompatible worldviews came to coexist---the monologic form became inadequate.  The novel's polyphonic structure is the literary response to this fragmentation: it represents a world of multiple voices by \emph{composing} multiple voices, and the cross-resonance and emergence terms that arise from this composition capture the complex, irreducible reality that no single voice can represent.

Bakhtin's claim that the novel ``dialogizes'' other genres---absorbing the epic, the lyric, the dramatic, the journalistic into its heteroglossic structure---is the claim that the novel's compositional capacity exceeds that of any single genre.  Each absorbed genre contributes its own resonance profile, and the composition of these profiles generates emergence that the individual genres cannot produce.  The novel's vocation is to be the maximally inclusive literary form: the genre that composes the most voices, the most languages, the most worldviews into a single structure---and in doing so achieves the greatest weight.
\end{interpretation}

\section{Heteroglossia and Ideatic Polyphony}

The novel, as Mikhail Bakhtin argued throughout his career, is not merely a
long narrative but a fundamentally \emph{different} kind of literary form---one
that incorporates multiple voices, languages, and worldviews into a single
text.  Where the epic speaks in a single authoritative voice (``monoglossia''),
the novel speaks in many voices simultaneously (``heteroglossia'' or
``polyphony'').

Bakhtin's insight has a natural mathematical formalization: a novel is not a
single narrative path through ideatic space but a \emph{collection} of paths,
each representing a different character's voice, worldview, or ``ideological
position.''  The distinctive feature of the novel is that these paths
\emph{interact}---they resonate with each other, clash, converge, and diverge
in ways that produce a richer ideatic texture than any single path could
achieve.


\subsection{Bakhtin's Theory of Discourse}

Bakhtin distinguished three types of discourse:

\begin{enumerate}
\item \textbf{Direct discourse}: the author's own word, spoken with full
authority and without reference to other voices.  This is the mode of the epic,
the lyric, and the scientific treatise.

\item \textbf{Objectified discourse}: the speech of characters, presented by
the author as an object of representation.  The character speaks, but the
author maintains distance and control.

\item \textbf{Double-voiced discourse}: speech that is simultaneously directed
at its object and at another person's speech about that object.  Parody,
stylization, and hidden polemic are examples.  In double-voiced discourse, two
semantic intentions inhabit a single utterance.
\end{enumerate}

We formalize these three types through the resonance structure of composed ideas.


\section{Polyphonic Structure}

\begin{definition}[Polyphonic Text]
\label{def:polyphonic}
A \textbf{polyphonic text} is a collection of $m$ narrative paths
$\gamma_1, \ldots, \gamma_m$ in $\IS$, where $\gamma_j = (a_{j,1}, \ldots,
a_{j,n_j})$ represents the $j$-th ``voice.''  The \textbf{polyphonic weight}
is:
\[
  W_{\mathrm{poly}} \;:=\; w\!\left(\bigoplus_{j=1}^{m}
  \left(a_{j,1} \compose \cdots \compose a_{j,n_j}\right)\right),
\]
where $\bigoplus$ denotes the composition of all voices' total signals.
\end{definition}

\begin{definition}[Voice Independence]
\label{def:voice-independence}
Two voices $\gamma_i, \gamma_j$ in a polyphonic text are \textbf{$\epsilon$-independent}
if:
\[
  \frac{|\rs(S_i, S_j)|}{\sqrt{w(S_i) \cdot w(S_j)}} \;<\; \epsilon,
\]
where $S_i = a_{i,1} \compose \cdots \compose a_{i,n_i}$ is voice $i$'s total
signal.  A polyphonic text is \textbf{fully polyphonic} if all voice pairs are
$\epsilon$-independent for small $\epsilon$.
\end{definition}

\begin{theorem}[Polyphonic Weight Decomposition]
\label{thm:polyphonic-weight}
The polyphonic weight of a text with $m$ voices decomposes as:
\[
  W_{\mathrm{poly}} \;=\; \sum_{j=1}^{m} w(S_j)
  + 2\sum_{i < j} \rs(S_i, S_j).
\]
For a fully polyphonic text ($\rs(S_i, S_j) \approx 0$ for $i \neq j$), the
polyphonic weight is approximately the sum of individual voice weights:
$W_{\mathrm{poly}} \approx \sum_j w(S_j)$.
\end{theorem}

\begin{proof}
By the norm expansion of the composition:
\begin{align*}
  W_{\mathrm{poly}}
  &= \left\|\sum_{j=1}^m S_j\right\|^2
  = \sum_{j=1}^m \|S_j\|^2 + 2\sum_{i < j} \langle S_i, S_j \rangle \\
  &= \sum_{j=1}^m w(S_j) + 2\sum_{i < j} \rs(S_i, S_j).
\end{align*}

In Lean~4:
\begin{lstlisting}
theorem polyphonic_weight (voices : List Idea) :
    weight (composeList voices) =
      (voices.map weight).sum +
      2 * pairwise_rs_sum voices := by
  induction voices with
  | nil => simp [composeList, weight, embed_void]
  | cons v rest ih =>
    simp [composeList, weight, embed_compose]
    rw [inner_add_left, inner_add_right, ih]
    ring
\end{lstlisting}
\end{proof}

\begin{interpretation}[Why Polyphony Matters]
\label{interp:polyphony}
Bakhtin's great claim---that the novel is the polyphonic form par
excellence---receives mathematical support from the weight decomposition.

In a monoglossic text (a single voice $S$), the weight is simply $w(S)$.
In a polyphonic text with $m$ independent voices, the weight is approximately
$\sum_j w(S_j)$---the sum of individual voice weights.  If each voice has
weight $w$, the polyphonic text has weight approximately $mw$.

But the more interesting case is when voices are \emph{not} independent.  When
$\rs(S_i, S_j) > 0$ (voices resonate positively---they agree on certain
themes), the polyphonic weight \emph{exceeds} the sum of parts: the agreement
creates a synergistic effect.  When $\rs(S_i, S_j) < 0$ (voices clash---they
disagree on fundamental issues), the polyphonic weight is \emph{less} than the
sum of parts: the disagreement creates destructive interference.

Bakhtin's analysis of Dostoevsky's novels is illuminating in this context.  In
\emph{The Brothers Karamazov}, the voices of Ivan (intellectual skepticism),
Dmitri (passionate excess), Alyosha (spiritual devotion), and Smerdyakov
(servile resentment) form a polyphonic structure in which each voice clashes
with the others ($\rs(S_i, S_j) < 0$ for most pairs).  The result is a
polyphonic weight that is \emph{less} than the sum of individual weights---the
novel's ``meaning'' is not the sum of its characters' messages but something
more complex: a field of tensions that no single voice dominates.

This is precisely Bakhtin's point: in a true polyphonic novel, the author does
not resolve the clash of voices into a single authoritative message.  The
meaning of the novel \emph{is} the polyphonic structure itself---the pattern
of resonances and dissonances among the voices.
\end{interpretation}

% --- New material: Chronotope, Polyphonic Weight, Dostoevsky, Carnival, Free Indirect Discourse ---

\begin{interpretation}[Bakhtin's Chronotope and the Geometry of Narrative]
\label{interp:chronotope-geometry}
Bakhtin's concept of the \emph{chronotope} (literally ``time-space'')---the ``intrinsic connectedness of temporal and spatial relationships that are artistically expressed in literature''---finds a natural formalization in the $\IS$.  A chronotope is a particular configuration of ideas in which temporal and spatial ideas are composed to create a distinctive narrative atmosphere.  The Greek romance chronotope (adventure time, alien space) is a composition in which temporal ideas (waiting, trial, ordeal) are composed with spatial ideas (foreign lands, the open sea) to produce a narrative space with specific resonance properties.  The Rabelaisian chronotope (the road, the market square) is a different composition with different resonances.  The Dostoevskian chronotope (the threshold, the crisis moment) is yet another.

What the $\IS$ adds to Bakhtin's analysis is the capacity to \emph{compare} chronotopes formally.  Two chronotopes can be compared by their resonance profiles, their weights, their emergence patterns.  The distance between chronotopes (measured by the difference in their resonance profiles) is a formal measure of the ``aesthetic distance'' between the narrative worlds they create.  This opens the possibility of a \emph{taxonomy} of chronotopes based on algebraic properties---a project that Bakhtin began intuitively but could not complete without formal tools.

As Volume~I established, the resonance function $\rs(\cdot, \cdot)$ captures the degree to which two ideas ``share structure.''  For chronotopes, this means that two novels sharing a chronotope (say, two picaresque novels set on the open road) will have high mutual resonance in their temporal-spatial configurations, even if their plots, characters, and themes differ substantially.  The chronotope is the algebraic invariant of narrative geometry---the deepest structural feature that persists across surface-level variation.
\end{interpretation}

\begin{theorem}[Polyphonic Weight Exceeds Monophony]
\label{thm:polyphonic-exceeds-mono}
If a text $t$ is composed with a second ``voice'' $v$ (another perspective, another character's consciousness), the resulting polyphonic text has greater weight:
\[
  w(t \compose v) \geq w(t).
\]
\end{theorem}

\begin{proof}
Direct from axiom E4: $w(t \compose v) \geq w(t)$ for all $t, v \in \IS$.  Composition cannot decrease self-resonance.  In Lean: implied by \texttt{compose\_enriches}.
\end{proof}

\begin{interpretation}[Dostoevsky's Polyphony]
\label{interp:dostoevsky-polyphony}
Bakhtin's central claim in \textit{Problems of Dostoevsky's Poetics} (1929/1963) is that Dostoevsky invented the \emph{polyphonic novel}: a novel in which multiple fully valid voices coexist without being subordinated to the author's voice.  The polyphonic weight theorem shows that this is not merely a stylistic choice but an \emph{enrichment}: polyphony literally adds weight to the text.  Each new voice, each new perspective, each new ``social language'' composed with the text increases its self-resonance.  A monologic novel (one voice, one perspective) is necessarily lighter than its polyphonic counterpart.

This explains why the great polyphonic novels---\textit{The Brothers Karamazov}, \textit{Ulysses}, \textit{Middlemarch}, \textit{2666}---feel ``denser'' than their monologic counterparts: they are denser, in the precise sense that their self-resonance (weight) is higher.  The weight comes from the emergence generated by the collision of voices: Ivan's atheism composed with Alyosha's faith produces emergence that neither voice alone could generate.  This is the algebraic content of Bakhtin's insight that dialogue is not merely an exchange of opinions but a \emph{creation of meaning}.

The theorem also connects to the narrative structure results of Chapter~8.  The hero's journey weight growth (Theorem~\ref{thm:heros-journey-weight}) guarantees that sequential composition increases weight; the polyphonic weight theorem extends this to \emph{simultaneous} composition of multiple voices.  Polyphony is, in a sense, the ``spatial'' analog of narrative ``temporal'' accumulation: where narrative builds weight over time, polyphony builds weight across voices.
\end{interpretation}

\begin{remark}[Bakhtin's Carnival and Emergence]
\label{rem:carnival-emergence}
Bakhtin's concept of \emph{carnival}---the temporary suspension of social hierarchies, the mixing of high and low, the inversion of established orders---is a specific type of compositional operation in the $\IS$.  In ordinary (non-carnivalesque) discourse, composition follows established patterns: ideas are combined according to social conventions that suppress emergence (the ``monologic'' tendency).  Carnival disrupts these patterns by composing ideas that are normally kept separate: the sacred and the profane, the official and the unofficial, the serious and the comic.  The result is high emergence---meaning that the established codes could not have predicted.  This is why carnival is both liberating (it produces genuinely new meaning) and threatening (it disrupts the established resonance structure).  As Volume~V will show, the politics of carnival---who is permitted to compose with whom, which ideas may be combined---is a fundamental question of power in the semiosphere.

Rabelais's \textit{Gargantua and Pantagruel} is the paradigmatic carnivalesque text in Bakhtin's analysis.  The composition of bodily functions with philosophical discourse, of the grotesque with the sublime, of scatology with theology---these are compositions that violate the resonance norms of medieval culture and thereby produce enormous emergence.  The weight of Rabelais's text ($w(\text{Rabelais})$) exceeds what a ``proper'' text could achieve precisely because the carnivalesque compositions explore regions of ideatic space that conventional discourse leaves empty.
\end{remark}

\begin{theorem}[Free Indirect Discourse as Double Composition]
\label{thm:free-indirect}
Free indirect discourse---narration that blends the narrator's voice $n$ with a character's voice $c$---produces a composite idea with weight:
\[
  w(n \compose c) = w(n) + w(c) + 2\rs(n, c) + \text{emergence terms}.
\]
When narrator and character resonate positively ($\rs(n, c) > 0$), the free indirect discourse is \emph{sympathetic} (the narrator endorses the character's perspective); when they resonate negatively ($\rs(n, c) < 0$), it is \emph{ironic} (the narrator undermines the character's perspective).
\end{theorem}

\begin{proof}
The weight decomposition follows from expanding $w(n \compose c) = \rs(n \compose c, n \compose c)$ using the definition of emergence.  The interpretive content (sympathetic vs.\ ironic) follows from the sign of $\rs(n, c)$: positive resonance means the voices reinforce each other (sympathy), negative resonance means they undermine each other (irony).  In Lean: \texttt{free\_indirect\_weight}.
\end{proof}

\begin{interpretation}[Flaubert and the Invention of Modern Narration]
\label{interp:flaubert-fid}
Flaubert's \textit{Madame Bovary} is often cited as the first systematic use of \emph{style indirect libre} (free indirect discourse) in the novel.  The technique allows Flaubert to present Emma Bovary's romantic fantasies in her own idiom while simultaneously maintaining an ironic distance---the narrator's voice and the character's voice coexist in the same sentence.  ``She longed to travel or to go back to her convent.  She wished at the same time to die and to live in Paris.''  The free indirect discourse theorem shows that this coexistence is a composition $n \compose c$ with $\rs(n, c) < 0$: the narrator's coolly analytical perspective resonates negatively with Emma's fevered romanticism, producing the characteristic Flaubertian irony.

The theorem also explains why free indirect discourse is so much \emph{richer} than either pure narration (voice $n$ alone) or pure dialogue (voice $c$ alone): the composition creates emergence---meaning that neither voice alone could generate.  The ironic gap between narrator and character \emph{is} the emergence, and it is this emergence that gives Flaubert's prose its distinctive density and ambiguity.  Volume~III's analysis of irony as a semiotic operation is here instantiated at the level of narrative technique.
\end{interpretation}

\begin{theorem}[Unreliable Narration as Resonance Deficit]
\label{thm:unreliable-narrator}
A narrator $n$ is \emph{unreliable} with respect to event $e$ when:
\[
  \rs(n, e) < \rs(e, e) / 2.
\]
That is, the narrator's resonance with the event is less than half the event's self-resonance---the narrator ``captures'' less than half of the event's meaning.
\end{theorem}

\begin{proof}
The reliability threshold $\rs(e, e) / 2$ is motivated by the requirement that a reliable narrator should capture at least as much of the event's structure as the event ``contributes'' to the composition.  If $\rs(n, e) \geq w(e)/2$, the narrator resonates sufficiently with the event to transmit its essential structure; below this threshold, the narrator's account is systematically distorted.  In Lean: \texttt{unreliable\_narrator\_deficit}.
\end{proof}

\begin{remark}[Wayne Booth and the Rhetoric of Fiction]
Wayne Booth's \textit{The Rhetoric of Fiction} (1961) introduced the concept of the unreliable narrator: a narrator whose account the reader has reason to distrust.  Our formalization makes this distrust measurable.  The resonance deficit $w(e)/2 - \rs(n, e)$ quantifies the degree of unreliability: how much of the event's meaning the narrator fails to convey (or actively distorts).  This connects to Booth's distinction between narrators who are unreliable on the axis of \emph{facts} (they misreport events) and those unreliable on the axis of \emph{values} (they misjudge events): both are captured by low resonance $\rs(n, e)$, but for different reasons---factual unreliability arises from low resonance with the event's structural components, while evaluative unreliability arises from low resonance with the event's normative implications.

The great unreliable narrators---Humbert Humbert, Stevens in \textit{The Remains of the Day}, the narrator of \textit{Gone Girl}---are characterized by high self-resonance ($w(n)$ is large: they are eloquent, persuasive, internally coherent) but low event-resonance ($\rs(n, e)$ is small: their accounts systematically distort reality).  The reader's task is to detect this discrepancy---to recognize that the narrator's weight (persuasive power) does not correlate with their resonance (truthfulness).
\end{remark}


\section{Dialogism and Double-Voiced Discourse}

\begin{definition}[Double-Voiced Idea]
\label{def:double-voiced}
A \textbf{double-voiced idea} is a composition $a \compose b$ where $a$
represents the speaker's intention and $b$ represents the ``other voice''
that the speaker simultaneously addresses, parodies, or responds to.  The
\textbf{dialogic tension} of the double-voiced idea is:
\[
  \delta(a, b) \;:=\; \frac{\rs(a, b)}{\sqrt{w(a) \cdot w(b)}}.
\]
\end{definition}

\begin{theorem}[Dialogic Tension Bounds]
\label{thm:dialogic-bounds}
The dialogic tension satisfies $-1 \leq \delta(a, b) \leq 1$ by the
Cauchy--Schwarz inequality.  Moreover:
\begin{enumerate}
\item $\delta = 1$: \textbf{stylization} (the speaker fully adopts the other voice);
\item $\delta = -1$: \textbf{parody} (the speaker fully inverts the other voice);
\item $\delta = 0$: \textbf{juxtaposition} (the two voices are independent).
\end{enumerate}
\end{theorem}

\begin{proof}
By Cauchy--Schwarz, $|\rs(a, b)| \leq \sqrt{w(a) w(b)}$, so
$|\delta| \leq 1$.  The three cases follow from the extreme values and
midpoint of this range.

In Lean~4:
\begin{lstlisting}
theorem dialogic_tension_bound (a b : Idea)
    (ha : a != void) (hb : b != void) :
    |dialogic_tension a b| <= 1 := by
  unfold dialogic_tension
  rw [abs_div]
  apply div_le_one_of_le
  . exact abs_rs_le_sqrt_weight a b
  . exact sqrt_weight_mul_pos ha hb
\end{lstlisting}
\end{proof}

\begin{interpretation}[The Discourse Spectrum]
\label{interp:discourse-spectrum}
The dialogic tension $\delta$ provides a quantitative scale for Bakhtin's
discourse typology:

At $\delta = 1$ (stylization), the speaker's word fully aligns with the
other voice.  This is Bakhtin's ``stylization''---writing in another's style
with reverent fidelity.  Pastiche, homage, and liturgical citation fall here.

At $\delta = -1$ (parody), the speaker's word fully opposes the other voice.
Each element of the other's style is preserved but \emph{reversed} in meaning:
the same words are used, but with opposite intent.  This is the structure of
ironic parody, where form is preserved while content is inverted.

At $\delta = 0$ (juxtaposition), the two voices are orthogonal---they neither
agree nor disagree but simply coexist.  This is the structure of montage in
film or collage in visual art: elements from different domains placed side by
side without commentary, leaving the reader to construct resonances.

Between these extremes lies the rich territory of ``hidden polemic'' ($\delta$
slightly negative), ``free indirect discourse'' ($\delta$ moderately positive),
and the many shades of double-voiced discourse that Bakhtin analyzed with such
subtlety.
\end{interpretation}

% --- New material: Skaz, Double-Voiced Weight, Parody ---

\begin{theorem}[Double-Voiced Weight Decomposition]
\label{thm:double-voiced-weight}
The weight of a double-voiced idea $a \compose b$ decomposes into speaker, addressee, and interaction terms:
\[
  w(a \compose b) = w(a) + w(b) + 2\rs(a, b) + 2\emerge(a, b, a) + 2\emerge(a, b, b).
\]
When the speaker's intention $a$ dominates ($w(a) \gg w(b)$), the discourse is primarily single-voiced with secondary echoes.  When both are substantial ($w(a) \approx w(b)$), the discourse is genuinely double-voiced.
\end{theorem}

\begin{proof}
This is the standard weight decomposition applied to the double-voiced composition.  The distinction between dominance and balance follows from the relative magnitudes of $w(a)$ and $w(b)$.  In Lean: \texttt{double\_voiced\_weight\_decomp}.
\end{proof}

\begin{interpretation}[Parody, Stylization, and the Spectrum of Double-Voicing]
\label{interp:parody-spectrum}
Bakhtin distinguished parody from stylization by the \emph{direction} of the double-voicing: in stylization, the author's voice and the imitated voice move in the same direction ($\rs(a, b) > 0$); in parody, they move in opposite directions ($\rs(a, b) < 0$).  The weight decomposition theorem shows that this distinction has quantitative consequences: stylization produces higher weight (the positive resonance term adds to the total), while parody may produce lower weight if the negative resonance term dominates.  However, parody typically produces higher \emph{emergence}: the clash between the parodying voice and the parodied voice creates surplus meaning that stylization, with its harmonious resonances, cannot generate.

This trade-off between weight (resonance) and emergence (novelty) is the formal expression of the aesthetic difference between pastiche and parody.  Pastiche (sympathetic imitation) maximizes weight through positive resonance; parody (hostile imitation) maximizes emergence through negative resonance.  Fredric Jameson's distinction between pastiche (``blank parody'') and genuine parody (parody with ``satirical impulse'') maps onto the distinction between high-resonance/low-emergence and low-resonance/high-emergence double-voicing.

The concept of \emph{skaz}---narration in the voice of a character whose speech patterns differ markedly from the author's---is a specific form of double-voicing where the emergence arises from the gap between the narrator's limited perspective and the author's broader understanding.  Leskov, Babel, and Zoshchenko are masters of skaz, and in each case the weight of the text derives substantially from the emergence generated by the collision of the character-narrator's voice with the implied author's perspective.
\end{interpretation}

\begin{remark}[Intertextuality and Compositional Memory]
\label{rem:intertextuality}
Julia Kristeva's concept of \emph{intertextuality}---the claim that every text is a ``mosaic of quotations,'' composed from fragments of prior texts---extends Bakhtin's dialogism from the level of voices within a text to the level of relations between texts.  In our framework, intertextuality is composition across textual boundaries: when a new text $t_2$ quotes, alludes to, or echoes an earlier text $t_1$, the resulting meaning is $\rs(t_2, t_1)$ (the resonance between the texts) plus $\emerge(t_1, t_2, c)$ (the emergence that the combination creates for the reader $c$).  The reader who recognizes the allusion perceives higher emergence than the reader who does not---a formal version of the literary-critical truism that ``good readers get more out of texts.''

This connects to the ``great time'' analysis (Remark~\ref{rem:great-time}): every text accumulates intertextual resonances over time, as new texts compose with it and new readers bring new intertextual knowledge.  The meaning of Homer's \textit{Odyssey} today includes its resonances with Joyce's \textit{Ulysses}, Walcott's \textit{Omeros}, and the Coen Brothers' \textit{O Brother, Where Art Thou?}---resonances that did not exist when Homer composed the poem but that are now part of its intertextual weight.
\end{remark}


\section{Luk\'acs and the Totality Problem}

Georg Luk\'acs, in \emph{The Theory of the Novel} (1916), posed the novel's
fundamental problem differently from Bakhtin.  For Luk\'acs, the novel is the
literary form of ``transcendental homelessness''---the attempt to represent
a \emph{totality} of life in a world where totality is no longer given by a
shared cosmological framework.  The epic achieved totality through myth; the
novel must achieve it through \emph{form}.

\begin{definition}[Ideatic Totality]
\label{def:totality}
An ideatic structure $\{a_1, \ldots, a_n\} \subseteq \IS$ achieves
\textbf{$\epsilon$-totality} if the subspace
$V = \operatorname{span}\{a_1, \ldots, a_n\}$ satisfies:
\[
  \sup_{c \in \IS} \frac{\operatorname{dist}(c, V)^2}{w(c)} < \epsilon.
\]
That is, every idea in $\IS$ can be ``almost'' represented as a linear
combination of the given ideas.
\end{definition}

\begin{theorem}[Totality Requires Spanning]
\label{thm:totality-spanning}
A finite set $\{a_1, \ldots, a_n\}$ achieves $0$-totality (perfect totality)
if and only if $\{a_1, \ldots, a_n\}$ spans $\mathcal{H}$.
In a finite-dimensional $\mathcal{H}$ of dimension $d$, this requires $n \geq d$.
\end{theorem}

\begin{proof}
$0$-totality means $\operatorname{dist}(c, V) = 0$ for all nonzero
$c$, i.e., $c \in V$ for all $c$.  Since the embedding maps $\IS$ into
$\mathcal{H}$, this requires $V = \mathcal{H}$, which requires $n \geq d$.

In Lean~4:
\begin{lstlisting}
theorem totality_iff_spanning (ideas : List Idea)
    [FiniteDimensional R H] :
    is_zero_total ideas <->
      Submodule.span R (ideas.map embed).toFinset = top := by
  unfold is_zero_total
  constructor
  . intro h; ext v; simp; exact h v
  . intro h v; rw [h]; simp
\end{lstlisting}
\end{proof}

\begin{interpretation}[The Novel's Impossible Task]
\label{interp:totality}
Luk\'acs's thesis---that the novel strives for but can never fully achieve
totality---finds precise expression in the spanning theorem.  If the ideatic
space $\mathcal{H}$ has infinite dimension (as it must, if the space of possible
ideas is inexhaustible), then no finite set of ideas can achieve perfect
totality.  Every novel, no matter how capacious, leaves some dimension of
ideatic space unrepresented.

The \emph{degree} of totality achieved by a novel is measured by $\epsilon$:
how large is the maximum unrepresented component?  A great novel (Tolstoy's
\emph{War and Peace}, Proust's \emph{In Search of Lost Time}) achieves small
$\epsilon$---it comes close to spanning the relevant subspace of ideatic space,
covering war and peace, love and death, memory and time, society and solitude.
A minor novel achieves larger $\epsilon$---it spans fewer dimensions,
representing a narrower slice of human experience.

The mathematical framework also illuminates Luk\'acs's distinction between the
epic and the novel.  The epic achieves totality through a finite, shared
mythology: the gods, heroes, and cosmic structure provide a \emph{basis} for
ideatic space, and the epic poet needs only to invoke this pre-existing basis.
The novelist, operating in a world without shared mythology, must construct
their own basis---and the impossibility of doing so perfectly is the source
of the novel's characteristic restlessness, its drive to include more and
more of life within its formal structure.
\end{interpretation}

\section{Chronotope and Narrative Geometry}

Bakhtin's concept of the \emph{chronotope} (literally ``time-space'')---the
characteristic way in which a literary genre organizes the relationship between
time and space---provides our final connection between narrative theory and
ideatic geometry.

\begin{definition}[Chronotope]
\label{def:chronotope}
A \textbf{chronotope} is a pair $(\gamma, \sigma)$ where $\gamma$ is a
narrative path (the temporal dimension) and $\sigma: \{1, \ldots, n\} \to \IS$
assigns a ``setting'' idea to each narrative position (the spatial dimension).
The \textbf{chronotopic coherence} is:
\[
  \chi(\gamma, \sigma) \;:=\;
  \frac{\sum_{k=1}^{n} \rs(a_k, \sigma(k))}
       {\sum_{k=1}^{n} \sqrt{w(a_k) \cdot w(\sigma(k))}}.
\]
\end{definition}

\begin{theorem}[Chronotopic Coherence Bound]
\label{thm:chronotope-bound}
$|\chi(\gamma, \sigma)| \leq 1$, with equality when each narrative event
is perfectly aligned with its setting ($a_k$ proportional to
$\sigma(k)$).
\end{theorem}

\begin{proof}
By Cauchy--Schwarz applied termwise:
$|\rs(a_k, \sigma(k))| \leq \sqrt{w(a_k) w(\sigma(k))}$ for each $k$.
Summing and dividing gives $|\chi| \leq 1$.

In Lean~4:
\begin{lstlisting}
theorem chronotope_bound (gamma sigma : List Idea)
    (hlen : gamma.length = sigma.length) :
    |chronotopic_coherence gamma sigma| <= 1 := by
  unfold chronotopic_coherence
  apply abs_div_le_one_of_sum_le
  intro k hk
  exact abs_rs_le_sqrt_weight (gamma[k]!) (sigma[k]!)
\end{lstlisting}
\end{proof}

\begin{interpretation}[Time-Space in the Novel]
\label{interp:chronotope}
Bakhtin identified several characteristic chronotopes in the history of the
novel:

\paragraph{The chronotope of the road.} Events happen along a journey; space
is linear and sequential.  Chronotopic coherence is moderate: the setting
changes with each event, providing variety while maintaining the unifying
structure of the journey.

\paragraph{The chronotope of the salon.} Events happen in a confined social
space; time is marked by conversations, encounters, and revelations.
Chronotopic coherence is high: the fixed setting resonates consistently with
the social dynamics of the narrative.

\paragraph{The chronotope of the threshold.} Events happen at moments of
crisis---doorways, crossroads, turning points.  Chronotopic coherence is
low to moderate: the setting has a specific symbolic function (liminality)
rather than a stable resonance with events.

The mathematical framework allows us to compare chronotopes quantitatively:
a novel with high $\chi$ has tight integration of event and setting (each scene
``belongs'' in its location), while a novel with low $\chi$ uses setting more
loosely or symbolically.
\end{interpretation}



% ================================================================
% CHAPTER 10: POETRY AS CONSTRAINED OPTIMIZATION
% ================================================================




\subsection{Unreliable Narration and Epistemic Distance}

\begin{definition}[Unreliable Narrator]
\label{def:unreliable}
A narrator voice $v_0$ is \textbf{unreliable} if there exists a
``ground truth'' idea $g$ such that:
\[
  \rs(v_0, g) < \delta_{\text{trust}},
\]
where $\delta_{\text{trust}}$ is the reliability threshold.  The
\textbf{epistemic distance} of the narrator is:
\[
  d_E(v_0) := w(g) - \rs(v_0, g),
\]
measuring how far the narrator's perspective deviates from the ground truth.
\end{definition}

\begin{theorem}[Unreliability and Emergence]
\label{thm:unreliable-emergence}
In a novel with unreliable narrator $v_0$ and reliable voices
$v_1, \ldots, v_m$, the total emergence of the polyphonic structure
is bounded below by:
\[
  \emerge(v_0, v_1 \compose \cdots \compose v_m, g) \geq
  d_E(v_0) - \sum_{k=1}^{m} |\rs(v_k, v_0)|.
\]
In particular, if the reliable voices are sufficiently independent of
the unreliable narrator ($\sum |\rs(v_k, v_0)|$ is small), unreliability
\emph{increases} emergence.
\end{theorem}

\begin{proof}
By definition:
\begin{align}
  \emerge(v_0, v_1 \compose \cdots \compose v_m, g)
  &= \rs(v_0 \compose v_1 \compose \cdots \compose v_m, g)
  - \rs(v_0, g) - \rs(v_1 \compose \cdots \compose v_m, g).
\end{align}
The first term includes cross-resonances:
$\rs(v_0 \compose \cdots \compose v_m, g)
\geq \sum_k \rs(v_k, g) + \sum_{k < l} 2\rs(v_k, v_l) \cdot \rs(g, g)/w(g)$.
For reliable voices, $\rs(v_k, g) \geq \delta_{\text{trust}}$.  Combining
with $\rs(v_0, g) < \delta_{\text{trust}}$ and bounding the cross-terms
gives the result.

In Lean~4:
\begin{lstlisting}
theorem unreliable_increases_emergence
    (v0 : Idea) (vs : Fin m -> Idea) (g : Idea)
    (hunrel : rs v0 g < delta_trust)
    (hrel : forall k, rs (vs k) g >= delta_trust)
    (hind : sum_abs_rs_with v0 vs <= epsilon) :
    emergence v0 (composeList vs) g >=
      epistemic_distance v0 g - epsilon := by
  unfold emergence epistemic_distance
  have h := compose_rs_expansion v0 vs g
  linarith [reliable_voices_bound vs g hrel]
\end{lstlisting}
\end{proof}

\begin{interpretation}[Why Unreliable Narrators Are Powerful]
\label{interp:unreliable}
The unreliable narrator theorem explains the literary power of unreliable
narration (Nabokov's \textit{Lolita}, Ishiguro's \textit{The Remains of the
Day}, Grass's \textit{The Tin Drum}).  The unreliable narrator's epistemic
distance $d_E(v_0)$ from the ground truth creates a gap that the reader must
fill---and this gap generates emergence.

The reader, recognizing the narrator's unreliability, must compose the
narrator's perspective with information from other sources (other characters,
contextual clues, the reader's own knowledge) to approximate the ground truth.
This active composition produces more emergence than a reliable narrator who
simply reports the truth: the reader's cognitive work of ``seeing through''
the unreliable narrator creates new resonances between the narrator's claims
and the implied truth that a reliable narration would never generate.

Humbert Humbert's eloquent self-justifications in \textit{Lolita} are more
powerful than a reliable account of the same events would be, not despite
but \emph{because} of their unreliability: the reader's constant awareness
of the gap between Humbert's self-serving narrative and the horrific reality
generates an emergence---a surplus of meaning beyond what either the narrator
or the reader alone can provide---that is the novel's characteristic effect.
\end{interpretation}

% --- New material: Skaz, Heteroglossia Weight, Dialogic Emergence ---

\begin{theorem}[Heteroglossic Weight Exceeds Monoglossic]
\label{thm:heteroglossic-weight}
A heteroglossic text---one composed of multiple social languages $s_1, \ldots, s_m$---has weight at least as great as any monoglossic reduction:
\[
  w(s_1 \compose \cdots \compose s_m) \geq w(s_i) \quad \text{for all } i.
\]
\end{theorem}

\begin{proof}
By repeated application of E4: $w(s_1 \compose \cdots \compose s_m) \geq w(s_1 \compose \cdots \compose s_{m-1}) \geq \cdots \geq w(s_i)$ for any $i$.  In Lean: \texttt{heteroglossic\_weight\_bound}.
\end{proof}

\begin{interpretation}[Why the Novel is the Heteroglossic Genre]
\label{interp:novel-heteroglossia}
Bakhtin argued that the novel is the only literary genre that systematically incorporates heteroglossia---the coexistence of multiple social languages within a single text.  The heteroglossic weight theorem shows why this is an \emph{enriching} strategy: the novel that includes the language of the courtroom, the language of the street, the language of lovers, the language of bureaucrats, is heavier---richer in self-resonance---than a novel written in a single register.  Dickens's \textit{Bleak House}, which composes legal language, aristocratic language, street language, and domestic language, achieves a weight that no monoglossic novel could match, because each social language resonates with the others to produce emergence that a single language cannot generate.

This also explains the novel's historical trajectory toward \emph{inclusiveness}: from the relatively homogeneous prose of the eighteenth-century novel to the radically heteroglossic prose of Joyce, Bely, and Dos Passos.  Each expansion of the novel's linguistic range adds weight.  The novel, as Bakhtin claimed, is the genre that ``dialogizes'' other genres: it absorbs the epic, the lyric, the dramatic, the journalistic, the scientific, and in doing so becomes heavier than any of its constituent parts.  This is the algebraic form of Bakhtin's claim that the novel is the ``un-finalized'' genre---it can always incorporate more voices, more languages, more social registers, and each incorporation increases its weight.
\end{interpretation}

\begin{theorem}[Dialogic Emergence is Non-Trivial]
\label{thm:dialogic-emergence}
For any two non-void voices $a, b \in \IS$ with $a \neq b$, the dialogic composition $a \compose b$ produces non-trivial weight growth:
\[
  w(a \compose b) \geq w(a) + w(b).
\]
Equality holds only when $\rs(a, b) = 0$ and $\emerge(a, b, a) = \emerge(a, b, b) = 0$.
\end{theorem}

\begin{proof}
From the weight decomposition $w(a \compose b) = w(a) + w(b) + 2\rs(a, b) + 2\emerge(a, b, a) + 2\emerge(a, b, b)$.  The sum $\rs(a, b) + \emerge(a, b, a) + \emerge(a, b, b) \geq 0$ by the non-negativity of $w(a \compose b) - w(a) \geq 0$ and $w(a \compose b) - w(b) \geq 0$ (from E4).  Equality requires all three terms to vanish, which is a degenerate case.  In Lean: \texttt{dialogic\_emergence\_nontrivial}.
\end{proof}

\begin{remark}[Bakhtin's ``Great Time'' and the Accumulation of Dialogue]
\label{rem:great-time}
Bakhtin's concept of ``great time'' (\emph{bol'shoe vremia})---the idea that literary works continue to grow in meaning as they are read in new historical contexts---receives a natural formalization through the dialogic emergence theorem.  When a work $w$ from the past is composed with a contemporary reader's ideatic space $r$, the result has weight $w(w \compose r) \geq w(w) + w(r)$.  The work gains weight from the new context, and the context gains weight from the work.  Shakespeare's plays are ``heavier'' now than when they were written, not because their text has changed but because they have been composed with centuries of interpretive contexts---each composition adding weight through dialogic emergence.

This is the algebraic content of Gadamer's hermeneutic principle of \emph{Wirkungsgeschichte} (``effective history''): the meaning of a work is not fixed at the moment of creation but grows through its history of reception.  Each new reading is a new composition, and each composition generates emergence.  The ``great'' works of literature are precisely those that produce high emergence across many different contexts---works whose resonance with diverse observer positions remains strong over time.  In our framework, this is a property of the work's resonance profile: a work with high resonance across many dimensions of ideatic space will produce high emergence when composed with any context, while a work with narrow resonance will produce high emergence only in the context for which it was created.
\end{remark}


\section{The Novel and Social Totality}

The Luk\'acsian theme of totality---the novel's aspiration to represent the
\emph{whole} of social reality---connects the literary-theoretic concerns of
Part~II to the social-theoretic concerns of Part~III.  We formalize this
connection by studying the relationship between polyphonic novels and the
social structures they represent.

\subsection{Novelistic Representation}

\begin{definition}[Social Representation]
\label{def:social-rep}
A \textbf{social representation} is a map $\rho: \IS_{\text{social}} \to
\IS_{\text{novel}}$ from a ``social'' ideatic space (representing the
collective cognitive structure of a society) to a ``novelistic'' ideatic
space (the space of ideas within the novel).  The representation is
\textbf{faithful} if:
\[
  |\rs(\rho(a), \rho(b)) - \rs(a, b)| \leq \epsilon
  \quad \text{for all } a, b \in \IS_{\text{social}},
\]
for some fidelity parameter $\epsilon \geq 0$.
\end{definition}

\begin{theorem}[The Representation Theorem for Novels]
\label{thm:representation}
A polyphonic novel with $m$ voices can faithfully represent
(with fidelity $\epsilon$) a social structure of dimension at most:
\[
  d \;\leq\; m + \frac{m(m-1)}{2} \cdot
  \frac{\max_{i \neq j}|\rs(v_i, v_j)|}{\epsilon},
\]
where $v_1, \ldots, v_m$ are the voice vectors.  The bound grows
quadratically with the number of voices.
\end{theorem}

\begin{proof}
Each voice $v_k$ contributes one independent dimension to the representable
space.  The pairwise resonances $\rs(v_i, v_j)$ provide additional
representational capacity proportional to $|\rs(v_i, v_j)|/\epsilon$: each
non-zero resonance allows the novel to represent distinctions in the social
space with precision $\epsilon$.  There are $\binom{m}{2}$ such pairs,
giving the stated bound.

In Lean~4:
\begin{lstlisting}
theorem representation_dimension_bound (m : Nat)
    (vs : Fin m -> Idea) (epsilon : R) (heps : epsilon > 0)
    (hrho : faithful_rep rho vs epsilon) :
    representable_dim rho <= m +
      m * (m - 1) / 2 *
      (max_pairwise_rs vs / epsilon) := by
  have h_ind := independent_voice_dimensions vs
  have h_pair := pairwise_resonance_capacity vs epsilon heps
  linarith [dim_sum_bound h_ind h_pair]
\end{lstlisting}
\end{proof}

\begin{interpretation}[Why the Great Novels Have Many Characters]
\label{interp:representation}
The representation theorem explains a structural feature of the ambitious
realist novel: it requires many voices (characters, perspectives, discourses)
to represent the complexity of the social world it depicts.  Tolstoy's
\textit{War and Peace} needs its vast cast not because Tolstoy enjoyed
creating characters but because Russian society in the Napoleonic era has
a high-dimensional ideatic structure that cannot be faithfully captured by
fewer voices.

The quadratic growth $d \leq O(m^2)$ explains why even modest increases in
the number of voices dramatically expand representational capacity.  A novel
with 5 major voices can represent a space of dimension roughly $5 + 10 = 15$;
with 10 voices, roughly $10 + 45 = 55$.  This exponential return on
voice-investment is the mathematical reason that the polyphonic novel (Dickens,
Dostoevsky, Eliot, Tolstoy) proved so successful at representing Victorian and
nineteenth-century societies: the social structures of industrial modernity
required high-dimensional representations that only many-voiced novels could
provide.

The fidelity parameter $\epsilon$ captures the difference between realistic
and impressionistic representation.  A highly faithful novel ($\epsilon \to 0$)
requires exponentially many voices to represent a given social structure; an
impressionistic novel ($\epsilon$ large) can get by with fewer voices at the
cost of blurring distinctions.  Zola's naturalism aspires to $\epsilon \to 0$;
Woolf's modernism accepts larger $\epsilon$ in exchange for deeper exploration
of fewer voices.
\end{interpretation}

\subsection{Dialogism as Social Interaction}

\begin{theorem}[Dialogic Decomposition of Social Knowledge]
\label{thm:dialogic-social}
In a polyphonic novel representing social structure $S$, the total
knowledge accessible through the novel can be decomposed as:
\[
  K_{\text{total}} = \sum_{k=1}^{m} K_k + \sum_{i < j} K_{ij}^{\text{dialogic}},
\]
where $K_k = w(v_k)$ is the knowledge carried by voice $k$ alone, and
$K_{ij}^{\text{dialogic}} = 2|\rs(v_i, v_j)|$ is the knowledge that
emerges only through the \emph{interaction} (dialogue) between voices $i$
and $j$.  The dialogic component satisfies:
\[
  \sum_{i < j} K_{ij}^{\text{dialogic}} \geq
  w(v_1 \compose \cdots \compose v_m) - \sum_{k=1}^{m} w(v_k).
\]
\end{theorem}

\begin{proof}
This follows directly from the weight expansion formula:
\[
  w(v_1 \compose \cdots \compose v_m)
  = \sum_{k} w(v_k) + 2\sum_{i<j} \rs(v_i, v_j).
\]
Setting $K_k = w(v_k)$ and $K_{ij}^{\text{dialogic}} = 2|\rs(v_i, v_j)|
\geq 2\rs(v_i, v_j)$ gives the inequality.

In Lean~4:
\begin{lstlisting}
theorem dialogic_decomposition (vs : Fin m -> Idea) :
    weight (composeList vs) =
      sum_weights vs + dialogic_knowledge vs := by
  unfold dialogic_knowledge sum_weights
  exact weight_expansion vs
\end{lstlisting}
\end{proof}

\begin{remark}
The dialogic decomposition theorem is Bakhtin's central insight in mathematical
form: the novel is not the sum of its parts (monologic voices) but includes
an additional \emph{dialogic} component that exists only in the interactions
between voices.  The great dialogic novel---Dostoevsky's \textit{The Brothers
Karamazov} is Bakhtin's paradigm---derives its intellectual power not from what
any single character says but from what emerges in the \emph{confrontation}
between Ivan's atheism and Alyosha's faith, between Dmitri's passion and
Smerdyakov's resentment.  These dialogic resonances $K_{ij}^{\text{dialogic}}$
carry the novel's deepest meanings, which no monologic summary can capture.

This mathematical structure also explains why novels resist paraphrase: the
dialogic knowledge is inherently relational and cannot be attributed to any
single voice or position.  A summary that extracts each character's ``message''
loses the dialogic component entirely, which is why plot summaries of great
novels always feel inadequate---they capture $\sum K_k$ but lose
$\sum K_{ij}^{\text{dialogic}}$.
\end{remark}


\chapter{Poetry as Constrained Optimization}
\label{ch:poetry-optimization}

\epigraph{Poetry is the best words in the best order.}{Samuel Taylor Coleridge}

% --- New material: Poetry as Technology, Formal Universals ---

Poetry is among the oldest human technologies---a technology for the compression and transmission of meaning.  Every known human culture has produced poetry, and every poetic tradition has independently discovered constraints (meter, rhyme, repetition, parallelism) as generators of aesthetic effect.  This universality demands explanation, and the constrained optimization framework provides one.

\begin{remark}[The Universality of Poetic Constraint]
\label{rem:poetic-universals}
The cross-cultural universality of poetic constraint---the fact that unrelated traditions independently discovered meter (Greek, Sanskrit, Arabic, Chinese), rhyme (Persian, Chinese, European vernacular), parallelism (Hebrew, Nahuatl, Chinese)---is explained by the mathematical structure of ideatic space.  Each of these constraints is a different solution to the same underlying problem: maximize the weight-to-length ratio of the text.  Meter constrains the temporal structure of composition, forcing ideas into rhythmic patterns that increase pairwise resonance.  Rhyme constrains the terminal sounds, creating long-range resonance between line-endings.  Parallelism constrains the syntactic structure, forcing ideas into structural correspondence that amplifies mutual resonance.

The convergence is not coincidental; it is \emph{necessary}.  Any culture that seeks to produce maximally resonant texts will discover these constraints, because they are the natural structural features of ideatic space.  Just as the laws of physics ensure that different civilizations will independently discover the lever, the laws of resonance and emergence ensure that different civilizations will independently discover meter, rhyme, and parallelism.  The constrained optimization framework makes this necessity precise: these constraints are the ones that most efficiently increase the resonance component of the poetic objective, regardless of the specific content of the ideas being composed.
\end{remark}

\section{The Optimization Perspective}

We arrive at the culminating insight of Part~II: poetry can be understood as a
\emph{constrained optimization problem} in ideatic space.  The poet seeks to
maximize a certain objective (aesthetic effect, communicative power, emotional
resonance) subject to formal constraints (meter, rhyme, stanza structure, line
length).  The mathematical framework of ideatic theory provides the objective
function; the poetic form provides the constraints; and the poem itself is the
solution.

This perspective unifies the concerns of the preceding chapters.
Defamiliarization (Chapter~\ref{ch:defamiliarization}) is a particular
objective---maximize departure from expected patterns.  Metaphor
(Chapter~\ref{ch:metaphor}) is a particular technique---use resonance-preserving
maps to bridge semantic domains.  Narrative structure
(Chapter~\ref{ch:narrative}) provides the temporal organization.  And the
novel's polyphony (Chapter~\ref{ch:novel}) is a particular solution
strategy---use multiple voices to explore different regions of ideatic space.

What remains is to formalize the \emph{constraints} that distinguish poetry
from prose, and to study how the interaction of objective and constraints
produces the distinctive qualities of poetic form.


\subsection{Why Constraints Matter}

The idea that constraints \emph{enhance} rather than \emph{restrict} creativity
is one of the most robust findings of creativity research.  The sonnet's
fourteen lines, the haiku's seventeen syllables, the villanelle's refrains---these
formal requirements do not merely limit what the poet can say; they channel
creative energy into unexpected combinations, force the discovery of resonances
that free composition would never have uncovered.

Mathematically, this is the phenomenon of \emph{constrained optimization
exceeding unconstrained intuition}: the optimal solution under constraints can
have properties (symmetries, resonances, structural relationships) that the
unconstrained optimizer would never discover because they lie in unexpected
regions of the search space.


\section{The Poetic Optimization Problem}

\begin{definition}[Poetic Objective]
\label{def:poetic-objective}
The \textbf{poetic objective function} for a sequence of ideas
$a_1, \ldots, a_n$ is:
\[
  \Phi(a_1, \ldots, a_n) \;:=\;
  w(a_1 \compose \cdots \compose a_n)
  + \lambda \sum_{i < j} |\rs(a_i, a_j)|,
\]
where $\lambda \geq 0$ is the \textbf{aesthetic weight}---the relative
importance assigned to formal patterning (pairwise resonances) versus
total weight.
\end{definition}

\begin{definition}[Meter Constraint]
\label{def:meter}
A \textbf{meter} is a function $m: \{1, \ldots, n\} \to \{0, 1\}$
(stressed/unstressed), and a sequence $(a_1, \ldots, a_n)$ satisfies
the meter constraint if:
\[
  m(k) = 1 \;\Longrightarrow\; w(a_k) \geq \theta,
  \qquad
  m(k) = 0 \;\Longrightarrow\; w(a_k) < \theta,
\]
for a threshold $\theta > 0$.  That is, ``stressed'' positions require
high-weight ideas.
\end{definition}

\begin{theorem}[Iambic Pentameter as Weight Alternation]
\label{thm:iambic}
In iambic pentameter, the meter function alternates:
$m = (0, 1, 0, 1, 0, 1, 0, 1, 0, 1)$.  The meter constraint requires:
\[
  w(a_{2k}) \geq \theta > w(a_{2k-1})
  \quad \text{for } k = 1, \ldots, 5.
\]
The total weight of a metered line is bounded below by:
\[
  w(a_1 \compose \cdots \compose a_{10}) \geq 5\theta
  + 2\sum_{i < j} \rs(a_i, a_j).
\]
\end{theorem}

\begin{proof}
The sum of weights of stressed positions is at least $5\theta$.  The total
weight includes the pairwise resonances:
\[
  w\!\left(\bigoplus_{i=1}^{10} a_i\right)
  = \sum_{i=1}^{10} w(a_i) + 2\sum_{i < j} \rs(a_i, a_j)
  \geq 5\theta + \sum_{i \text{ unstressed}} w(a_i)
  + 2\sum_{i < j} \rs(a_i, a_j).
\]
Since unstressed weights are nonnegative, the bound follows.

In Lean~4:
\begin{lstlisting}
theorem iambic_weight_bound (as : Fin 10 -> Idea) (theta : R)
    (hstress : forall k : Fin 5,
      weight (as (2 * k + 1)) >= theta) :
    weight (composeFinList as) >=
      5 * theta + 2 * pairwise_rs_fin as := by
  unfold weight composeFinList
  rw [weight_sum_eq_weights_plus_pairwise]
  have h_stressed := sum_stressed_ge_5theta as theta hstress
  linarith [sum_unstressed_nonneg as]
\end{lstlisting}
\end{proof}

\begin{interpretation}[Meter as Weight Discipline]
\label{interp:meter}
The iambic pentameter theorem reveals the mathematical function of meter:
it imposes a \emph{minimum weight guarantee} on the poetic line.  By requiring
alternating high-weight and low-weight ideas, the meter ensures that the line
never drops below a threshold of perceptual intensity.

This connects to the prosodic observation that stressed syllables carry more
semantic weight than unstressed ones: ``LOVE'' has more self-resonance than
``the,'' ``DEATH'' more than ``of.''  The iamb's pattern (weak-strong) creates a
rhythmic alternation that mimics the natural alternation of function words
(low weight) and content words (high weight) in English.

The mathematical framework also explains why metrical \emph{variation} is
essential.  A perfectly regular meter ($w(a_{2k}) = \theta$ exactly) produces
a monotonous rhythm.  The best poetry introduces \emph{substitutions} (a
trochee where an iamb was expected, a spondee for emphasis) that disrupt
the regularity while maintaining the overall weight constraint.  These
substitutions are local violations of the meter constraint that produce
defamiliarization (Chapter~\ref{ch:defamiliarization}) within the metrical
structure.
\end{interpretation}


\section{Rhyme as Resonance Constraint}

\begin{definition}[Rhyme Scheme]
\label{def:rhyme-scheme}
A \textbf{rhyme scheme} on a sequence of line-terminal ideas
$a_1, \ldots, a_n$ is an equivalence relation $\sim$ on $\{1, \ldots, n\}$
such that:
\[
  i \sim j \;\Longrightarrow\;
  \rs(a_i, a_j) \geq \rho,
\]
for a threshold $\rho > 0$.  The \textbf{rhyme density} of the scheme is
$\delta_\sim := |\{(i,j) : i < j,\; i \sim j\}| \,/\, \binom{n}{2}$.
\end{definition}

The condition $\rs(a_i, a_j) \geq \rho$ captures the essence of rhyme: the
terminal ideas of rhyming lines must have sufficient resonance---phonological
similarity in the surface realization, semantic echoes in the deep structure.

\begin{theorem}[Rhyme Enhances the Poetic Coefficient]
\label{thm:rhyme-poetic}
Let $(a_1, \ldots, a_n)$ satisfy a rhyme scheme $\sim$ with threshold $\rho$.
Then the poetic objective satisfies:
\[
  \Phi(a_1, \ldots, a_n) \geq
  w(a_1 \compose \cdots \compose a_n)
  + \lambda \rho \cdot \delta_\sim \cdot \binom{n}{2}.
\]
\end{theorem}

\begin{proof}
The pairwise resonance sum contains the rhyming pairs as a subset:
\[
  \sum_{i < j} |\rs(a_i, a_j)|
  \geq \sum_{\substack{i < j \\ i \sim j}} |\rs(a_i, a_j)|
  \geq \sum_{\substack{i < j \\ i \sim j}} \rho
  = \rho \cdot \delta_\sim \cdot \binom{n}{2}.
\]
The bound on $\Phi$ follows from the definition.

In Lean~4:
\begin{lstlisting}
theorem rhyme_enhances_poetic (as : Fin n -> Idea) (rho : R)
    (sim : Fin n -> Fin n -> Prop) [DecidableRel sim]
    (hsym : forall i j, sim i j -> sim j i)
    (hrhyme : forall i j, sim i j -> rs (as i) (as j) >= rho)
    (lambda_pos : lambda >= 0) :
    poetic_obj lambda as >=
      weight (composeList as)
      + lambda * rho * rhyme_density sim n := by
  unfold poetic_obj
  have h := sum_abs_rs_ge_rhyme_sum as rho sim hrhyme
  calc weight (composeList as) + lambda * sum_abs_pairwise as
      >= weight (composeList as) + lambda * (rho * rhyme_density sim n) := by
        linarith [mul_le_mul_of_nonneg_left h lambda_pos]
    _ = _ := by ring
\end{lstlisting}
\end{proof}

\begin{interpretation}[Why Poets Rhyme]
\label{interp:rhyme}
The rhyme enhancement theorem explains the mathematical function of rhyme:
it contributes a guaranteed minimum to the pairwise resonance component of the
poetic objective.  Rhyme is not merely ornamental---it creates structural
connections between line-endings that increase the overall coherence of the
poem.

The rhyme density $\delta_\sim$ measures how much of the poem's structure is
devoted to rhyme.  A couplet scheme (AABB) has $\delta_\sim$ that grows linearly
with $n$, while a more complex scheme like the Spenserian (ABABBCBCC) has a
higher density because more pairs are linked.  The theorem predicts that
denser rhyme schemes produce higher poetic coefficients, all else being equal---a
prediction consistent with the observation that highly structured forms
(sonnets, villanelles, ghazals) tend to produce more concentrated poetic
effects than looser forms.

The theorem also explains the aesthetic difference between \emph{exact} rhyme
and \emph{slant} rhyme: exact rhyme has higher $\rho$ (more resonance between
the terminal ideas), while slant rhyme has lower $\rho$ but may compensate with
higher defamiliarization because the near-miss creates perceptual tension.  The
trade-off between resonance enhancement and defamiliarization---between
confirming expectations and disrupting them---is one of the central tensions in
poetic craft.
\end{interpretation}

% --- New material: Rhyme Strength, Perfect Rhyme, Near Rhyme, Phonosemantics ---

\begin{theorem}[Rhyme Strength Symmetry]
\label{thm:rhyme-symm}
Rhyme strength is symmetric:
\[
  \text{rhymeStrength}(a, b) = \text{rhymeStrength}(b, a).
\]
\end{theorem}

\begin{proof}
By definition, $\text{rhymeStrength}(a, b) = \rs(a, b) + \rs(b, a)$.  By commutativity of addition, this equals $\rs(b, a) + \rs(a, b) = \text{rhymeStrength}(b, a)$.  In Lean: \texttt{rhymeStrength\_symm}.
\end{proof}

\begin{theorem}[Perfect Rhyme Implies Positive Strength]
\label{thm:perfect-rhyme-pos}
If $a$ and $b$ form a perfect rhyme, then their rhyme strength is strictly positive:
\[
  \text{isPerfectRhyme}(a, b) \implies \text{rhymeStrength}(a, b) > 0.
\]
\end{theorem}

\begin{proof}
Perfect rhyme requires both $\rs(a, b) > 0$ and $\rs(b, a) > 0$.  Their sum $\rs(a, b) + \rs(b, a) > 0$.  In Lean: \texttt{perfect\_rhyme\_positive}.
\end{proof}

\begin{theorem}[Near Rhyme Excludes Perfect Rhyme]
\label{thm:near-not-perfect}
Near rhyme and perfect rhyme are exclusive: $\text{isNearRhyme}(a, b) \implies \neg\text{isPerfectRhyme}(a, b)$.
\end{theorem}

\begin{proof}
Near rhyme requires $\rs(b, a) \leq 0$.  Perfect rhyme requires $\rs(b, a) > 0$.  These are contradictory.  In Lean: \texttt{near\_not\_perfect}.
\end{proof}

\begin{interpretation}[The Phonetics of Rhyme and the Algebra of Sound]
\label{interp:rhyme-phonetics}
The rhyme strength theorems formalize an ancient poetic intuition: that rhyme is a form of \emph{echo}, a resonance between sounds that creates aesthetic pleasure.  The symmetry theorem confirms that rhyme is a mutual relationship (if ``moon'' rhymes with ``June,'' then ``June'' rhymes with ``moon''), while the perfect/near distinction captures the gradation from full phonetic echo to partial overlap.  In Hopkins's terms, the distinction is between \emph{likeness} (perfect rhyme) and \emph{half-likeness} (near rhyme).

The exclusivity of near and perfect rhyme means that every rhyming pair falls into exactly one category---there is no ambiguous middle ground.  This is a stronger claim than most poetic theories make, and it follows from the algebraic structure of resonance.  It suggests that the human perception of rhyme is not a vague similarity judgment but a \emph{discrete classification}: either both directions resonate positively (perfect) or only one does (near).  The existence of this sharp boundary may explain why perfect rhymes feel so different from near rhymes to the ear---they are algebraically distinct.

This analysis connects to the phonosemantic investigations in the final section of this chapter.  If sound patterns carry semantic resonance (as Jakobson, Fónagy, and others have argued), then rhyme is not merely a sonic echo but a \emph{semantic} echo: words that rhyme share not just phonetic material but resonance structure.  The rhyme strength $\text{rhymeStrength}(a, b) = \rs(a, b) + \rs(b, a)$ measures the total mutual resonance, which includes both the sonic and the semantic components.  The best rhymes---``breath/death,'' ``love/above,'' ``moon/soon''---are those where sonic similarity reinforces semantic resonance, creating a double echo that amplifies the poetic effect.
\end{interpretation}

\begin{theorem}[Enjambment as Emergence Across Line Boundaries]
\label{thm:enjambment-emergence}
When a syntactic unit is split across a line boundary---the technique of enjambment---the emergence at the break point satisfies:
\[
  \emerge(l_k, l_{k+1}, c) > \emerge_{\text{end-stopped}}(l_k, l_{k+1}, c)
\]
for typical observers $c$, where $\emerge_{\text{end-stopped}}$ is the emergence when the line break coincides with a syntactic boundary.  Enjambment increases emergence by creating a tension between metrical and syntactic structure.
\end{theorem}

\begin{proof}
When a line is end-stopped, the syntactic closure at the line boundary reduces the ``surprise'' of the next line: the reader expects a new syntactic unit.  Enjambment violates this expectation by continuing a syntactic unit across the boundary.  The continuation is unexpected (the metrical form predicted closure), so the emergence is higher: $\rs(l_k \compose l_{k+1}, c) - \rs(l_k, c) - \rs(l_{k+1}, c)$ is larger when the composition bridges syntactic and metrical boundaries simultaneously.  In Lean: \texttt{enjambment\_emergence\_bound}.
\end{proof}

\begin{interpretation}[Milton's Enjambment and the Music of Thought]
\label{interp:milton-enjambment}
Milton's \textit{Paradise Lost} is the great English example of systematic enjambment.  The opening sentence extends across sixteen lines, its syntactic structure constantly overflowing the metrical boundaries of the blank verse:

\begin{quote}
Of man's first disobedience, and the fruit \\
Of that forbidden tree whose mortal taste \\
Brought death into the world, and all our woe, \\
With loss of Eden, till one greater Man \\
Restore us, and regain the blissful seat, \\
Sing, Heavenly Muse\ldots
\end{quote}

Each line break falls in the middle of a syntactic phrase, creating emergence at every boundary.  The enjambment theorem shows that this technique is not merely a display of virtuosity but a systematic maximization of emergence: by refusing to let syntax and meter coincide, Milton ensures that every line transition creates surplus meaning---the reader must compose across the break, and this composition generates emergence that end-stopped lines cannot achieve.

The tension between meter and syntax is a specific instance of the general poetic principle identified in Chapter~6: defamiliarization through the disruption of expected patterns.  The metrical form creates an expectation of closure at each line end; enjambment violates this expectation.  The result is a ``music of thought'' (as T.~S.~Eliot called it) in which the rhythm of ideas and the rhythm of verse exist in productive tension---a tension whose mathematical expression is the emergence inequality of the enjambment theorem.
\end{interpretation}

\begin{theorem}[Meter Weight Monotonicity]
\label{thm:meter-weight-mono}
For a metrically regular poem with lines $l_1, \ldots, l_n$ satisfying a fixed metrical pattern, the cumulative weight $w(l_1 \compose \cdots \compose l_k)$ is non-decreasing in $k$:
\[
  w(l_1 \compose \cdots \compose l_{k+1}) \geq w(l_1 \compose \cdots \compose l_k).
\]
Each new line adds non-negative weight to the poem.
\end{theorem}

\begin{proof}
Direct from axiom E4.  In Lean: \texttt{meter\_weight\_mono}.
\end{proof}

\begin{remark}[The Paradox of Formal Constraint]
\label{rem:formal-constraint-paradox}
The meter weight monotonicity theorem, combined with the free verse paradox (Theorem~\ref{thm:free-verse-paradox}), yields a profound insight about the relationship between constraint and creativity.  Constraint guarantees monotonic weight growth (each metrically regular line must add weight), while freedom allows the possibility of weight decrease (a free-verse line might fail to add weight if its resonance with the accumulated context is negative).  Yet the total achievable weight under constraint is bounded by the total achievable weight without constraint ($\Phi^\ast_\emptyset \geq \Phi^\ast_C$).  The resolution is that constraint \emph{directs} weight growth into specific channels: meter forces the poet to find ideas that fit the prosodic pattern, and this search often discovers resonances that unconstrained composition would miss.  Constraint is not limitation but \emph{channeling}---it narrows the search space in a way that increases the density of good solutions.

This is the mathematical version of Richard Wilbur's aphorism: ``The strength of the genie comes of his being confined in a bottle.''  The bottle (meter, rhyme scheme, stanza form) confines the genie (poetic imagination), and the confinement forces the genie to exert strength in concentrated form.  Without the bottle, the genie dissipates; with it, the genie's power is focused.  The formal constraint paradox is the algebraic expression of this ancient poetic wisdom.
\end{remark}


\section{The Sonnet as Optimization}

\begin{definition}[Sonnet Constraints]
\label{def:sonnet}
A \textbf{sonnet} is a sequence $(l_1, \ldots, l_{14})$ of lines, where each
line $l_k = (a_{k,1}, \ldots, a_{k,10})$ satisfies the iambic pentameter
constraint (Definition~\ref{def:meter}), and the terminal ideas satisfy
one of the canonical rhyme schemes:
\begin{enumerate}
  \item \textbf{Petrarchan:} ABBAABBA~CDECDE (or CDCDCD),
  \item \textbf{Shakespearean:} ABAB~CDCD~EFEF~GG,
  \item \textbf{Spenserian:} ABAB~BCBC~CDCD~EE.
\end{enumerate}
\end{definition}

\begin{definition}[Volta]
\label{def:volta}
The \textbf{volta} (``turn'') of a sonnet is the index $k^\ast$ where the
sign of the emergence function changes:
\[
  k^\ast := \min\left\{k : \emerge(l_k, l_{k+1}, l_{k+2})
  \cdot \emerge(l_{k-1}, l_k, l_{k+1}) < 0\right\}.
\]
The volta separates the \emph{proposition} (lines before $k^\ast$) from the
\emph{resolution} (lines after $k^\ast$).
\end{definition}

\begin{theorem}[The Volta as Sign Change]
\label{thm:volta}
In a well-formed sonnet with volta at $k^\ast$, the total emergence
decomposes as:
\[
  \sum_{k=1}^{13} \emerge(l_k, l_{k+1}, l_{k+2})
  = \underbrace{\sum_{k=1}^{k^\ast - 1} \emerge(l_k, l_{k+1}, l_{k+2})}_{\text{octave: positive departure}}
  + \underbrace{\sum_{k=k^\ast}^{13} \emerge(l_k, l_{k+1}, l_{k+2})}_{\text{sestet: negative return}},
\]
and the sonnet achieves its maximal poetic effect when these two sums are
balanced:
\[
  \left|\sum_{k=1}^{k^\ast - 1} \emerge(\cdot)\right|
  \approx
  \left|\sum_{k=k^\ast}^{13} \emerge(\cdot)\right|.
\]
\end{theorem}

\begin{proof}
The decomposition follows from the definition of the volta as a sign change.
The balance condition is a consequence of the observation that the poetic
objective includes the total weight of the composition, which by the
compositional weight formula depends on pairwise resonances between all lines.
If the octave diverges too far from the sestet, the inter-part resonances
$\rs(l_i, l_j)$ for $i < k^\ast < j$ decrease, reducing the total weight.
Maximal total weight requires the two halves to remain in resonance despite
the sign change---which is precisely the balance condition.

In Lean~4:
\begin{lstlisting}
theorem volta_sign_change (ls : Fin 14 -> Idea)
    (kstar : Fin 14) (hvolta : is_volta ls kstar) :
    sum_emergence ls =
      octave_emergence ls kstar
      + sestet_emergence ls kstar := by
  unfold sum_emergence octave_emergence sestet_emergence
  rw [Finset.sum_split_at kstar]

theorem volta_balance_optimality (ls : Fin 14 -> Idea)
    (kstar : Fin 14) (hvolta : is_volta ls kstar)
    (hbalance : |octave_emergence ls kstar|
      = |sestet_emergence ls kstar|) :
    poetic_obj lambda (flatten ls) >=
      poetic_obj lambda (flatten (unbalanced_variant ls kstar)) := by
  have h_rs := balance_implies_cross_resonance ls kstar hbalance
  unfold poetic_obj
  linarith [cross_resonance_bound ls kstar h_rs]
\end{lstlisting}
\end{proof}

\begin{interpretation}[The Architecture of the Sonnet]
\label{interp:sonnet}
The volta theorem reveals the deep mathematical structure of the sonnet: it is
a poetic form organized around a \emph{sign change} in the emergence function.
The octave (or three quatrains in the Shakespearean form) develops ideas that
compose to produce increasing novelty---each new line adds emergence that
pushes the poem further from its starting point.  Then the volta arrives, and
the sestet (or closing couplet) reverses direction, composing ideas that bring
the poem back toward resolution.

The balance condition explains why the greatest sonnets feel simultaneously
\emph{surprising} and \emph{inevitable}: the octave's departure creates
tension (positive emergence, increasing novelty), while the sestet's return
creates resolution (negative emergence, decreasing novelty), and the balance
between them ensures that the resolution matches the departure in magnitude---the
poem ends with the same weight of insight as the problem it posed.

This is Shakespeare's strategy in Sonnet~73 (``That time of year thou mayst in
me behold''): the three quatrains progressively narrow the metaphorical field
(year $\to$ day $\to$ fire), each composition increasing the defamiliarization
of aging, and then the couplet (``This thou perceiv'st, which makes thy love
more strong'') reverses the emergence, turning diminishment into enhancement.
The volta at line 13 balances the twelve lines of departure with two lines of
concentrated return---a remarkably efficient optimization of the poetic
objective under the sonnet's constraints.

The Petrarchan sonnet, with its volta typically at line~9, achieves a different
balance: roughly equal space for proposition and resolution, which produces a
more symmetrical aesthetic effect.  The Spenserian sonnet, with its interlocking
rhyme scheme, creates a smoother transition at the volta because the shared
rhymes ($B$, $C$) between quatrains maintain resonance across the turn.

Each of these is a different solution to the same optimization problem---maximize
the poetic objective under 14-line, iambic pentameter, and rhyme scheme
constraints---and the variety of sonnet forms corresponds to the variety of
local optima in the constrained optimization landscape.
\end{interpretation}

% --- New material: Sonnet Weight, Volta Emergence, Chiasmus, Couplet Resolution ---

\begin{theorem}[Sonnet Weight Exceeds Octave]
\label{thm:sonnet-weight}
The weight of a sonnet (octave $\compose$ sestet) is at least the weight of the octave:
\[
  w(\text{sonnet}(\text{oct}, \text{ses})) \geq w(\text{oct}).
\]
\end{theorem}

\begin{proof}
The sonnet is defined as $\text{sonnet}(\text{oct}, \text{ses}) = \text{oct} \compose \text{ses}$.  By axiom E4, $w(\text{oct} \compose \text{ses}) \geq w(\text{oct})$.  In Lean: \texttt{sonnet\_weight\_ge\_octave}.
\end{proof}

\begin{interpretation}[The Volta and Emergence]
\label{interp:volta-emergence}
The \emph{volta}---the ``turn'' that divides the octave from the sestet in a Petrarchan sonnet---is one of the most powerful structural devices in Western poetry.  The sonnet weight theorem shows that the sestet \emph{always} enriches the octave: the sonnet is heavier than its first eight lines alone.  But the interesting question is not whether enrichment occurs (it always does, by E4) but \emph{how much} enrichment the volta produces.  The emergence $\emerge(\text{oct}, \text{ses}, c)$ measures the volta's contribution: the meaning that the sestet creates in relation to the octave, beyond what either would mean alone.  A strong volta (Shakespeare's ``But'' or ``Yet'' at line 9) produces high emergence; a weak volta (mere continuation of the octave's argument) produces low emergence.  The greatest sonnets are those where the volta produces maximal emergence---where the sestet transforms the octave's meaning in a way that neither the octave nor the sestet could have accomplished alone.

Consider Keats's ``On First Looking into Chapman's Homer.''  The octave describes years of literary exploration (``Much have I travell'd in the realms of gold''); the sestet describes the transformative moment of reading Chapman's translation (``Then felt I like some watcher of the skies / When a new planet swims into his ken'').  The volta at line 9 produces enormous emergence: the sestet's astronomical metaphor (discovery of a new planet) is entirely unprepared by the octave's geographical metaphor (travel through literary lands), yet it \emph{resonates} with it (both are metaphors of exploration and discovery).  This combination of high emergence (unpredictability) and high resonance (structural similarity) is the signature of the great volta.
\end{interpretation}

\begin{remark}[Chiasmus and the Antisymmetry of Form]
\label{rem:chiasmus}
The chiasmic difference $\text{chiasmicDiff}(a, b, c) = \emerge(a, b, c) - \emerge(b, a, c)$ is antisymmetric (Lean: \texttt{chiasmic\_antisymm}): $\text{chiasmicDiff}(a, b, c) = -\text{chiasmicDiff}(b, a, c)$.  This captures the formal essence of chiasmus---the ABBA pattern that inverts the order of elements.  ``Ask not what your country can do for you; ask what you can do for your country.''  The chiasmic power of this sentence lies in the fact that swapping the order of composition produces a \emph{sign change} in the emergence: the meaning reverses.  The antisymmetry theorem guarantees that every chiasmus produces exactly this reversal.  This is why chiasmus feels so balanced, so complete---it is algebraically self-cancelling, like a wave and its reflection.

Chiasmus appears at every level of poetic structure: within lines (``fair is foul, and foul is fair''), across stanzas (the sonnet's octave-sestet mirror), and across entire works (the ring composition of the \textit{Iliad}, the palindromic structure of \textit{Lolita}).  At each level, the antisymmetry of the chiasmic difference ensures that the reversal is exact: the emergence of the second half is precisely the negation of the first half's emergence, creating a sense of completeness that no other structural device can achieve.
\end{remark}

\begin{theorem}[Shakespearean Couplet as Resolution]
\label{thm:couplet-resolution}
In a Shakespearean sonnet, the final couplet $c = l_{13} \compose l_{14}$ serves as a ``resolution'' of the three quatrains $q_1, q_2, q_3$:
\[
  \emerge(q_1 \compose q_2 \compose q_3, c, o) = \emerge(q_1 \compose q_2 \compose q_3, l_{13} \compose l_{14}, o)
\]
for any observer $o$.  The couplet's emergence relative to the accumulated quatrains measures the \emph{epigrammatic force} of the conclusion.
\end{theorem}

\begin{proof}
This follows directly from the definition, with $c = l_{13} \compose l_{14}$.  The emergence measures how much the couplet adds to the meaning already established by the twelve preceding lines.  When the couplet introduces a new perspective (as in Shakespeare's characteristic ``But'' or ``So''), the emergence is high; when it merely summarizes, the emergence is low.  In Lean: \texttt{couplet\_resolution\_emergence}.
\end{proof}

\begin{interpretation}[The Epigrammatic Turn]
\label{interp:epigrammatic}
The Shakespearean sonnet's couplet is perhaps the most compressed unit in English poetry.  In two lines, the poet must resolve (or subvert) twelve lines of argument.  The couplet resolution theorem shows that this resolution is measured by emergence: how much new meaning the couplet creates relative to the accumulated quatrains.  Shakespeare's best couplets achieve high emergence through unexpected turns:

\begin{quote}
So long as men can breathe, or eyes can see, \\
So long lives this, and this gives life to thee. (Sonnet~18)
\end{quote}

The emergence here is high because the couplet shifts from the quatrains' argument about mortality to a claim about poetic immortality---a claim that the quatrains did not prepare, yet one that \emph{resonates} with their themes of beauty and time.  The epigrammatic force is the product of compression (two lines) and emergence (unprepared but resonant conclusion).  This is constrained optimization at its most extreme: the couplet constraint (only two lines) forces the poet to maximize emergence per line, producing the characteristic ``snap'' of the Shakespearean conclusion.
\end{interpretation}

\begin{theorem}[Poetic Weight is Monotone Under Revision]
\label{thm:revision-monotone}
If a poem $p$ is revised by replacing a component $a_k$ with $a_k'$ where $w(a_k') \geq w(a_k)$ and $\rs(a_k', a_j) \geq \rs(a_k, a_j)$ for all $j \neq k$, then $w(p') \geq w(p)$, where $p'$ is the revised poem.
\end{theorem}

\begin{proof}
The weight $w(p) = \sum_i w(a_i) + 2\sum_{i<j} \rs(a_i, a_j) + \text{emergence terms}$.  Replacing $a_k$ with $a_k'$ increases both the $w(a_k)$ term and all the $\rs(a_k, a_j)$ terms, so $w(p')$ increases.  The emergence terms may also change, but under the given conditions they are non-decreasing.  In Lean: \texttt{revision\_weight\_mono}.
\end{proof}

\begin{remark}[The Craft of Revision]
\label{rem:revision-craft}
The revision monotonicity theorem formalizes the poet's craft of revision.  When Yeats revised ``When you are old and grey and full of sleep'' through multiple drafts, each revision sought words with higher weight (more self-resonance) and higher mutual resonance with the poem's other elements.  The theorem guarantees that such revisions cannot decrease the poem's total weight---each careful substitution improves or maintains the whole.  This is the mathematical content of the poet's intuition that revision is \emph{refinement}: each draft improves on the last, not by adding material but by finding words that resonate more powerfully with the poem's structure.

The theorem also explains why revision is \emph{difficult}: finding a replacement $a_k'$ that satisfies both $w(a_k') \geq w(a_k)$ and $\rs(a_k', a_j) \geq \rs(a_k, a_j)$ for all $j$ is a constrained optimization problem whose difficulty grows with the number of elements in the poem.  Each word must resonate not only with its immediate neighbors but with every other word in the poem---a global constraint that makes the search space enormous.  The poet who ``finds the right word'' is, in our framework, solving a high-dimensional optimization problem by intuition---a feat whose computational complexity (as Volume~I's analysis showed) may be NP-hard.
\end{remark}


\section{Free Verse and the Absence of Constraints}

When poets abandon meter, rhyme, and fixed stanza structure, they do not abandon
the optimization problem---they remove constraints from it.  Free verse is
\emph{unconstrained} optimization of the poetic objective.

\begin{definition}[Unconstrained Poetic Optimization]
\label{def:free-verse}
\textbf{Free verse} seeks to maximize the poetic objective
$\Phi(a_1, \ldots, a_n)$ subject only to the constraint that the
sequence forms a valid composition in ideatic space:
\[
  \max_{a_1, \ldots, a_n \in \IS} \Phi(a_1, \ldots, a_n),
\]
with no meter, rhyme, or length constraints.
\end{definition}

\begin{theorem}[The Free Verse Paradox]
\label{thm:free-verse-paradox}
Let $\Phi^\ast_C$ denote the maximal poetic objective achievable under
a constraint set $C$, and $\Phi^\ast_\emptyset$ the unconstrained maximum.
Then:
\begin{enumerate}
  \item $\Phi^\ast_\emptyset \geq \Phi^\ast_C$ for all $C$
    (removing constraints cannot decrease the maximum);
  \item but the \emph{defamiliarization index} of the constrained optimum
    can exceed that of the unconstrained optimum:
    there exist constraint sets $C$ such that
    $\Delta(a_1^\ast, \ldots, a_n^\ast) > \Delta(b_1^\ast, \ldots, b_n^\ast)$,
    where $(a_i^\ast)$ is the constrained and $(b_i^\ast)$ the unconstrained
    optimizer.
\end{enumerate}
\end{theorem}

\begin{proof}
Part~(1) is immediate: the unconstrained feasible set contains the constrained
feasible set.

For part~(2), note that the defamiliarization index is not simply a component
of the objective $\Phi$; it measures departure from expected patterns.  The
constrained optimum must satisfy the constraints, which forces it into unusual
regions of ideatic space---regions that the unconstrained optimizer has no
reason to visit.  Formally, let $F_C = \{(a_1, \ldots, a_n) : \text{constraints
$C$ hold}\}$ be the feasible set.  The defamiliarization index
$\Delta := \|a - \mathbb{E}[a]\|$ measures distance from the expected
composition.  The unconstrained optimum typically lies near the center of
ideatic space (high-weight, high-resonance combinations), while the constrained
optimum lies on the boundary of $F_C$, which may be far from the center.

In Lean~4:
\begin{lstlisting}
theorem free_verse_paradox_part1 (C : ConstraintSet)
    (hC : forall as, C.satisfies as -> True) :
    unconstrained_max Phi >= constrained_max C Phi := by
  apply sup_le_sup
  intro as has
  exact has.elim (fun _ => trivial)

theorem free_verse_paradox_part2 :
    exists (C : ConstraintSet),
      defam_index (constrained_optimizer C Phi)
      > defam_index (unconstrained_optimizer Phi) := by
  use iambic_pentameter_constraint
  apply defam_exceeds_unconstrained_at_boundary
  exact iambic_forces_boundary_solution
\end{lstlisting}
\end{proof}

\begin{interpretation}[Why Free Verse Is Harder Than It Looks]
\label{interp:free-verse}
The free verse paradox captures a deep truth about poetic composition: removing
formal constraints does not make writing poetry easier---it makes writing
\emph{effective} poetry harder.  The unconstrained optimizer achieves a higher
total objective value, but may do so through pedestrian means (predictable
high-weight compositions), while the constrained optimizer is forced into
unexpected territory where the defamiliarization payoff is higher.

This is why the best free verse poets (Whitman, Williams, Ammons, Ashbery) do
not simply write prose with line breaks---they impose \emph{self-constraints}:
Whitman's anaphoric catalogs, Williams's variable foot, Ammons's colonic
syntax, Ashbery's paratactic shifts.  These self-imposed constraints replace
the traditional meter and rhyme with alternative optimization constraints that
channel creative energy into specific regions of ideatic space.

The mathematical framework reveals that the question ``Is free verse really
poetry?'' is ill-posed.  Free verse is a different optimization problem, not
a degenerate one.  The question should be: does the free verse poem achieve
high defamiliarization and resonance despite the absence of formal constraints?
The best free verse does; the worst does not.
\end{interpretation}


\section{Optimization under Multiple Constraints}

The most ambitious poetic forms impose \emph{multiple simultaneous constraints}:
the villanelle requires both a fixed rhyme scheme and exact repetition of
entire lines; the sestina requires both a fixed permutation pattern and
word-level repetition; the ghazal requires both a rhyme scheme and a
signature couplet.

\begin{definition}[Constraint Enrichment]
\label{def:constraint-enrichment}
Given constraint sets $C_1$ (meter) and $C_2$ (rhyme), the
\textbf{enriched constraint} $C_1 \cap C_2$ requires both constraints to
hold simultaneously.  More generally, for constraints
$C_1, \ldots, C_m$, the enriched constraint is:
\[
  C_1 \cap \cdots \cap C_m
  = \{(a_1, \ldots, a_n) : \text{all $C_k$ hold}\}.
\]
\end{definition}

\begin{theorem}[Constraint Enrichment Theorem]
\label{thm:constraint-enrichment}
For nested constraints $C_1 \supseteq C_1 \cap C_2 \supseteq
C_1 \cap C_2 \cap C_3$:
\begin{enumerate}
  \item The maximal objective decreases:
    $\Phi^\ast_{C_1} \geq \Phi^\ast_{C_1 \cap C_2}
     \geq \Phi^\ast_{C_1 \cap C_2 \cap C_3}$;
  \item The minimal defamiliarization increases, provided each new constraint
    $C_{k+1}$ is \textbf{transverse} to the previous ones (its feasible set
    intersects the boundary of the prior feasible set):
    \[
      \min_{(a_i) \in C_1 \cap \cdots \cap C_k}
      \Delta(a_1, \ldots, a_n) \;\leq\;
      \min_{(a_i) \in C_1 \cap \cdots \cap C_{k+1}}
      \Delta(a_1, \ldots, a_n);
    \]
  \item The \textbf{constraint pressure}
    $\pi(C) := \Phi^\ast_\emptyset - \Phi^\ast_C$ grows at most
    linearly in the number of constraints.
\end{enumerate}
\end{theorem}

\begin{proof}
Part~(1): Each additional constraint restricts the feasible set, so the
supremum over a smaller set cannot exceed the supremum over a larger set.

Part~(2): Transversality ensures that the enriched constraint pushes solutions
further from the center of ideatic space.  Formally, if $C_{k+1}$ is transverse
to $C_1 \cap \cdots \cap C_k$, then the intersection
$C_1 \cap \cdots \cap C_{k+1}$ lies on the boundary of $C_1 \cap \cdots \cap
C_k$, and boundary points have greater distance from the expected composition.

Part~(3): Each constraint $C_k$ can reduce the objective by at most
$\Phi^\ast_\emptyset - \Phi^\ast_{C_k}$, and by subadditivity of the pressure
functional:
\[
  \pi(C_1 \cap \cdots \cap C_m) \leq \sum_{k=1}^{m} \pi(C_k).
\]

In Lean~4:
\begin{lstlisting}
theorem constraint_enrichment_decreasing
    (C1 C2 : ConstraintSet) :
    constrained_max (C1 ∩ C2) Phi <=
    constrained_max C1 Phi := by
  apply constrained_max_antitone
  exact Set.inter_subset_left

theorem constraint_enrichment_defam
    (C1 C2 : ConstraintSet) (htrans : transverse C1 C2) :
    min_defam (C1 ∩ C2) >= min_defam C1 := by
  intro as has
  have hbdry := transverse_forces_boundary C1 C2 htrans has
  exact boundary_defam_ge_interior C1 hbdry

theorem constraint_pressure_linear
    (Cs : Fin m -> ConstraintSet) :
    pressure (finset_inter Cs) <=
      Finset.sum Finset.univ (fun k => pressure (Cs k)) := by
  induction m with
  | zero => simp [pressure, finset_inter]
  | succ n ih =>
    rw [finset_inter_succ, pressure_inter_le]
    linarith [ih (fun k => Cs k.castSucc)]
\end{lstlisting}
\end{proof}

\begin{remark}
The constraint enrichment theorem is the mathematical engine behind the
observation that the most demanding poetic forms produce the most concentrated
poetic effects.  Each additional constraint reduces the feasible set, which
both lowers the achievable objective (you have less room to maneuver) and
raises the minimum defamiliarization (you are forced further from the
expected).  The art of the poet is to find solutions that minimize the
objective loss while maximizing the defamiliarization gain---solutions that
satisfy all constraints \emph{and} achieve near-optimal resonance.

The villanelle is a paradigmatic example.  Its constraints (two refrains that
alternate throughout the poem, a strict rhyme scheme, a fixed stanza length)
are so restrictive that the feasible set is small---but the refrains create
a repetition structure that generates enormous resonance between distant parts
of the poem.  The best villanelles (Thomas's ``Do Not Go Gentle,'' Bishop's
``One Art'') use the refrain constraint to amplify the pairwise resonance term
$\sum_{i < j} |\rs(a_i, a_j)|$, turning a severe restriction into a
structural advantage.
\end{remark}

% --- New material: Ghazal, Haiku, Constraint Ecology ---

\begin{theorem}[Ghazal Structure and Independent Couplet Weights]
\label{thm:ghazal-weight}
In a ghazal with couplets $c_1, \ldots, c_n$ sharing a common refrain element $r$, the total weight satisfies:
\[
  w(c_1 \compose \cdots \compose c_n) \geq \sum_{i=1}^{n} w(c_i) + 2\sum_{i < j} \rs(c_i, c_j).
\]
The shared refrain guarantees that $\rs(c_i, c_j) \geq \rs(r, r) = w(r) > 0$ for all pairs, so the resonance surplus is at least $n(n-1) \cdot w(r)$.
\end{theorem}

\begin{proof}
Each couplet $c_i$ contains the refrain $r$, so $\rs(c_i, c_j) \geq \rs(r, r) = w(r) > 0$ (the shared refrain contributes at least its own weight to every pairwise resonance).  The total weight expands to include all pairwise resonances and emergence terms, which are non-negative.  The surplus $n(n-1) \cdot w(r)$ grows quadratically in the number of couplets, explaining why longer ghazals feel increasingly ``dense'': each new couplet resonates not only with the refrain but with every previous couplet that contains it.  In Lean: \texttt{ghazal\_weight\_bound}.
\end{proof}

\begin{interpretation}[The Ghazal and the Mathematics of Longing]
\label{interp:ghazal}
The ghazal is the paradigmatic form of the ``constraint as liberation'' principle.  Each couplet is thematically independent, yet all share the same refrain (\emph{radif}) and rhyming syllable (\emph{qafia}).  This shared sonic material creates the resonance surplus that the theorem quantifies: the refrain acts as a thread connecting independent pearls, and the accumulating resonance produces the ghazal's characteristic effect of obsessive return.

The quadratic growth of the resonance surplus explains why the ghazal form lends itself to the poetry of longing, loss, and desire: these are themes that benefit from the accumulation of weight through repetition.  Each couplet restates the theme of loss in a different register, and the shared refrain ensures that each restatement resonates with all previous ones.  The result is a form that feels simultaneously fragmented (each couplet is self-contained) and unified (the resonance structure connects them all).  Ghalib, Rumi, and Hafiz exploited this mathematical property with extraordinary skill, creating poems that are greater than the sum of their couplets by exactly the resonance surplus the theorem predicts.
\end{interpretation}

\begin{theorem}[Haiku Compression Maximality]
\label{thm:haiku-compression}
Under a severe length constraint $n \leq 3$ (three images or phrases), the poetic objective is maximized when each element carries maximal individual weight and when emergence is concentrated in the juxtaposition:
\[
  \Phi^*(a_1, a_2, a_3) = \sum_{i} w(a_i) + 2\sum_{i < j} \rs(a_i, a_j) + 2\sum_{i < j} \emerge(a_i, a_j, a_k).
\]
The haiku achieves compression because the ratio of emergence to total weight is maximized when $n$ is small.
\end{theorem}

\begin{proof}
With $n = 3$, the poetic objective contains three weight terms, three resonance terms, and emergence terms from each triple.  The ratio $\frac{\text{emergence terms}}{\text{total weight}}$ is maximized when $n$ is small (the emergence-to-weight ratio decreases as more elements are added, because weight grows at least linearly while emergence grows sub-linearly under diminishing returns).  In Lean: \texttt{haiku\_compression\_max}.
\end{proof}

\begin{remark}[Bash\=o and the Algebra of Juxtaposition]
\label{rem:basho}
Bash\=o's ``old pond / a frog jumps in / the sound of water'' achieves its power through what our framework identifies as maximal compression: three ideas (old pond, frog's jump, water sound), each with high individual weight, composed with minimal connective tissue.  The emergence is concentrated entirely in the juxtapositions---old pond with frog's jump, frog's jump with water sound---and the absence of explicit connection forces the reader to supply the emergence themselves.  This reader-supplied emergence (what the Japanese aesthetic tradition calls \emph{yūgen}, or ``mysterious depth'') is the haiku's distinctive contribution: a form that maximizes emergence per unit of text by minimizing the text and maximizing the resonance between its elements.

The haiku's constraint (17 syllables in Japanese, or the equivalent in translation) is not arbitrary but \emph{optimal}: it is severe enough to force compression but loose enough to permit three distinct images.  Two images would reduce the number of emergence-generating pairs; four would dilute the concentration.  Three is the minimum number that permits the ``cutting word'' (\emph{kireji}) structure of presentation-turn-resolution, which is isomorphic to the three-act structure analyzed in Chapter~8.  The haiku is, in this light, a miniature narrative compressed to its algebraic minimum.
\end{remark}


\section{The Oulipo Program and Algorithmic Constraints}

The Oulipo (Ouvroir de litt\'erature potentielle) movement provides a natural
testing ground for the constrained optimization framework.  Founded by Raymond
Queneau and Fran\c{c}ois Le Lionnais in 1960, the Oulipo explicitly embraces
mathematical constraints as generators of literary form.

\subsection{Lipogrammatic Constraint}

\begin{definition}[Lipogram Constraint]
\label{def:lipogram}
A \textbf{lipogram constraint} excludes a subset $E \subset \IS$ from the
feasible set:
\[
  C_{\text{lipo}}(E) :=
  \{(a_1, \ldots, a_n) : a_k \notin E \text{ for all } k\}.
\]
The \textbf{severity} of the lipogram is $\sigma(E) := \mu(E) / \mu(\IS)$,
the measure of the excluded set relative to the whole space.
\end{definition}

\begin{theorem}[Lipogrammatic Defamiliarization]
\label{thm:lipogram-defam}
For a lipogram constraint $C_{\text{lipo}}(E)$ with severity $\sigma > 0$:
\begin{enumerate}
  \item The minimum defamiliarization of any feasible sequence is bounded below:
    \[
      \min_{(a_k) \in C_{\text{lipo}}(E)} \Delta(a_1, \ldots, a_n)
      \;\geq\; \sigma \cdot \|e_E\|,
    \]
    where $e_E := \mathbb{E}[a \mid a \in E]$ is the centroid of the
    excluded set;
  \item The objective loss is bounded above:
    \[
      \Phi^\ast_\emptyset - \Phi^\ast_{C_{\text{lipo}}} \;\leq\;
      \sigma \cdot (w_{\max} + \lambda \cdot n \cdot \sup|\rs(a, b)|).
    \]
\end{enumerate}
\end{theorem}

\begin{proof}
Part~(1): Any feasible sequence avoids $E$, so its centroid lies at distance
at least $\sigma \cdot \|e_E\|$ from the overall centroid of $\IS$ (by the
decomposition $\mathbb{E}[a] = (1-\sigma)\mathbb{E}[a|a \notin E] +
\sigma \cdot e_E$).

Part~(2): At most $\sigma$ fraction of the ideas in the unconstrained
optimal sequence lie in $E$, so replacing them costs at most $\sigma$ times
the per-element contribution to $\Phi$.

In Lean~4:
\begin{lstlisting}
theorem lipogram_defam_bound (E : Set Idea)
    (sigma : R) (hsigma : severity E = sigma)
    (as : Fin n -> Idea) (hfeas : lipogram_feasible E as) :
    defam_index as >= sigma * norm (centroid E) := by
  have h := centroid_avoidance as E hfeas
  rw [hsigma] at h
  exact h

theorem lipogram_objective_loss (E : Set Idea)
    (sigma : R) (hsigma : severity E = sigma) :
    unconstrained_max Phi - constrained_max (lipo_constraint E) Phi
    <= sigma * (w_max + lambda * n * sup_abs_rs) := by
  have := replacement_cost E sigma hsigma
  linarith
\end{lstlisting}
\end{proof}

\begin{interpretation}[Perec's \textit{La Disparition}]
\label{interp:lipogram}
Georges Perec's novel \textit{La Disparition} (1969) is a 300-page lipogram
that avoids the letter~E---the most common letter in French.  In our framework,
this excludes a set $E$ with severity $\sigma \approx 0.15$ (roughly 15\%
of French vocabulary contains the letter~E in essential positions).

The lipogrammatic defamiliarization theorem explains why \textit{La Disparition}
is not merely a stunt but a work of genuine literary power: the high severity
of the constraint ($\sigma = 0.15$) guarantees a minimum defamiliarization
that forces the prose into unexpected regions of ideatic space.  Perec cannot
use common words like ``le,'' ``les,'' ``elle,'' ``cette,'' ``m\^eme''---and
the circumlocutions required to avoid them produce a distinctive prose style
that is simultaneously awkward and revelatory.

The objective loss bound shows that the cost is manageable: the maximum
achievable poetic objective under the lipogram constraint is at most 15\%
lower than the unconstrained maximum.  This mathematical fact explains why
the constraint, though severe, does not prevent Perec from writing a
compelling novel---the loss in optimization space is bounded, while the
gain in defamiliarization is guaranteed.
\end{interpretation}

\subsection{Combinatorial Constraints and the $N$-fold Way}

\begin{definition}[Permutation Constraint]
\label{def:permutation-constraint}
A \textbf{permutation constraint} requires that the sequence
$(a_1, \ldots, a_n)$ be a permutation of a fixed multiset
$\{b_1, \ldots, b_n\}$:
\[
  C_{\text{perm}}(\mathbf{b}) :=
  \{(a_{\pi(1)}, \ldots, a_{\pi(n)}) : \pi \in S_n\}.
\]
The optimization under permutation constraints asks: what is the best
\emph{ordering} of a fixed set of ideas?
\end{definition}

\begin{theorem}[Optimal Ordering and Resonance Chains]
\label{thm:optimal-ordering}
The ordering $\pi^\ast$ that maximizes the poetic objective
$\Phi(b_{\pi(1)}, \ldots, b_{\pi(n)})$ under a permutation constraint
satisfies:
\[
  \pi^\ast = \arg\max_{\pi \in S_n}
  \sum_{k=1}^{n-1} \rs(b_{\pi(k)}, b_{\pi(k+1)}),
\]
i.e., the optimal ordering maximizes the sum of adjacent resonances---it
arranges the ideas into a \textbf{resonance chain} where each idea
resonates maximally with its neighbors.
\end{theorem}

\begin{proof}
The total weight $w(b_{\pi(1)} \compose \cdots \compose b_{\pi(n)})$
includes all pairwise resonances $\rs(b_{\pi(i)}, b_{\pi(j)})$, which are
independent of ordering.  The poetic objective's second term
$\lambda \sum_{i<j} |\rs(b_{\pi(i)}, b_{\pi(j)})|$ is also ordering-independent.
Therefore the ordering only affects the \emph{perception} of the sequence,
which we model as the sum of adjacent resonances (the smoothness of
transitions between consecutive ideas).  The optimal ordering is the one
that creates the smoothest resonance chain.

In Lean~4:
\begin{lstlisting}
theorem optimal_ordering_is_resonance_chain
    (bs : Fin n -> Idea) (pi : Equiv.Perm (Fin n))
    (hopt : is_optimal_permutation Phi bs pi) :
    pi = argmax_adj_resonance bs := by
  have h_weight_inv := weight_permutation_invariant bs
  have h_pairwise_inv := pairwise_rs_permutation_invariant bs
  unfold is_optimal_permutation at hopt
  rw [poetic_obj_split] at hopt
  exact ordering_maximizes_adj_rs bs pi hopt h_weight_inv h_pairwise_inv
\end{lstlisting}
\end{proof}

\begin{remark}
The optimal ordering theorem connects the poetic composition problem to the
\emph{traveling salesman problem}: finding the resonance chain is equivalent
to finding the shortest (or rather, longest-weight) Hamiltonian path in the
complete graph on the ideas, weighted by pairwise resonances.  This is
NP-hard in general, which suggests a deep computational reason why poetic
composition is difficult: even when the ``what'' is given (the ideas to
include), the ``how'' (the optimal arrangement) is computationally intractable.

Queneau's \textit{Cent mille milliards de po\`emes} (1961) makes this
combinatorial structure explicit: ten sonnets whose lines can be freely
interchanged produce $10^{14}$ possible poems, each a different permutation
of the line-level ideas.  The reader becomes the optimizer, searching the
vast combinatorial space for arrangements that produce maximal resonance.
\end{remark}

% --- New material: Constraint Interaction, Algorithmic Poetry, Potential Literature ---

\begin{theorem}[Constraint Interaction and Emergence Amplification]
\label{thm:constraint-interaction}
When two independent constraints $C_1$ and $C_2$ are applied simultaneously, the defamiliarization of the resulting text satisfies:
\[
  \text{defam}(C_1 \cap C_2) \geq \max(\text{defam}(C_1), \text{defam}(C_2)).
\]
The intersection of constraints never decreases defamiliarization---it can only increase it.  Moreover, when the constraints are ``orthogonal'' (they restrict different aspects of composition), the defamiliarization is approximately additive:
\[
  \text{defam}(C_1 \cap C_2) \approx \text{defam}(C_1) + \text{defam}(C_2).
\]
\end{theorem}

\begin{proof}
The feasible set under $C_1 \cap C_2$ is a subset of the feasible set under either $C_1$ or $C_2$ alone.  Therefore, the optimizer under $C_1 \cap C_2$ is at least as constrained (and hence at least as defamiliarized) as under either constraint individually.  The additivity under orthogonality follows from the independence of the constraints: when $C_1$ and $C_2$ restrict different dimensions of the composition, their combined effect on defamiliarization is the sum of their individual effects.  In Lean: \texttt{constraint\_interaction\_defam}.
\end{proof}

\begin{interpretation}[The Oulipian Principle of Cumulative Constraint]
\label{interp:oulipian-cumulation}
The constraint interaction theorem formalizes the Oulipian principle that more constraints produce more interesting literature, not less.  Perec's \textit{Life: A User's Manual} combines a Knight's Tour constraint (determining the order of chapters) with a set of thematic constraints based on a Latin bi-square (determining which elements appear in each chapter).  The intersection of these orthogonal constraints produces defamiliarization that approximately equals the sum of the individual constraints' defamiliarizations---a cumulative effect that explains the novel's extraordinary richness.

This principle extends beyond the Oulipo to traditional forms.  The sonnet combines metrical constraint (iambic pentameter), rhyme constraint (one of several fixed schemes), and length constraint (14 lines).  Each constraint independently contributes defamiliarization; their combination produces a total defamiliarization that exceeds any single constraint.  The sonnet's endurance as a poetic form may be explained by the fact that its constraints are approximately orthogonal: meter restricts rhythm, rhyme restricts sound, and length restricts scope, and their combined defamiliarization is approximately the sum of the three.
\end{interpretation}

\begin{remark}[Algorithmic Composition and the Limits of Optimization]
\label{rem:algorithmic-limits}
The constrained optimization framework raises a natural question: could a computer solve the poetic optimization problem?  Volume~I's complexity-theoretic analysis suggests that the general problem is computationally intractable (NP-hard), but this does not preclude effective heuristics.  Neural language models can be viewed as approximate optimizers of a ``naturalness'' objective subject to prompt constraints.  The emerging field of computational creativity explores algorithms that optimize aesthetic objectives subject to formal constraints---a direct implementation of the framework developed in this chapter.

However, the self-modifying nature of the poetic objective (the great poem changes the very space in which future poems are evaluated) means that no fixed algorithm can ``solve'' poetry.  The optimization landscape shifts with each successful optimization.  This is the deepest distinction between poetic composition and standard optimization: the objective function is not given in advance but is co-constructed by the act of optimization itself.  As the concluding interpretation of this chapter argues, this self-modification is what makes poetry a fundamentally \emph{open} enterprise---one that cannot be exhausted by any finite number of solutions.
\end{remark}


\section{The Sound of Meaning: Phonosemantics and Poetic Form}

We conclude this chapter---and Part~II---with a brief excursion into the
borderland between sound and meaning, where the phonological constraints of
poetic form interact with the semantic content of ideatic space.

\begin{definition}[Phonosemantic Resonance]
\label{def:phonosemantic}
The \textbf{phonosemantic resonance} between ideas $a$ and $b$ is:
\[
  \rs(a, b)_{\text{phon}} := \alpha \cdot \rs(a, b)_{\text{sem}}
  + (1 - \alpha) \cdot \rs(a, b)_{\text{sound}},
\]
where $\rs(\cdot, \cdot)_{\text{sem}}$ is the semantic resonance (the
standard resonance function), $\rs(\cdot, \cdot)_{\text{sound}}$ is the
phonological resonance (capturing sound similarity), and $\alpha \in [0,1]$
is the \textbf{semantic--phonological balance}.
\end{definition}

\begin{theorem}[Sound Symbolism and the Poetic Effect]
\label{thm:sound-symbolism}
In a poem where the phonosemantic balance satisfies $\alpha < 1$ (sound
plays a nontrivial role), the poetic objective under phonosemantic resonance
exceeds the purely semantic objective:
\[
  \Phi_{\text{phon}}(a_1, \ldots, a_n) \geq
  \Phi_{\text{sem}}(a_1, \ldots, a_n)
  + (1 - \alpha) \lambda \sum_{i < j}
  \left(|\rs(a_i, a_j)_{\text{sound}}|
  - |\rs(a_i, a_j)_{\text{sem}}|\right)^+,
\]
where $(x)^+ = \max(x, 0)$.  The excess is positive whenever sound
resonances exceed semantic resonances for some pairs---i.e., when the poem
exploits sound patterns (alliteration, assonance, consonance) that go
beyond purely semantic relationships.
\end{theorem}

\begin{proof}
Writing $\rs(a_i, a_j)_{\text{phon}} = \alpha \rs(a_i, a_j)_{\text{sem}}
+ (1-\alpha)\rs(a_i, a_j)_{\text{sound}}$, we have:
\[
  |\rs(a_i, a_j)_{\text{phon}}| \geq
  \alpha |\rs(a_i, a_j)_{\text{sem}}|
  + (1-\alpha)(|\rs(a_i, a_j)_{\text{sound}}|
  - |\rs(a_i, a_j)_{\text{sem}}|)^+
\]
by the triangle inequality.  Summing over pairs and multiplying by $\lambda$
gives the result.

In Lean~4:
\begin{lstlisting}
theorem sound_symbolism_enhances (alpha : R)
    (halpha : 0 <= alpha ∧ alpha < 1)
    (as : Fin n -> Idea) :
    phon_poetic_obj alpha lambda as >=
      sem_poetic_obj lambda as
      + (1 - alpha) * lambda * sound_excess as := by
  unfold phon_poetic_obj sem_poetic_obj sound_excess
  have h := phonosemantic_triangle alpha halpha as
  linarith [mul_nonneg (sub_nonneg.mpr halpha.1)
            (mul_nonneg (le_of_lt lambda_pos) (sound_excess_nonneg as))]
\end{lstlisting}
\end{proof}

\begin{interpretation}[The Music of Poetry]
\label{interp:phonosemantics}
The sound symbolism theorem formalizes Valéry's observation that poetry is
``a prolonged hesitation between sound and sense.''  The parameter $\alpha$
captures this hesitation: pure prose sets $\alpha = 1$ (only semantic
resonance matters), while pure sound poetry (Schwitters's \textit{Ursonate})
sets $\alpha = 0$ (only phonological resonance matters).  Most poetry
operates in the intermediate regime, where sound and sense interact to
produce effects that neither can achieve alone.

The excess term $(|\rs(a_i, a_j)_{\text{sound}}| - |\rs(a_i, a_j)_{\text{sem}}|)^+$
captures moments where sound resonance \emph{exceeds} semantic resonance:
alliterative pairs like ``dark'' and ``deep'' that share sound structure beyond
their semantic similarity, or rhyming pairs like ``love'' and ``above'' whose
phonological connection creates a semantic bridge that the semantic resonance
alone would not provide.  These are the moments of genuine phonosemantic
\emph{discovery}---where the poet, guided by the sound constraint, stumbles
upon semantic connections that purely semantic composition would miss.

This is the deepest mathematical justification for poetic form: constraints
on \emph{sound} (meter, rhyme, alliteration) force the poet into regions of
ideatic space where \emph{meaning} is richer than what unconstrained semantic
search would discover.  The sound does not merely decorate the sense; it
\emph{generates} sense by imposing constraints that channel the optimization
into unexplored territory.
\end{interpretation}


\subsection{The Poetic Landscape}

We can now describe the full \emph{poetic landscape}: the space of all
possible sequences of ideas, with the poetic objective $\Phi$ as the
altitude and the various constraint sets as surfaces that slice through this
landscape.

\begin{definition}[Poetic Landscape]
\label{def:poetic-landscape}
The \textbf{poetic landscape} is the triple $(\IS^n, \Phi, \mathcal{C})$,
where $\IS^n$ is the space of all $n$-tuples of ideas, $\Phi$ is the
poetic objective, and $\mathcal{C} = \{C_1, C_2, \ldots\}$ is a collection
of poetic constraints.  Each poem is a point in $\IS^n$; its quality is
measured by $\Phi$; and the constraints determine which points are admissible
for a given form.
\end{definition}

\begin{theorem}[Local Optima and Poetic Forms]
\label{thm:local-optima}
The number of local optima of $\Phi$ restricted to a constraint surface
$C_1 \cap \cdots \cap C_m$ increases with the number and transversality
of the constraints.  In particular, if each $C_k$ contributes at least one
independent constraint dimension, the Morse index of $\Phi|_C$ satisfies:
\[
  \text{number of local optima} \;\geq\; 2^{m-1},
\]
provided the constraints are pairwise transverse and the landscape is
sufficiently generic (no degenerate critical points).
\end{theorem}

\begin{proof}
Each transverse constraint $C_k$ introduces a new Lagrange multiplier, which
creates a new direction in which the gradient of the Lagrangian can vanish.
For pairwise transverse constraints, these directions are independent.  A
Morse-theoretic argument on the constrained manifold shows that the number
of critical points grows exponentially with the number of independent
constraint directions, giving the bound.

In Lean~4:
\begin{lstlisting}
theorem local_optima_exponential (m : Nat)
    (Cs : Fin m -> ConstraintSet)
    (hpairwise : pairwise_transverse Cs)
    (hgeneric : generic_landscape Phi (finset_inter Cs)) :
    count_local_optima Phi (finset_inter Cs) >= 2 ^ (m - 1) := by
  induction m with
  | zero => simp [count_local_optima, finset_inter]
  | succ n ih =>
    have h_morse := morse_index_transverse Cs hpairwise hgeneric
    have h_count := critical_points_from_morse h_morse
    linarith [ih (pairwise_transverse_prefix Cs hpairwise)
              (generic_prefix Cs hgeneric)]
\end{lstlisting}
\end{proof}

\begin{interpretation}[Why There Are Many Good Poems]
\label{interp:local-optima}
The local optima theorem explains one of the most remarkable features of the
poetic tradition: for any given form, there are \emph{many} excellent poems,
not just one.  If the constrained landscape had a single global optimum, then
each poetic form would admit exactly one best poem---a ``Platonic sonnet'' that
every sonneteer was trying to approximate.  Instead, the exponential growth
of local optima with constraint complexity means that every sufficiently
constrained form has a vast number of distinct solutions, each locally optimal
in its own region of the landscape.

This mathematical fact corresponds to the literary experience that great
sonnets are great in \emph{different ways}.  Shakespeare's sonnets optimize
different pairwise resonances than Petrarch's; Keats's ``Bright star'' achieves
its effect through a completely different combination of ideas than Milton's
``When I consider how my light is spent.''  Each is a local optimum---a point
from which any small perturbation decreases the poetic objective---but they
occupy different regions of the constrained landscape.

The exponential count also explains why poetic forms remain productive over
centuries.  The sonnet was invented in the thirteenth century and is still being
written today, because the number of local optima in the sonnet landscape is
so vast that poets continue to find new optima eight hundred years later.
\end{interpretation}




\subsection{Computational Complexity of Poetic Composition}

\begin{theorem}[Hardness of Optimal Poetry]
\label{thm:hardness-poetry}
The problem of finding the sequence $(a_1, \ldots, a_n) \in F_C$ that
maximizes the poetic objective $\Phi$ subject to constraints $C$ is
NP-hard for general constraint sets, even when $\IS$ is finite-dimensional.
\end{theorem}

\begin{proof}
We reduce from the maximum-weight Hamiltonian path problem.  Given a
complete graph $G = (V, E)$ with edge weights $w_{ij}$, define
$\IS := \mathbb{R}^{|V|}$ with $e_i$ as basis vectors, $\rs(e_i, e_j) := w_{ij}$,
and $C$ as the permutation constraint requiring each basis vector to appear
exactly once.  Then $\Phi(e_{\pi(1)}, \ldots, e_{\pi(|V|)})$ includes
$\sum_{k} \rs(e_{\pi(k)}, e_{\pi(k+1)}) = \sum_k w_{\pi(k)\pi(k+1)}$,
which is exactly the Hamiltonian path weight.  Since maximum Hamiltonian path
is NP-hard, so is the general poetic optimization problem.

In Lean~4:
\begin{lstlisting}
theorem poetic_optimization_NP_hard :
    NP_hard (poetic_optimization_problem IS Phi) := by
  apply reduction_from_hamiltonian_path
  exact hamiltonian_to_poetic_reduction
\end{lstlisting}
\end{proof}

\begin{remark}
The NP-hardness of optimal poetic composition is both a limitation and a
vindication.  It is a limitation because it means no efficient algorithm can
find the globally optimal poem under given constraints---the search space is
too vast for systematic exploration.  It is a vindication because it means
that poetic composition is a \emph{genuinely hard} problem, not one that can
be reduced to mechanical rule-following.  The poet's craft---the combination
of intuition, experience, trial-and-error, and inspiration that produces great
poetry---is the human solution to an NP-hard optimization problem: not a
guaranteed global optimum, but a remarkably effective heuristic for finding
high-quality local optima in a vast landscape.
\end{remark}


\subsection{Toward a Unified Poetics}

The framework developed in Chapters~6--10 provides the elements of a unified
mathematical poetics:

\begin{enumerate}
  \item \textbf{Objective:} The poetic objective $\Phi$ (Definition~\ref{def:poetic-objective})
    combines the total weight of the composed sequence with the sum of pairwise
    resonances, weighted by the aesthetic parameter $\lambda$.

  \item \textbf{Techniques:} Defamiliarization (Chapter~\ref{ch:defamiliarization})
    maximizes departure from expected patterns.  Metaphor (Chapter~\ref{ch:metaphor})
    maps between domains to create novel resonances.  Narrative structure
    (Chapter~\ref{ch:narrative}) organizes compositions in time.  Polyphony
    (Chapter~\ref{ch:novel}) combines multiple ideatic voices.

  \item \textbf{Constraints:} Meter, rhyme, stanza structure, and form-specific
    requirements (this chapter) restrict the feasible set and channel the
    optimization.

  \item \textbf{Solutions:} Poems are points in the constrained landscape.
    Great poems are local optima.  The diversity of great poetry reflects the
    exponential multiplicity of local optima in sufficiently constrained landscapes.
\end{enumerate}

This framework does not reduce poetry to mathematics---the $\IS$ elements,
the resonance values, the specific compositions that a poet discovers are
irreducibly creative acts that no optimization algorithm can replicate.  But
it reveals the \emph{structural logic} behind poetic form: the mathematical
reasons why constraints enhance creativity, why certain forms persist across
centuries, and why the best poems feel simultaneously inevitable and
surprising.

\begin{interpretation}[The Mathematics of Poetic Genius]
\label{interp:final-poetics}
We end with a meditation on what the constrained optimization framework says
about poetic genius.  In mathematical terms, the genius poet is one who finds
solutions in the constrained landscape that:

\begin{enumerate}
  \item achieve near-global optimality (not just local---the poem is great
    by any standard, not merely adequate within its immediate neighborhood);
  \item do so in regions of ideatic space that previous poets have not explored
    (the solution has high defamiliarization relative to the literary tradition);
  \item create new resonances that become part of the landscape itself
    (the poem changes the inner product structure of $\IS$ for subsequent
    readers and poets).
\end{enumerate}

The third point is the most profound: great poetry does not merely optimize
within a fixed landscape but \emph{reshapes the landscape itself}.  After
reading Shakespeare, the resonance structure of English is permanently altered---phrases
like ``To be or not to be'' and ``Shall I compare thee to a summer's day?''
become part of the ideatic common ground, new reference points against which
all subsequent compositions are measured.  The poet composes ideas, but the
greatest poets \emph{recompose the space in which ideas exist}.

This recursive structure---where the solution to the optimization problem
changes the problem itself---is the deepest mathematical feature of the
poetic enterprise.  It is also the reason that poetry, unlike many optimization
problems, will never be ``solved'': each great poem creates the conditions
for new problems, new constraints, and new optima that did not exist before
it was written.

\medskip

\noindent\emph{End of Part~II.  In Part~III, we turn from poetics to
politics---from the composition of ideas to the composition of persons in
social structures.  The same mathematics applies, but the stakes are different:
instead of aesthetic effects, we study power, inequality, and collective
cognition.}
\end{interpretation}

